\documentclass[a4paper,10pt]{article}
\usepackage[top=24mm,bottom=24mm,left=26mm,right=26mm,headheight=14pt]{geometry}
\usepackage{fontspec,amsmath,amsthm,unicode-math}
\usepackage[spanish,es-nodecimaldot,es-noquoting]{babel}
\usepackage{microtype,xcolor,graphicx,tikz,booktabs,array,needspace,fancyhdr}
\usetikzlibrary{arrows.meta,calc,positioning}
\usepackage[unicode,hidelinks]{hyperref}
\setmainfont{STIXTwoText-Regular.otf}[BoldFont=STIXTwoText-Bold.otf,ItalicFont=STIXTwoText-Italic.otf,BoldItalicFont=STIXTwoText-BoldItalic.otf]
\setmathfont{STIXTwoMath-Regular.otf}
\newtheorem{theorem}{Teorema}[section]
\newtheorem{proposition}[theorem]{Proposición}
\newtheorem{lemma}[theorem]{Lema}
\newtheorem{corollary}[theorem]{Corolario}
\theoremstyle{definition}\newtheorem{definition}[theorem]{Definición}
\theoremstyle{remark}\newtheorem{remark}[theorem]{Observación}
\newcommand{\RR}{\mathbb R}\newcommand{\QQ}{\mathbb Q}
\newcommand{\ZZ}{\mathbb Z}\newcommand{\CC}{\mathbb C}
\newcommand{\Gr}{\operatorname{Gr}}\newcommand{\Res}{\operatorname{Res}}
\newcommand{\APP}{\mathrm{APP}}\newcommand{\TPK}{\mathrm{TPK}}
\newcommand{\TRIT}{\mathrm{TRIT}}\newcommand{\id}{\operatorname{id}}
\newcommand{\diag}{\operatorname{diag}}\newcommand{\conv}{\operatorname{conv}}
\numberwithin{equation}{section}
\setlength{\parskip}{4pt plus 1pt}\setlength{\parindent}{1em}
\setlength{\emergencystretch}{2em}
\renewcommand{\normalsize}{\fontsize{10.5}{13.5}\selectfont}\normalsize
\clubpenalty=10000\widowpenalty=10000\displaywidowpenalty=10000
\AddToHook{cmd/section/before}{\Needspace{10\baselineskip}}
\AddToHook{cmd/subsection/before}{\Needspace{6\baselineskip}}
\AddToHook{env/theorem/before}{\Needspace{6\baselineskip}}
\AddToHook{env/proposition/before}{\Needspace{6\baselineskip}}
\pagestyle{fancy}\fancyhf{}
\fancyhead[L]{\footnotesize Geometría positiva y dualidad reticular del registro dodecafásico}
\fancyfoot[C]{\thepage}\renewcommand{\headrulewidth}{.2pt}
\hypersetup{pdftitle={Geometría positiva y dualidad reticular del registro dodecafásico},pdfauthor={Oumar Haidara Fall; Rubén Ramos Balsa}}
\begin{document}
\begin{center}
{\small Manuscrito de recopilación para transferencia académica}\par\vspace{6mm}
{\LARGE\bfseries Geometría positiva y dualidad reticular\\del registro dodecafásico}\par\vspace{3mm}
{\large Realizaciones grassmannianas y amplituhedrales,\\memoria nonádica y límite energético}\par\vspace{5mm}
Oumar Haidara Fall \textperiodcentered\ Rubén Ramos Balsa\par
{\fontsize{8}{10}\selectfont Coautoría sin prelación de contribución\par
Integración y recopilación documental generadas por ChatGPT - Astra.}
\end{center}
\begin{abstract}
El registro dodecafásico de Holografía Modular Triádica determina una
familia de realizaciones geométricas relacionadas por reconstrucción,
refinamiento y transporte. Se demuestra que dos memorias nonádicas y una
carga uniforme recuperan íntegramente sus doce coordenadas; el cono de
registros positivos adquiere así una descripción lineal en variables de
memoria y una cota cuantitativa de estabilidad. La partición ordenada
resultante produce una medición planar de frontera en el Grassmanniano
positivo y, mediante la curva de momentos, amplituedros de rango uno.
Sus formas canónicas admiten triangulaciones con cobertura explícita y
cancelación de residuos interiores. El flujo asociado a la dirección
algebraica del registro induce una deformación efectiva de las razones
cruzadas y una dualidad convexa compatible con la subdivisión heredada.
La misma partición determina una energía nodal cuya reducción de Schur
conserva exactamente la energía efectiva y cuyo complemento reconstruye
el detalle; el sistema compatible de energías acotadas tiene realización
en un espacio de Sobolev de primer orden. En la realización excepcional
de veinticuatro dimensiones, el flujo conserva volumen y simetría ternaria,
y transforma la inversión del parámetro en dualidad reticular. Se obtiene
la correspondiente identidad de inversión theta y una realización
autodual de signatura partida. La orientación temporal del TRIT, sus tres
álgebras cuadráticas y la doble proyección del estado enriquecido sitúan
estos resultados en su genealogía constructiva. Los enlaces con acción,
Hamiltonianos y operadores de masa se formulan mediante sus mapas e
invariantes específicos.
\end{abstract}
\clearpage
\tableofcontents
\clearpage

\section{Genealogía aritmética y sistema de realizaciones geométricas}
\label{sec:genealogy}
La construcción empieza por las dos hojas aritméticas de APP. Si
\(p,q\in\{1,\ldots,9\}\), cada operación conserva una división completa,
\[
 p+q=R_\Sigma+9q_\Sigma,\qquad
 pq=R_\Pi+9q_\Pi.
\]
TRIT añade el régimen local y la orientación. El selector trítico de fase
operativa elige la operación compatible con ese estado, mientras ambas hojas
permanecen en el soporte. Sobre el espacio de estados enriquecidos, el
Topological Phase Kernel compone selección, transporte y actualización:
\[
 U_t=\operatorname{Upd}_t\circ\operatorname{Tra}_t\circ\operatorname{Sel}_t.
\]
El registro conserva los residuos, cocientes, acarreos, hojas, rutas,
fronteras, memoria y supervivencia que utiliza cada prolongación. La
holonomía nonádica \(\Gamma_9=U_{t+8}\circ\cdots\circ U_t\) pertenece
a esta dinámica enriquecida. Su publicación reducida de memoria se efectúa
después. El retorno de fase conserva el avance del registro y por ello
distingue estados situados a profundidades diferentes.

La prolongación
\[
 w_6\longrightarrow w_{12}\longrightarrow w_{18}\longrightarrow
 w_{24}\longrightarrow w_{30}\longrightarrow R_{36}\longrightarrow G_9
\]
organiza los refinamientos compatibles del mismo estado. Las cinco
construcciones del continuo emergen conjuntamente, con sus aspectos
espectral, solenoidal, gauge, cohomológico y determinantal. Los lectores
de ese estado publican correlativamente \(\pi,\varphi,e,\alpha\).
Su origen común admite un orden interno de evaluación: las coordenadas
regionales ya generadas actúan junto con el registro dodecafásico en los
lectores de clausura de \(\alpha\). En este artículo, Grassmannianos,
polígonos y formas diferenciales son realizaciones posteriores de datos
generados por esta cadena. Sus parámetros quedan fijados antes de toda
comparación geométrica o metrológica.

El problema concreto consiste en describir lo que conserva una realización
positiva y lo que puede variar al cambiar su forma. Una misma lista positiva
admite varios usos matemáticos: pesos de aristas, longitudes de intervalos,
coeficientes de una forma cuadrática o coordenadas de un lector material.
Cada uso posee un dominio, una operación y unos invariantes propios. La
expresión \emph{multimorfología} designa aquí el sistema de estas realizaciones
con sus mapas explícitos de compatibilidad. Una identidad de sus números
requiere una comprobación menor que una identidad de sus operadores; el
desarrollo siguiente formula y demuestra cada nivel por separado.

La base de procedencia es el artículo de generación coinductiva de las
constantes y estructura algebraica del electrón, junto con el desarrollo
de extrema y media razón holográfica. Sus construcciones anteriores fijan
el registro utilizado abajo. Los resultados geométricos de este texto comienzan
con esa salida finita y prueban sus consecuencias mediante datos explícitos. Las
proposiciones de transferencia material declaran adicionalmente los
lectores que emplean. Esta organización preserva el alcance específico de
cada prueba.

\section{Orientación trítica, prolongación compatible y doble proyección}
\label{sec:trit-projection}

El carácter generado de las coordenadas reales se expresa en la relación
entre profundidad aritmética y evaluación. El dato primitivo es una historia
admisible de operaciones con sus residuos, cocientes y orientaciones. Una
lectura arquimediana asocia a sus prefijos intervalos cada vez más estrechos;
una lectura incidencial conserva relaciones de frontera del mismo estado.
La doble proyección del corpus \cite{hmtI,HMT-IX} reúne esas dos operaciones
sin convertir una evaluación escalar en la totalidad de la historia. El
artículo utiliza ese antecedente para estudiar las realizaciones posteriores
del registro dodecafásico.

\subsection{Cilindros posicionales y determinación de la coordenada real}

Para un bloque ternario \(b=(b_1,\ldots,b_6)\), con representantes
\(b_j\in\{0,1,2\}\) de \(\mathbb F_3\), la evaluación entera
\(\nu(b)=\sum_{j=1}^6b_j3^{6-j}\) pertenece a
\(\{0,\ldots,728\}\). Si \(b_{n+1}\) es el bloque admisible siguiente,
la prolongación posicional satisface
\[
 N_{n+1}=729N_n+\nu(b_{n+1}),\qquad
 a_n=\frac{N_n}{729^n},\qquad b_n^+=\frac{N_n+1}{729^n}.
\]
Aquí \(N_0\) es la parte entera que determina el lector regional. La letra
\(b_n^+\) designa un extremo real y se distingue del bloque ternario \(b\).
El cilindro de representación es \([a_n,b_n^+)\); su clausura se utiliza
en el argumento de existencia del límite. El acarreo conserva el paso de
un prefijo al anterior mediante división euclídea.

\begin{proposition}[Límite de los centros cilíndricos]
Para una prolongación compatible, las clausuras de los cilindros están
encajadas y determinan una única coordenada \(x\). Si
\(c_n=(a_n+b_n^+)/2\), entonces
\[
 x=\lim_{n\to\infty}a_n=\lim_{n\to\infty}b_n^+
 =\lim_{n\to\infty}c_n,\qquad
 |x-c_n|\leq\frac1{2\,729^n}.
\]
\end{proposition}
\begin{proof}
La desigualdad \(0\leq\nu(b_{n+1})<729\) sitúa el cilindro hijo en
la clausura del padre. Sus extremos son racionales, su anchura es
\(729^{-n}\) y permanecen dentro del primer intervalo cerrado acotado.
El encajamiento y la anchura decreciente producen una única clase de
Cauchy de extremos. En la realización arquimediana, esa clase tiene el
representante \(x\); la distancia al punto medio es a lo sumo la mitad
de la anchura. La convención de frontera asigna una escritura cuando dos
expansiones posicionales representan el mismo límite.
\end{proof}

La media aritmética interviene, por tanto, en una evaluación precisa: el
centro de cada intervalo que acota la coordenada. El límite conserva la
exactitud del valor aunque la anchura desaparezca. La representación
enriquecida conserva adicionalmente la sucesión de prefijos y su incidencia.
La proyección numérica y la proyección excepcional son dos aplicaciones
con codominios diferentes; una conserva el valor límite y la otra transporta
las relaciones que habilitan la construcción reticular. El paso a
\(\RR\) se sitúa así en la realización del lector, después de la
generación aritmética de los datos que evalúa.

Para hacer explícita la compatibilidad, sean \(p_{n+1,n}\) las restricciones
del estado enriquecido y \(x_n=p_{n+1,n}(x_{n+1})\). Si \(I_n(x_n)\)
es el cilindro numérico y \(\mathcal B_n(x_n)\) la incidencia de frontera,
los dos transportes conservan simultáneamente
\[
 \overline{I_{n+1}(x_{n+1})}\subseteq\overline{I_n(x_n)},\qquad
 r_{n+1,n}\mathcal B_{n+1}(x_{n+1})=\mathcal B_n(x_n).
\]
La segunda igualdad utiliza los mapas incidenciales del corpus; la
proposición~\ref{prop:lf-incidence-naturality} explicita su realización
mediante marcos y códigos transportados. Sus cinco
construcciones consustanciales ---espectral, solenoidal, gauge, cohomológica
y determinantal--- pertenecen a ese mismo sistema compatible. El estudio
geométrico posterior conserva esta unidad de origen al especificar cuál
de sus observaciones utiliza en cada teorema.

\subsection{Trayectorias orientadas y temporalidad de la actualización}

Sea \(x\xrightarrow{\rm adm}y\) la relación de admisibilidad del TPK en
el espacio enriquecido \(\widehat X\). Una trayectoria parametrizada por
un conjunto discretamente ordenado \(\mathcal T\) es
\[
 \gamma:\mathcal T\longrightarrow\widehat X,
 \qquad\gamma(t)\xrightarrow{\rm adm}\gamma(t^+).
\]
La relación de compatibilidad determina las transiciones; la parametrización
permite ordenarlas. En el espacio \(E=\mathcal T\times\widehat X\),
con \(p(t,x)=t\), la trayectoria determina la sección
\(\widetilde\gamma(t)=(t,\gamma(t))\), que satisface
\(p\circ\widetilde\gamma=\id_{\mathcal T}\).
Esta precisión conserva la diferencia entre un estado, su trayectoria y
la coordenada utilizada para describir la evolución.

La firma trítica \(\tau\in\{+1,0,-1\}\) tipa los generadores mediante
\[
 J_\tau^2=-\tau I.
\]
En el lector temporal del corpus, \(+1\) representa retorno y memoria;
\(0\), actualización de umbral; y \(-1\), apertura de prolongaciones
conjugadas. Su interpretación como pasado, presente y futuro se refiere
a las funciones del antecedente conservado, la transición efectiva y la
continuación admisible. Los tres tiempos verbales adquieren así un
contenido operatorio sobre historias.

Si \(\mathcal P_n(x)\) es el conjunto de continuaciones admisibles de
longitud \(n\), la supresión del último paso define
\(r_n:\mathcal P_{n+1}(x)\to\mathcal P_n(x)\). Una continuación
ilimitada es una familia \((\gamma_n)\) con
\(r_n(\gamma_{n+1})=\gamma_n\). La restricción reconstruye antecedentes;
la actualización \(U_t\) actúa en la transición; la prolongación mantiene
la frontera del prefijo. La existencia de estas familias se utiliza en
los dominios supervivientes fijados por el teorema de prolongación del
corpus.

Sobre una firma finita \(\sigma=(\tau_1,\ldots,\tau_n)\), la inversión
orientada es
\[
 C_{\rm rev}\sigma=(-\tau_n,\ldots,-\tau_1),
 \qquad C_{\rm rev}^2\sigma=\sigma.
\]
La trayectoria conjugada conserva también las transformaciones de hoja,
posición y frontera de su transporte. Por ello la bidireccionalidad se
expresa en la historia completa, mientras esta fórmula publica su firma.
El retorno de fase de \(\Gamma_9\) admite un incremento de memoria:
estados con fase idéntica pueden ocupar profundidades diferentes.

\subsection{Álgebras cuadráticas y actividad del umbral parabólico}

En un régimen fijo, la conjugación de \(aI+bJ_\tau\) cambia el signo
de \(b\), y su norma algebraica es
\[
 (aI+bJ_\tau)(aI-bJ_\tau)=(a^2+\tau b^2)I.
\]
Los tipos elíptico, parabólico e hiperbólico corresponden a las tres
relaciones cuadráticas. Sus grupos locales se representan por
\[
 e^{tJ_{+1}}=\cos t\,I+\sin t\,J_{+1},\quad
 e^{tJ_0}=I+tJ_0,\quad
 e^{tJ_{-1}}=\cosh t\,I+\sinh t\,J_{-1}.
\]
La prueba separa términos pares e impares en la serie exponencial; en el
tipo nilpotente sólo permanecen los grados cero y uno. Cada familia
satisface la ley de composición en su parámetro y la inversión
\(e^{-tJ_\tau}=(e^{tJ_\tau})^{-1}\).

El valor trítico cero permite, en particular, un generador distinto de
cero con cuadrado nulo. Para
\(N=\left(\begin{smallmatrix}0&1\\0&0\end{smallmatrix}\right)\),
\[
 e^{tN}(q,p)=(q+tp,p).
\]
El umbral tiene una acción lineal efectiva sobre su fibra. Su parámetro,
su desplazamiento y la memoria de la transición permanecen determinados
aunque la firma de régimen sea cero. Éste es el contenido preciso que
permite desarrollar narrativamente el presente como actualización activa.

La forma elíptica positiva tiene niveles circulares en coordenadas
ortonormales y elípticos tras una transformación lineal; la forma hiperbólica
posee dos direcciones isotrópicas y niveles hiperbólicos; el tipo parabólico
posee una forma degenerada y transporte unipotente. Una geometría esférica
requiere además el espacio homogéneo o la métrica que la realiza. La
clasificación anterior fija los generadores que pueden intervenir en esa
realización y mantiene separados el tipo algebraico y la curvatura de un
espacio posterior.

El calendario de modos concreta la conexión con la autoescala. La regla
\(+1:m\mapsto\Sigma\), \(-1:m\mapsto\Pi\), \(0:m\mapsto m\)
aplicada al bloque cíclico \((+1,-1,0)\) produce
\((\Sigma,\Pi,\Pi)\). Contar sus tres transiciones, en el orden
\((\Sigma,\Pi)\), determina
\begin{equation}
 F_{\rm av}=\begin{pmatrix}0&1\\1&1\end{pmatrix}.
 \label{eq:fav}
\end{equation}
Su rayo propio positivo, normalizado como \((1,x)^\mathsf T\), satisface
\(x=\lambda\), \(1+x=\lambda x\), y por ello
\(x^2-x-1=0\). La raíz positiva es la coordenada de autoescala publicada
como \(\varphi\). La matriz nace del calendario de operaciones; la
expresión radical y las identidades exponenciales siguientes son sus
realizaciones posteriores.

% Articulación posterior a la generación de pi, phi y e.
% Contenido acumulativo; no reemplaza el desarrollo previo del TRIT.
\subsection{Realizaciones exponenciales de las coordenadas generadas}
\label{euler-trit-articulacion}

Las coordenadas regionales ya determinadas permiten reconocer las
realizaciones concretas de la familia cuadrática del TRIT. El orden es
sustantivo: la construcción de las regiones precede a la evaluación de una
exponencial en los parámetros \(\pi\) o \(2\log\varphi\). Estas expresiones
identifican la función de las coordenadas en operadores construidos desde
APP y TPK; no intervienen en la selección de los prefijos que las producen.

\subsubsection{Ciclo ternario, transporte nilpotente e incidencia de hojas}

La multiplicación \(a(r)=4r\) en \(\mathbb Z/9\mathbb Z\) tiene orden
\(3\). Sobre las unidades determina las órbitas \(\{1,4,7\}\) y
\(\{2,8,5\}\). En el módulo de aumento de cualquiera de ellas, formado por
las combinaciones enteras cuyos coeficientes suman cero, la base
\((e_r-e_{4r},e_{4r}-e_{7r})\) da
\[
 C=\begin{pmatrix}0&-1\\1&-1\end{pmatrix},
 \qquad C^2+C+I=0.
\]
La forma restringida tiene matriz de Gram
\(\bigl(\begin{smallmatrix}2&-1\\-1&2\end{smallmatrix}\bigr)\).
En las evaluaciones residuales módulo \(9\), la acción distingue las hojas:
\(S(ax,ay)=aS(x,y)\), mientras \(P(ax,ay)=a^{-1}P(x,y)\).
La realización cíclica conserva, por tanto, una orientación contragrediente
entre suma y producto. Los cocientes del levantamiento entero se conservan
en el registro aritmético correspondiente.

El transporte parabólico elemental se representa mediante
\[
 N=\begin{pmatrix}0&1\\0&0\end{pmatrix}.
\]
La incidencia areal--volumétrica ya construida en \eqref{eq:fav} proporciona
\[
 F_{\rm av}=\begin{pmatrix}0&1\\1&1\end{pmatrix},
 \qquad A=F_{\rm av}^2=\begin{pmatrix}1&1\\1&2\end{pmatrix}.
\]
Los tres operadores actúan en espacios reales declarados: el módulo de
aumento tensorizado con \(\mathbb R\), el plano del transporte nilpotente y
el plano de incidencia modal. Compartir dimensión matricial no identifica
estos espacios ni sus coordenadas.

\begin{proposition}[Generadores concretos de los tres tipos]
\label{euler-trit-generadores}
Los operadores
\[
 J_{\mathrm e}=\frac{2C+I}{\sqrt3},\qquad
 J_{\mathrm p}=N,\qquad
 J_{\mathrm h}=\frac{2A-3I}{\sqrt5}
\]
satisfacen, respectivamente,
\[
 J_{\mathrm e}^{\,2}=-I,\qquad
 J_{\mathrm p}^{\,2}=0,\qquad
 J_{\mathrm h}^{\,2}=I.
\]
\end{proposition}
\begin{proof}
La relación de \(C\) implica \((2C+I)^2=-3I\).
La multiplicación directa da \(N^2=0\).
Finalmente, \(A\) tiene traza \(3\) y determinante \(1\), por lo que
\(A^2-3A+I=0\) y \((2A-3I)^2=5I\).
\end{proof}

\begin{proposition}[Tres evaluaciones de la identidad exponencial]
\label{euler-trit-evaluaciones}
En las realizaciones anteriores,
\begin{equation}
 \boxed{
 \exp(\pi J_{\mathrm e})+I=0,\qquad
 \exp(J_{\mathrm p})=I+J_{\mathrm p},\qquad
 \exp(2\log\varphi\,J_{\mathrm h})=F_{\rm av}^{\,2}.}
 \label{euler-trit-tres-evaluaciones}
\end{equation}
\end{proposition}
\begin{proof}
La especialización elíptica de la exponencial trítica da
\(\cos\pi\,I+\sin\pi\,J_{\mathrm e}=-I\).
La nilpotencia elimina exactamente los términos de grado al menos \(2\)
en la segunda expresión. Para la tercera, \(\varphi^2=\varphi+1\) implica
\[
 \cosh(2\log\varphi)=\frac32,\qquad
 \sinh(2\log\varphi)=\frac{\sqrt5}{2}.
\]
Por tanto, la exponencial vale
\(\frac32I+\frac{\sqrt5}{2}J_{\mathrm h}=A\).
\end{proof}

La primera igualdad representa la media vuelta elíptica; la vuelta completa
satisface \(\exp(2\pi J_{\mathrm e})=I\). La segunda mantiene un término
lineal no nulo: el régimen parabólico expresa transporte, no ausencia de
evolución. La tercera realiza un avance hiperbólico unimodular. Sus
autovalores son \(\varphi^2\) y \(\varphi^{-2}\); una dirección se expande y
la otra se contrae conservando el determinante. Las tres evaluaciones
emplean parámetros distintos de una familia exponencial común.

En esa familia, \(e\) interviene mediante la evaluación escalar
\(\exp(I)=eI\) y la ley \(\exp(sI)\exp(tI)=\exp((s+t)I)\).
Su papel no se restringe al tipo parabólico. La coordenada \(e\), obtenida
previamente de su región, es reconocida aquí como la evaluación unitaria
de la exponencial; esta identidad no sustituye su genealogía coinductiva.

\subsubsection{Resolución polinómica de los regímenes}

Las tres realizaciones admiten una presentación compacta sobre la suma
directa de sus espacios. Defínanse
\[
 \mathbb J=J_{\mathrm e}\oplus J_{\mathrm p}\oplus J_{\mathrm h},
 \qquad Q=\mathbb J^2.
\]
Las relaciones anteriores implican
\[
 \mathbb J^6=\mathbb J^2,\qquad Q^3=Q.
\]
El uso de \(Q\) mantiene separado este cuadrado operatorial del registro
dodecafásico \(K\).

\begin{proposition}[Proyectores polinómicos de régimen]
\label{euler-trit-proyectores}
Las aplicaciones
\[
 p_{\mathrm e}=\frac{Q^2-Q}{2},\qquad
 p_{\mathrm p}=I-Q^2,\qquad
 p_{\mathrm h}=\frac{Q^2+Q}{2}
\]
son idempotentes mutuamente ortogonales, suman \(I\) y recuperan los tres
sumandos de la representación.
\end{proposition}
\begin{proof}
Los tres polinomios toman los valores \((1,0,0)\), \((0,1,0)\) y
\((0,0,1)\) en los puntos \(-1,0,+1\), respectivamente. Las diferencias
entre sus cuadrados y ellos mismos, así como sus productos cruzados, son
múltiplos de \(X^3-X\). Sustituir \(Q\), que anula ese polinomio, demuestra
las identidades. En los tres bloques, \(Q\) actúa como \(-I,0,I\).
\end{proof}

Se obtiene así
\begin{align}
 \exp(t\mathbb J)
 ={}&p_{\mathrm e}(\cos t\,I+\sin t\,\mathbb J)
   +p_{\mathrm p}(I+t\mathbb J)\notag\\
   &+p_{\mathrm h}(\cosh t\,I+\sinh t\,\mathbb J).
 \label{euler-trit-exponencial-total}
\end{align}
Cada sumando conserva su relación cuadrática y la conjugación invierte
su evolución. La representación total reúne los tres regímenes, pero no
define por sí sola un operador de transición entre ellos. Tal transición
requiere los mapas de transporte del TPK con sus dominios, orientación y
memoria. La articulación exponencial permite reconocer sus tipos sin
reemplazar esa dinámica por una identificación de las fibras.

% PROCEDENCIA FOCAL (no impresa)
% Resultado recuperado: tres generadores y tres evaluaciones exponenciales.
% Propietario completo:
% BASE/manuscrito/sections/hmt/09b_algebras_cuadraticas.tex:289-516.
% Antecedente causal de C y acción contragrediente:
% BASE/manuscrito/sections/hmt/03d_esqueleto_ciclotomico_app.tex:13-47,145-246.
% Los tres tipos y la orientación están en:
% BASE/manuscrito/sections/hmt/02_trit_geometrias.tex:83-196.
% Resolución polinómica recuperada:
% output/INVESTIGACION_HOLONOMIA_NONADICA_CRISTAL_APERIODICO_20260903/
% ALGEBRA_HOLONOMICA_UNIVERSAL_Y_EULER_TRITICO.md, secciones 8-11 y 22.
% Los cuatro propietarios se leyeron completos. BASE designa:
% /Users/ruben/Documents/HMT2/
% EL_CIERRE_HOLOGRAFICO_DEL_INFINITO_SUCESOR_102_CAPITULOS_2026-08-28_EN_TRABAJO/
% No se incorpora la identidad P9(alpha)=delta del documento especializado:
% esa identificación no pertenece a este bloque.
% Tampoco se identifica una dirección nilpotente con la vacancia electrónica,
% ni se presume una representación global del transporte por la suma directa.
% COMPROBACIONES: Python estándar, enteros y Fraction.
% C^2+C+I=0; (2C+I)^2=-3I; N^2=0; (2Fav^2-3I)^2=5I.
% 81 pares de residuos verifican las dos acciones contragredientes.
% Los 3 proyectores se verificaron en Q=-1,0,+1; prueba general en texto.
% 6 labels únicos y entornos equilibrados. Sin compilación.



\section{Extracción dodecafásica y proyección isotípica de la dirección}
\label{sec:register}
Antes de evaluar números, los lectores distinguen regiones dinámicas.
La simetría especular del corpus actúa sobre el bloque
\(B=(u^\pi,u^e,u^\varphi)\in(\mathbb F_3^6)^3\) por
\[
 \mathfrak cB=(u^e,u^\pi,u^\varphi).
\]
Con \(q=u^\pi+u^e+u^\varphi\),
\(a=u^\pi+u^e-u^\varphi\) y \(d=u^\pi-u^e\), las fórmulas
\[
 u^\varphi=2(q-a),\quad s=q-u^\varphi,\quad
 u^\pi=2(s+d),\quad u^e=2(s-d)
\]
son identidades componente a componente en \(\mathbb F_3\). El espejo
conserva el eje y el agregado en una fibra e invierte el reparto
antisimétrico. El transporte nonádico hace evolucionar la fibra misma.
La quinta fase tiene dato de frontera \(111111\), firma residual
\(000\), una solución local y una extensión: la marca de
\(\varphi\) en cinco identifica el eje regional de esa operación.
En las fases seis, siete y ocho los agregados regionales son
\(012100\), \(101001\) y \(112000\), respectivamente. Estos datos
pertenecen al dominio ternario; los lectores de historias compatibles
producen después sus evaluaciones reales.

El registro de doce ventanas y los bloques regionales conservan así
sus dominios y su enlace en el mismo estado. La realización positiva
que sigue utiliza el registro publicado, con su orden y su escala.
Para transportar el espejo regional a una deformación geométrica debe
componerse primero su acción en ese registro; la igualdad de genealogía
por sí sola expresa menos que la conmutatividad de ese cuadrado. Esta
precisión permite conservar la dinámica regional completa al estudiar
una de sus realizaciones finitas.

El calendario divide las ciento ocho posiciones en doce ventanas nonádicas.
La extracción terminal tiene la composición
\[
 x_{\rm term}\longmapsto\mathcal R_{12}(x_{\rm term})
 \longmapsto(b^{90},b^{120},Q)\longmapsto U_{\rm sgn}\longmapsto K.
\]
Las dos primeras flechas conservan la historia orientada en los lectores
del corpus. Para mostrar íntegramente la reconstrucción finita que utiliza
este artículo, sus salidas enteras son
\begin{align*}
 b^{90}&=(-495,-116,-684,108,601,-90,-173,-88,313,560,-397,461),\\
 b^{120}&=(-425,-281,-481,671,-255,30,475,-543,680,251,6,-128),
 \qquad Q=6263.
\end{align*}
Se fija el orden cíclico de doce posiciones. Sea
\(B_r=\sum_{j=0}^3 b^{120}_{r+3j}\) y
\(d_{r,j}=b^{90}_{r+3j}\), para \(r=1,2,3\). Entonces
\[
 s_1=\frac{Q+2B_1+B_2}{3},\quad s_2=s_1-B_1,\quad s_3=s_2-B_2,
 \qquad
 U^{(r)}=(s_r,d_{r,1}-d_{r,3},d_{r,0}+d_{r,2},d_{r,0}-d_{r,2}).
\]
La evaluación da
\[
 U^{(1)}=(2378,-452,-668,-322),\quad
 U^{(2)}=(1406,998,-204,-28),\quad
 U^{(3)}=(2479,-551,-371,-997).
\]
La matriz
\[
 H_4=\begin{pmatrix}1&1&1&1\\1&1&-1&-1\\1&-1&1&-1\\1&-1&-1&1\end{pmatrix},
 \qquad H_4^2=4I_4,
\]
reconstruye cada órbita mediante \(H_4U^{(r)}/4\). Restituir el orden
de las posiciones produce
\begin{equation}
 K=(234,543,140,729,659,824,621,58,914,794,146,601),
 \qquad Q=\sum_{j=1}^{12}K_j=6263.
 \label{eq:K-main}
\end{equation}
Cada componente es positiva. La división por tres procede de las tres
órbitas y la división por cuatro de la inversa de Hadamard. El control
directo adicional es \((I-S_3)K=b^{90}\),
\((I-S_4)K=b^{120}\), donde \((S_aK)_j=K_{j+a}\) con índices cíclicos.
Estas identidades hacen reproducible la reconstrucción sin seleccionar
coeficientes mediante una geometría objetivo.

La carta incidencial del corpus realiza las doce posiciones como cosets
de \(A_5/C_5\). Si \(\rho\) es esa representación y \(\chi_3\) su
carácter tridimensional real, su proyector es
\[
 P_3=\frac3{60}\sum_{a\in A_5}\chi_3(a^{-1})\rho(a),\qquad
 P_{11}=I-\frac1{12}\mathbf1\mathbf1^{\mathsf T}.
\]
La enumeración y la orientación de cosets se conservan en la fuente de
esta carta. Para que la extensión geométrica sea evaluable directamente,
se escribe su salida exacta \(u=P_3P_{11}K\), con \(r=\sqrt5\):
\begin{equation}
\begin{split}
u=\bigl(&-495/4-233r/10,\quad403/4+169r/10,\\
&-403/4-169r/10,\quad495/4+233r/10,\\
&601/4+1031r/10,\quad203/4+211r/5,\\
&-203/4-211r/5,\quad-601/4-1031r/10,\\
&313/4+569r/10,\quad162+219r/20,\\
&-162-219r/20,\quad-313/4-569r/10\bigr).
\end{split}
\label{eq:u-main}
\end{equation}
Se comprueba por suma y producto en \(\QQ(\sqrt5)\) que
\[
 \mathbf1^{\mathsf T}u=0,\qquad
 \|u\|^2=\frac{6638585+2275584\sqrt5}{20}>0.
\]
En particular, sus coordenadas son no constantes. La dirección \(u\),
el registro íntegro \(K\) y el proyector
\(P_{10}=P_{11}-uu^{\mathsf T}/\|u\|^2\) conservan diferentes cantidades
de información. Esta distinción permite utilizar la dirección para un
flujo sin sustituir por ella el registro que emplean los lectores de acción.

\section{Reconstrucción de la geometría positiva desde la memoria nonádica}
\label{sec:memory}

La extracción dodecafásica conserva más información que una dirección de
deformación. La dirección selecciona un movimiento; el registro completo
determina también la partición sobre la que ese movimiento actúa. La
recuperación del registro a partir de sus memorias permite formular esta
distinción como un teorema de reconstrucción. El resultado procede de la
composición de las compresiones nonádicas del corpus \cite{hmtX}; su aplicación
geométrica conserva la carga uniforme y los complementos de ambas etapas.

En esta sección se numeran las coordenadas de cero a once. Sea
\(\mathsf S:\RR^{12}\to\RR^{12}\) el desplazamiento
\((\mathsf Sx)_i=x_{i+1}\), con índices módulo doce. En la carta euclídea
posterior al estado enriquecido, definimos
\[
 T_j=\frac{8I+\mathsf S^j}{9},\qquad
 D_j=\frac{\sqrt8}{9}(\mathsf S^j-I),\qquad j\in\{3,4\}.
\]
El coeficiente ocho cuenta las posiciones ordinarias de la carta nonádica;
la novena posición contiene el transporte. La ortogonalidad del desplazamiento
implica por multiplicación directa
\[
 T_j^*T_j+D_j^*D_j=I.
\]
La primera memoria es el complemento de la primera compresión; la segunda
se toma sobre el estado que ya atravesó esa compresión. Por tanto,
\begin{equation}
 q=\mathbf1^*K,\qquad m_0=D_3K,\qquad m_1=D_4T_3K,
 \qquad MK=(m_0,m_1).
 \label{eq:mem-data}
\end{equation}
El orden de la composición forma parte del objeto. La aplicación
\(M:\RR^{12}\to\RR^{24}\) conserva dos observaciones correlacionadas,
producidas por una misma trayectoria lineal.

\begin{theorem}[Inversión exacta de la publicación de memoria]
\label{thm:memory-inversion}
La aplicación \(K\mapsto(q,MK)\) es un isomorfismo lineal entre
\(\RR^{12}\) y \(\RR\times\operatorname{im}M\). Su inversa se obtiene
mediante
\begin{align}
 T_3^{-1}&=\frac{512I-64\mathsf S^3+8\mathsf S^6-\mathsf S^9}{455},
 \label{eq:memory-T-inverse}\\
 b&=-\frac9{\sqrt8}m_0,\qquad
 c=-\frac9{\sqrt8}T_3^{-1}m_1,\qquad w_i=c_i-b_{i+1},\notag\\
 B_0&=0,\qquad B_i=\sum_{j=0}^{i-1}w_j,\qquad
 K_i=\frac{q+\sum_{j=0}^{11}B_j}{12}-B_i.
 \label{eq:memory-inverse}
\end{align}
\end{theorem}
\begin{proof}
Escribiendo \(X=\mathsf S^3\), se tiene \(X^4=I\) y
\((8I+X)(512I-64X+8X^2-X^3)=4095I\). Dividir por nueve
prueba \eqref{eq:memory-T-inverse}. Todos los operadores utilizados son
polinomios del mismo desplazamiento y conmutan. Así,
\[
 b=(I-\mathsf S^3)K,\qquad c=(I-\mathsf S^4)K,
 \qquad w_i=K_i-K_{i+1}.
\]
La suma telescópica da \(B_i=K_0-K_i\); sumar esta igualdad sobre
las doce posiciones proporciona \eqref{eq:memory-inverse}. Además,
\(MK=0\) equivale a invariancia bajo \(\mathsf S^3\) y
\(\mathsf S^4\). Como \(\mathsf S=\mathsf S^4(\mathsf S^3)^{-1}\),
el núcleo es exactamente \(\RR\mathbf1\). La carga recupera esa
componente y demuestra la inyectividad. La descomposición
\(K=(q/12)\mathbf1+K^\perp\) prueba la sobreyectividad sobre el
codominio declarado.
\end{proof}

La forma geométrica queda determinada por la información que una lectura
comprimida parecía haber desplazado. El complemento vectorial contiene las
diferencias entre posiciones; la carga contiene la componente común. Su
reunión reconstruye las doce coordenadas antes de normalizarlas como
longitudes. Esta afirmación relaciona operaciones efectivas: compresión,
archivo de complementos e inversión. La conservación de una norma expresa
una propiedad distinta y, por sí sola, más débil que esa recuperabilidad.

\subsection{Cono de memoria compatible y estabilidad de los menores}

Sea \(L\) la pseudoinversa de \(M\), con imagen
\(\mathbf1^\perp\). Equivalentemente,
\[
 L=\left(M^*M\big|_{\mathbf1^\perp}\right)^{-1}M^*,
 \qquad LM=P_{11},\qquad K=\frac q{12}\mathbf1+Lm.
\]
La notación de inversa restringida se interpreta como un operador que toma
valores en \(\mathbf1^\perp\). La positividad estricta del registro
equivale a pertenecer al cono abierto
\begin{equation}
 \mathcal M_+=\left\{(q,m)\in\RR\times\operatorname{im}M:
 \frac q{12}\mathbf1+Lm>0\right\}.
 \label{eq:memory-positive-cone}
\end{equation}
Las desigualdades son coordenadas. El cono tiene dimensión doce; su sección
de carga positiva fija tiene dimensión once. En particular, las veinticuatro
entradas de las dos memorias satisfacen relaciones de compatibilidad.

\begin{proposition}[Constante óptima de reconstrucción en norma euclídea]
\label{prop:memory-stability}
Se cumple \(\|L\|=9/4\). Para perturbaciones arbitrarias de las observaciones,
\begin{equation}
 \|\delta K\|^2\leq\frac{|\delta q|^2}{12}
 +\frac{81}{16}\bigl(\|\delta m_0\|^2+\|\delta m_1\|^2\bigr).
 \label{eq:memory-stability}
\end{equation}
\end{proposition}
\begin{proof}
La base de Fourier diagonaliza \(\mathsf S\). Si \(z^{12}=1\),
el autovalor de \(M^*M\) en el modo correspondiente es
\[
 \lambda(z)=\frac8{81}|z^3-1|^2
 +\frac8{81}|z^4-1|^2\left|\frac{8+z^3}{9}\right|^2.
\]
Evaluar los doce modos da
\[
 \operatorname{Spec}(M^*M)=
 \left\{0^{[1]},(16/81)^{[2]},(8/27)^{[2]},(32/81)^{[1]},
 (952/2187)^{[4]},(1256/2187)^{[2]}\right\}.
\]
Los exponentes entre corchetes indican multiplicidad. El menor autovalor
positivo es \(16/81\), de donde \(\|L\|=9/4\). La componente uniforme
\((\delta q/12)\mathbf1\) es ortogonal a \(L\delta m\); la identidad
pitagórica y la norma de \(L\) producen \eqref{eq:memory-stability}.
\end{proof}

Para pasar a la geometría intervalar recuperamos la numeración de uno a doce
y ponemos \(d_j=K_j/q\), \(t_0=0\),
\(t_j=\sum_{r=1}^j d_r\). El cono \eqref{eq:memory-positive-cone}
determina exactamente \(0=t_0<t_1<\cdots<t_{12}=1\), y por tanto
\[
 \Delta_{ij}=t_j-t_i=\sum_{i<r\leq j}d_r>0\qquad(i<j).
\]
Son los menores de la matriz de caminos que se construye en la sección
\ref{app:network-cluster}. Su dependencia respecto de la memoria es
explícita. Los puntos \((t_j,t_j^2)\) quedan igualmente determinados y
producen la realización amplituhedral de rango uno desarrollada después.

Si la raíz del miembro derecho de \eqref{eq:memory-stability} es menor que
\(\min_j K_j\), la reconstrucción perturbada permanece en el cono positivo.
Para el registro publicado, \(\min_jK_j=58\); a carga fija basta
\[
 \sqrt{\|\delta m_0\|^2+\|\delta m_1\|^2}<\frac{232}{9}.
\]
La cota conserva el orden de los vértices y el signo de todos los menores
ordenados. Describe estabilidad de esta carta matemática, con la
normalización indicada. Para observaciones fuera de \(\operatorname{im}M\),
la reconstrucción utiliza la proyección \(MLm\); el residuo
\((I-ML)m\) se conserva separadamente cuando interesa recuperar la
observación íntegra.

\subsection{Recuperación nodal y reducción a la frontera}

La partición obtenida tiene también una realización cuadrática:
\[
 \mathcal E_d(f)=\sum_{j=1}^{12}\frac{|f_j-f_{j-1}|^2}{d_j}
 =f^*L_df.
\]
El laplaciano nodal recupera cada longitud por
\(d_j=-1/(L_d)_{j-1,j}\). Si se conservan únicamente los extremos
globales, la minimización elimina los nodos interiores y produce, en
cambio, la matriz \(\left(\begin{smallmatrix}1&-1\\-1&1\end{smallmatrix}\right)\),
puesto que \(\sum_jd_j=1\). El lector nodal distingue la forma intervalar;
el lector de dos extremos sólo observa su longitud total.

La demostración de esta eliminación se formula íntegramente en la sección
energética: la prolongación afín minimizante \(Hf\) y el detalle \(z\),
nulo en los nodos antiguos, satisfacen
\[
 \mathcal E_{\widetilde d}(Hf+z)
 =\mathcal E_d(f)+\mathcal E_{\widetilde d}(z).
\]
El segundo término es no negativo y se anula exactamente para el detalle
nulo. La lectura efectiva conserva la energía mínima; el par \((f,z)\)
conserva además la configuración refinada. Así, positividad de menores,
positividad energética y recuperabilidad de memoria se relacionan mediante
mapas definidos sobre la misma partición, manteniendo sus respectivos
espacios de llegada.

% Procedencia HMT: output/EXTREMA_MEDIA_RAZON_NARRACION_Y_REVISION_20260916/
% ARTICULO/ES/source/libro/01_acoplamiento.tex, lib01:cadena-registro y lib01:K.
% Dirección: ARTICULO/ES/source/sections/k_direccion_dimensional.tex.
% Estatuto de las redes, cartas y flujos de este apéndice: FORMALIZACION_NUEVA.
% La extracción de K y el proyector P_3 son RESULTADO_RECUPERADO.
\section{Medición planar de frontera, coordenadas de Plücker y mutaciones positivas}
\label{app:network-cluster}

La publicación dodecafásica del estado enriquecido admite una realización
positiva cuya red, coordenadas y cambios de carta pueden escribirse
completamente. La composición de partida es
\[
 \mathrm{APP}\longrightarrow\mathrm{TRIT}\longrightarrow\mathrm{TPK}
 \longrightarrow x_{\rm term}\longrightarrow\mathcal R_{12}
 \longrightarrow(b^{90},b^{120},Q_{\rm TPK})
 \longrightarrow U_{\rm sgn}\longrightarrow K.
\]
Esta lectura actúa sobre el estado enriquecido que conserva las dos hojas,
la orientación, el acarreo, la memoria y la frontera. El registro ya generado
es
\[
 K=(234,543,140,729,659,824,621,58,914,794,146,601),
 \qquad Q=\sum_{r=1}^{12}K_r=6263.
\]
Su extracción se recupera exactamente de los dos canales y la carga total:
\[
 b^{90}=(-495,-116,-684,108,601,-90,-173,-88,313,560,-397,461),
\]
\[
 b^{120}=(-425,-281,-481,671,-255,30,475,-543,680,251,6,-128).
\]
En las órbitas de posiciones \((r,r+3,r+6,r+9)\), para \(r=1,2,3\),
sean \(B_r=\sum_{j=0}^3b^{120}_{r+3j}\) y
\(d_{r,j}=b^{90}_{r+3j}\). Se obtienen
\[
 s_1=(Q+2B_1+B_2)/3,\quad s_2=s_1-B_1,\quad s_3=s_2-B_2,
\]
\[
 U^{(r)}=(s_r,d_{r,1}-d_{r,3},d_{r,0}+d_{r,2},d_{r,0}-d_{r,2}),
 \qquad K|_r=\tfrac14H_4U^{(r)},
\]
\[
 H_4=\begin{pmatrix}1&1&1&1\\1&1&-1&-1\\1&-1&1&-1\\1&-1&-1&1\end{pmatrix},
 \qquad H_4^2=4I_4.
\]
La red construida a continuación realiza los valores positivos producidos
por esa extracción. Su forma combinatoria es una formalización posterior
del registro, con origen y orientación declarados.

\subsection{Red dirigida acíclica y determinantes de caminos no intersectantes}

Para un registro positivo de longitud \(L\), ponemos
\[
 d_r=K_r/Q>0,\qquad t_0=0,\qquad
 t_j=\sum_{r=1}^jd_r\quad(1\le j\le L).
\]
Así \(t_L=1\). Las \(L\) separaciones producen \(L+1\) puntos marcados;
en el caso dodecafásico son trece puntos. La matriz que los realiza es
\begin{equation}
 C(d)=\begin{pmatrix}1&1&\cdots&1\\t_0&t_1&\cdots&t_L\end{pmatrix}.
 \label{eq:nc-path-matrix}
\end{equation}
El Grassmanniano \(\operatorname{Gr}(2,L+1)\) identifica matrices de rango
dos por cambios invertibles de las dos filas. Su parte positiva está formada
por las clases que admiten todos sus menores ordenados de orden dos
estrictamente positivos.

Construimos una red dirigida con vértices \(a_j,b_j\), para
\(0\le j\le L\), y sumideros \(q_j\). Hay aristas horizontales
\(a_{j-1}\to a_j\) y \(b_{j-1}\to b_j\), de peso uno; aristas verticales
\(a_j\to b_j\), de peso \(d_j\), para \(j\ge1\); y aristas
\(b_j\to q_j\), de peso uno. Las fuentes ordenadas son \(b_0,a_0\).
Las dos filas se dibujan horizontalmente, las aristas verticales descienden
y los sumideros salen por debajo de la fila inferior. Este dibujo es planar
y la orientación es acíclica. La medición de frontera \(M_{ij}\) es la
suma de los productos de pesos de todos los caminos desde la fuente \(i\)
al sumidero \(q_j\).
La red de caminos tiene dos fuentes y \(L+1\) sumideros, es decir,
\(L+3\) conexiones exteriores. La matriz que realiza tiene \(L+1\)
columnas. La semilla plábica introducida después tiene precisamente
\(L+1\) vértices de frontera y constituye otra presentación de sus
coordenadas; estas dos cuentas de frontera se mantienen explícitas.

\begin{theorem}[Medición de frontera y menores del registro]
\label{thm:nc-network}
La medición de frontera de esta red es \(M=C(d)\). Para
\(0\le i<j\le L\),
\begin{equation}
 \Delta_{ij}(C(d))=t_j-t_i=\sum_{r=i+1}^{j}d_r>0.
 \label{eq:nc-minor-gap}
\end{equation}
El mismo número es la suma de pesos de las parejas de caminos
vértice-disjuntos \(b_0\to q_i\), \(a_0\to q_j\).
\end{theorem}
\begin{proof}
La fuente inferior tiene un único camino a cada sumidero, de peso uno.
Un camino desde \(a_0\) a \(q_j\) desciende exactamente una vez, por una
arista de índice \(1\le r\le j\); su peso es \(d_r\). Esto prueba la
fórmula de la matriz. El camino inferior hasta \(q_i\) ocupa todos los
vértices \(b_0,\ldots,b_i\). Por tanto, un camino de la otra fuente a
\(q_j\) es disjunto de él precisamente cuando desciende en
\(i<r\le j\). Su suma de pesos es el miembro derecho de
\eqref{eq:nc-minor-gap}. La pareja con destinos intercambiados siempre
comparte un vértice. El desarrollo del determinante cancela todas las
parejas que se cruzan y conserva exactamente las disjuntas. Aquí se ha
verificado directamente, para la red concreta, el mecanismo determinantal
de caminos no intersectantes.
\end{proof}

La identificación de las mediciones de redes planas con las celdas del
Grassmanniano no negativo proporciona su reconocimiento geométrico posterior;
véase A. Postnikov,
\href{https://arxiv.org/abs/math/0609764}{\emph{Total positivity,
Grassmannians, and networks}}. La prueba anterior da, además, una fórmula
finita para cada menor sin utilizar un valor metrológico.

La familia \(d_r>0\), \(\sum_rd_r=1\), tiene dimensión \(L-1\).
Su aplicación al Grassmanniano es inyectiva: si las filas de \(C(d)\) y
\(C(d')\) generan el mismo plano, las segundas filas difieren por una
transformación afín; las condiciones \(t_0=t'_0=0\) y
\(t_L=t'_L=1\) fijan esa transformación. En particular, para \(L=12\)
se obtiene una familia de dimensión once dentro de la celda positiva de
dimensión veintidós. El punto está en la celda superior; la familia conserva
su dimensión propia. Los trece puntos y las doce separaciones tienen así
papeles matemáticos distintos.

Permitir \(d_r\ge0\) conserva la no negatividad. Cuando \(d_r=0\), las
columnas consecutivas \(r-1\) y \(r\) coinciden y forman elementos
paralelos del matroide. El rango sigue siendo dos porque \(t_0=0\) y
\(t_L=1\). Las bases son exactamente los pares \(i<j\) que atraviesan
alguna separación positiva. Esta degeneración produce estratos de incidencia
consecutiva. Una columna nula, que daría un lazo, corresponde a una operación
distinta de la coincidencia de dos columnas no nulas.

\subsection{Semilla de triangulación y relación de intercambio}

Sea \(N=L+1\) y considérense los vértices \(0,\ldots,L\) de un polígono
convexo en su orden cíclico. A cada lado o diagonal \((i,j)\) se asigna
la coordenada \(A_{ij}=\Delta_{ij}(C(d))\), con índices ordenados.
La triangulación en abanico desde el vértice cero tiene diagonales
\((0,j)\), \(2\le j\le L-1\). Sus lados son variables congeladas y
sus diagonales variables mutables. El quiver se obtiene orientando
cíclicamente los tres lados de cada triángulo y cancelando los pares de
flechas opuestas. Esta es una semilla de tipo \(A_{N-3}\); para el registro
dodecafásico, su tipo es \(A_{10}\).

La semilla tiene una realización plábica explícita. Se coloca un vértice
negro dentro de cada triángulo; en los vértices del polígono se coloca un
vértice blanco si incide en alguna diagonal, y negro en caso contrario.
Se unen los centros a las tres esquinas de sus triángulos, se retiran los
lados originales y se contraen las aristas entre vértices del mismo color.
Se añade una hoja de frontera a cada vértice marcado. Esta construcción
recoge el algoritmo de triangulación de Kodama y Williams,
\href{https://people.math.harvard.edu/~williams/papers/KW-ArxivVersion.pdf}
{\emph{KP solitons and total positivity for the Grassmannian}, algoritmo 12.1}.
Las etiquetas de coordenadas de esta semilla se evalúan en los menores ya
producidos por la red de dos filas. Para una cara mutable, la coordenada
de tipo \(\mathcal X\) es el cociente de los productos de las
coordenadas \(\mathcal A\) incidentes que prescribe su columna del
quiver; en un cuadrilátero se escribirá explícitamente como una razón
doble. Cambiar el grafo cambia la carta de esos menores, manteniendo su
procedencia.

\begin{theorem}[Intercambio positivo de la realización del registro]
\label{thm:nc-ptolemy}
Para cuatro vértices \(a<b<c<d\),
\begin{equation}
 A_{ac}A_{bd}=A_{ab}A_{cd}+A_{ad}A_{bc}.
 \label{eq:nc-ptolemy}
\end{equation}
En consecuencia, el cambio de diagonal \((a,c)\leftrightarrow(b,d)\)
se evalúa mediante
\[
 A_{bd}=\frac{A_{ab}A_{cd}+A_{ad}A_{bc}}{A_{ac}}>0.
\]
Todas las triangulaciones proporcionan cartas positivas compatibles sobre
la misma realización \(C(d)\).
\end{theorem}
\begin{proof}
Al sustituir \(A_{ij}=t_j-t_i\), la igualdad se reduce a
\[
 (t_c-t_a)(t_d-t_b)
 =(t_b-t_a)(t_d-t_c)+(t_d-t_a)(t_c-t_b).
\]
La expansión de ambos miembros prueba la identidad. Las diferencias son
positivas; por tanto, cada flip conserva una evaluación positiva y el
cociente está definido. La conexión de las triangulaciones por flips
permite pasar entre ellas mediante estos intercambios. La interpretación
cluster de la relación de Plücker es el reconocimiento posterior desarrollado
por J. Scott en
\href{https://arxiv.org/abs/math/0311148}{\emph{Grassmannians and Cluster
Algebras}}. La prueba algebraica de la evaluación concreta ya está contenida
en la identidad mostrada.
\end{proof}

\subsection{Refinamiento de las separaciones y naturalidad de los intercambios}

Sea \(B\ge2\). En profundidad \(n\), cada separación \(d_r\) se divide
en \(B^n\) separaciones consecutivas iguales:
\[
 d^{(n)}_{(r-1)B^n+c}=d_r/B^n,
 \qquad 1\le c\le B^n.
\]
Esto define \(LB^n+1\) puntos y una matriz \(C_n\). Si \(\iota_n(j)=Bj\)
identifica un punto con su descendiente de igual posición, entonces
\begin{equation}
 C_{n+1}|_{\iota_n(\{0,\ldots,LB^n\})}=C_n,
 \qquad
 \Delta_{\iota_n(i),\iota_n(j)}(C_{n+1})=\Delta_{ij}(C_n).
 \label{eq:nc-refinement}
\end{equation}
La red refinada reemplaza cada arista vertical de peso \(d\) por una
secuencia de \(B\) posiciones de descenso, con pesos \(d/B\). La suma
de caminos entre extremos heredados sigue siendo \(d\). La composición de
refinamientos conserva esta igualdad porque las sumas finitas se agrupan
asociativamente.

\begin{proposition}[Naturalidad sobre extremos heredados]
\label{prop:nc-inherited}
Todo intercambio \eqref{eq:nc-ptolemy} entre cuatro extremos heredados
conmuta exactamente con el refinamiento. Toda triangulación antigua se
extiende a una del polígono refinado preservando sus triángulos interiores:
los polígonos añadidos entre dos extremos consecutivos se triangulan por
fuera de esos triángulos. Los flips entre diagonales antiguas conservan
entonces sus relaciones de intercambio.
\end{proposition}
\begin{proof}
La primera conclusión resulta al sustituir cada menor por su igualdad en
\eqref{eq:nc-refinement}. Para la segunda, se conservan las cuerdas que
unen extremos antiguos consecutivos y se triangula cada bolsillo exterior
formado por las nuevas marcas. Estas regiones tienen interiores disjuntos.
Los dos triángulos incidentes en una diagonal antigua permanecen iguales,
y, por tanto, también su cuadrilátero y su identidad de Ptolomeo.
\end{proof}

Esta naturalidad especifica los extremos sobre los que actúa. Un lado
antiguo, al recibir vértices nuevos en su arco de frontera, pasa a ser una
diagonal del polígono mayor. Por ello una variable antiguamente congelada
puede ser mutable en el problema ampliado. La compatibilidad anterior es
una compatibilidad exacta del sistema de evaluaciones y de los flips
heredados. Una inclusión de patrones de semillas con igual régimen de
coeficientes se obtiene manteniendo congeladas esas cuerdas y las nuevas
diagonales auxiliares. Esto declara la operación que conserva la semilla;
el incremento de profundidad y una mutación siguen siendo operaciones
distintas.

\subsection{Acción diagonal proyectiva e invariancia de las razones cruzadas}

Si \(\lambda_j>0\), la acción por columnas
\(C\mapsto C\operatorname{diag}(\lambda_0,\ldots,\lambda_L)\) da
\[
 A_{ij}\longmapsto\lambda_i\lambda_j A_{ij}.
\]
Preserva la positividad y todos los estratos definidos por anulación de
menores. La razón doble asociada a un cuadrilátero es
\begin{equation}
 X_{abcd}=\frac{A_{ab}A_{cd}}{A_{ad}A_{bc}}>0.
 \label{eq:nc-crossratio}
\end{equation}
Los cuatro factores \(\lambda_a\lambda_b\lambda_c\lambda_d\) se
cancelan exactamente. En consecuencia, toda acción diagonal positiva por
columnas deja fijo \(X_{abcd}\), incluso si depende de un parámetro y
tiene determinante uno. La modificación del peso de las columnas conserva
el punto proyectivo de cada columna; sus razones dobles expresan esa
conservación. Para que el registro cambie la forma proyectiva debe actuar
sobre las separaciones ordenadas.

\subsection{Deformación positiva de las separaciones dirigida por el registro}

El corpus fija la representación de las doce posiciones y su proyector
ortogonal \(P_3\) sobre el sector tridimensional. Con
\(P_{11}=I_{12}-J_{12}/12\), la dirección recuperada es
\[
 u=P_3P_{11}K=P_3K\in\mathbf1^\perp.
\]
En la carta declarada, poniendo \(r=\sqrt5\), su evaluación exacta es
\begin{align*}
u=(&-495/4-233r/10,\ 403/4+169r/10,\ -403/4-169r/10,\\
   &495/4+233r/10,\ 601/4+1031r/10,\ 203/4+211r/5,\\
   &-203/4-211r/5,\ -601/4-1031r/10,\ 313/4+569r/10,\\
   &162+219r/20,\ -162-219r/20,\ -313/4-569r/10).
\end{align*}
Se conservan \(\sum_ru_r=0\) y
\(\|u\|^2=(6638585+2275584\sqrt5)/20>0\).

Para \(s\in\mathbb R\), definimos la deformación de los pesos mediante
\begin{equation}
 d_r(s)=\frac{K_re^{su_r}}{\sum_{\ell=1}^{12}K_\ell e^{su_\ell}},
 \qquad t_j(s)=\sum_{r=1}^jd_r(s),\qquad C(s)=C(d(s)).
 \label{eq:nc-gap-flow}
\end{equation}
La fórmula utiliza el registro y su dirección ya generados. Es una nueva
realización analítica del lector; el parámetro \(s\) mide su deformación.
Cuando se usa la salida interna \(\rho_A=\pi_{\rm HMT}/\varphi_{\rm HMT}^2\),
la escritura \(s=\tau\log\rho_A\) expresa exactamente la misma familia.
Esa reparametrización actúa después de la generación de las constantes.
Una identificación de \(s\) o \(\tau\) con un tiempo físico requiere la
ley temporal del lector correspondiente.

\begin{theorem}[Cambio de forma, positividad y compatibilidad con la profundidad]
\label{thm:nc-shape-flow}
La familia \eqref{eq:nc-gap-flow} pertenece a
\(\operatorname{Gr}_{>0}(2,13)\) para todo \(s\) real y conserva
\(t_0=0\), \(t_{12}=1\). Para \(i<j\), sean
\[
 \mu_{ij}(s)=
 \frac{\sum_{r=i+1}^jK_ru_re^{su_r}}{\sum_{r=i+1}^jK_re^{su_r}},
 \qquad
 \mu(s)=\frac{\sum_rK_ru_re^{su_r}}{\sum_rK_re^{su_r}}.
\]
Entonces
\begin{equation}
 \frac{d}{ds}\log A_{ij}(s)=\mu_{ij}(s)-\mu(s),
 \quad
 \frac{d}{ds}\log X_{abcd}(s)
 =\mu_{ab}(s)+\mu_{cd}(s)-\mu_{ad}(s)-\mu_{bc}(s).
 \label{eq:nc-crossratio-derivative}
\end{equation}
La deformación conmuta con cada refinamiento que asigna a los hijos de la
separación \(r\) el peso \(d_r/B\) y la misma velocidad \(u_r\).
\end{theorem}
\begin{proof}
Todos los exponenciales y todos los \(K_r\) son positivos. La normalización
hace que la suma de separaciones sea uno; los menores son sumas de
separaciones positivas por \eqref{eq:nc-minor-gap}. Derivar el logaritmo de
\[
 A_{ij}(s)=\frac{\sum_{r=i+1}^jK_re^{su_r}}
                      {\sum_rK_re^{su_r}}
\]
da la primera identidad. La combinación de cuatro logaritmos en
\eqref{eq:nc-crossratio} cancela las cuatro medias globales con sus signos
y da la segunda. Finalmente, reemplazar \(K_r\) por \(B\) copias de
\(K_r/B\) con idéntica \(u_r\) conserva el numerador y el denominador
cuando se agrupan los hijos. La igualdad vale para todos los parámetros,
para cualquier profundidad y para toda composición de refinamientos.
\end{proof}

Hay un testigo exacto de que esta familia cambia la forma. Para los primeros
cuatro extremos,
\begin{equation}
 X_{0123}(0)=\frac{1560}{23711},\qquad
 \left.\frac{d}{ds}\log X_{0123}(s)\right|_{s=0}
 =-\frac{309899}{917}-\frac{268596}{4585}\sqrt5<0.
 \label{eq:nc-nontrivial-witness}
\end{equation}
La cuenta utiliza únicamente las primeras tres separaciones
\((234,543,140)\) y sus velocidades ya recuperadas. La desigualdad
estricta prueba un movimiento proyectivo efectivo. Las identidades de
Ptolomeo siguen siendo exactas durante ese movimiento, pues se aplican a
las columnas \((1,t_j(s))\) en cada parámetro. El registro controla así
una deformación positiva de la forma y, simultáneamente, un sistema de
cartas que permanece compatible con ella.

El alcance demostrado reúne una red planar de rango dos, su medición
explícita, una semilla plábica de triangulación, todas sus mutaciones por
flips, la naturalidad de los extremos heredados y una deformación de
razones dobles dirigida por \(K\). Cada operación tiene un dominio y un
resultado verificables. La comparación de esta realización con la
dinámica temporal enriquecida y con la rama material utiliza sus respectivos
mapas de lectura; ninguna identidad de coordenadas sustituye esos mapas.

\section{Geometrías amplituhedrales de rango uno y compatibilidad de residuos}
\label{sec:k-poligono-amplituhedral}

El registro dodecafásico permite reunir explícitamente datos positivos,
dominio grassmanniano, imagen cinemática, cobertura, grado sobre las celdas
y forma canónica. El primer resultado concierne al amplituedro
bidimensional de parámetros \((n,k,m)=(13,1,2)\); después se prolonga la
construcción a los polítopos cíclicos de rango uno y dimensión superior.
Se trata de una realización geométrica posterior del registro producido
por APP--TRIT--TPK. La operación adicional declarada es la elevación por
potencias de una partición racional ordenada. Esta elección fija una
realización particular; las correspondencias físicas posteriores conservan
sus lectores y condiciones propios.

La procedencia del registro conserva la composición
\[
 x_{\mathrm{term}}\longmapsto\mathcal R_{12}(x_{\mathrm{term}})
 \longmapsto(b^{90},b^{120},Q_{\mathrm{TPK}})
 \longmapsto U_{\mathrm{sgn}}\longmapsto K.
\]
APP conserva ambas hojas y sus residuos y cocientes; TRIT orienta el
régimen de cada transición; TPK transporta y actualiza el estado con
acarreo, frontera, ruta y memoria. La publicación \(K\) se lee en el mismo
sistema de estados enriquecidos que genera conjuntamente el continuo.
Su extracción firmada y su inversión
de Hadamard producen
\begin{equation}
 K=(234,543,140,729,659,824,621,58,914,794,146,601),
 \qquad S_K=\sum_{j=1}^{12}K_j=6263.
 \label{eq:polygon-K}
\end{equation}
La prueba geométrica utiliza la positividad de esas doce coordenadas, su
orden y su total. Sus coeficientes proceden de aquella extracción y no de
la selección retrospectiva de un valor físico.

\subsection{Partición racional y elevación positiva}

Definimos trece puntos racionales de partición:
\begin{equation}
 t_0=0,\qquad t_i=\frac{K_1+\cdots+K_i}{S_K}\quad(1\leq i\leq12),
 \qquad Z_i=(1,t_i,t_i^2)^\mathsf T,\qquad Z=(Z_0\ \cdots\ Z_{12}).
 \label{eq:polygon-Z}
\end{equation}
Sus numeradores, con denominador común \(6263\), son
\[
 (0,234,777,917,1646,2305,3129,3750,3808,4722,5516,5662,6263).
\]
Así, \(0=t_0<t_1<\cdots<t_{12}=1\). La partición etiquetada y el total
conservan el registro mediante la inversa exacta
\begin{equation}
 K_i=S_K(t_i-t_{i-1}).
 \label{eq:polygon-recover-K}
\end{equation}
La partición conserva las proporciones; retener \(S_K\) conserva también
la escala entera.

\begin{proposition}[Datos cinemáticos estrictamente positivos]
\label{prop:polygon-vandermonde}
La matriz \(Z\) tiene rango tres y todos sus menores ordenados de tamaño
tres son positivos. Para \(0\leq a<b<c\leq12\),
\begin{equation}
 \langle abc\rangle:=\det(Z_a,Z_b,Z_c)
 =(t_b-t_a)(t_c-t_a)(t_c-t_b)>0.
 \label{eq:polygon-minors}
\end{equation}
\end{proposition}
\begin{proof}
Se resta la primera columna de las dos últimas y se desarrolla por la
primera fila. Cada diferencia \(t_j^2-t_a^2\) factoriza como
\((t_j-t_a)(t_j+t_a)\), lo que da el producto indicado. Las sumas parciales
estrictamente crecientes prueban su positividad. Cualquiera de esos
menores acredita el rango tres.
\end{proof}

\subsection{Imagen, fibras y celdas de grado uno}

El Grassmanniano \(\operatorname{Gr}_{\geq0}(1,13)\) es el espacio
proyectivo de vectores \(c=(c_0,\ldots,c_{12})\) no negativos y no nulos,
módulo escala positiva. La carta \(\sum c_j=1\) lo identifica con el
símplex cerrado de dimensión doce. Definimos
\begin{equation}
 \Phi_Z:\operatorname{Gr}_{\geq0}(1,13)\longrightarrow\mathbb P^2,
 \qquad[c]\longmapsto\left[\sum_{j=0}^{12}c_jZ_j\right].
 \label{eq:polygon-map}
\end{equation}
En la carta \(Y_0=1\), sean \(v_j=(t_j,t_j^2)\) y
\(P_K=\operatorname{conv}(v_0,\ldots,v_{12})\).

\begin{theorem}[Amplituedro poligonal del registro]
\label{thm:polygon-amplituhedron}
La imagen cerrada de \(\Phi_Z\) es exactamente el polígono convexo
\(P_K=\mathcal A_{13,1,2}(Z)\), de dimensión dos y trece vértices. La imagen
del Grassmanniano estrictamente positivo es el interior de \(P_K\).
Las once celdas
\begin{equation}
 \mathcal S_i=\{[c]:c_0,c_i,c_{i+1}>0,\
 c_j=0\text{ si }j\notin\{0,i,i+1\}\},\qquad 1\leq i\leq11,
 \label{eq:polygon-cells}
\end{equation}
se aplican biyectivamente, con grado uno y orientación positiva, sobre los
interiores de \(T_i=[v_0,v_i,v_{i+1}]\). Sus clausuras cubren \(P_K\) y sus
interiores son disjuntos dos a dos.
\end{theorem}
\begin{proof}
En la normalización \(\sum c_j=1\), la aplicación es
\(c\mapsto\sum c_jv_j\), cuya imagen es la envolvente convexa. Los menores
de \eqref{eq:polygon-minors} muestran que cada triple ordenado tiene
orientación positiva. Cada lado consecutivo deja todos los otros vértices
estrictamente a su izquierda; también lo hace el lado final de \(v_{12}\)
a \(v_0\). Todos los puntos son, por tanto, vértices extremos de un polígono
convexo.

Toda combinación con \(c_j>0\) es interior, pues en cualquier recta de
soporte alguno de los vértices está estrictamente en el semiplano
interior. Recíprocamente, sean \(y\) interior y
\(b=13^{-1}\sum v_j\). Para \(\varepsilon>0\) suficientemente pequeño,
\(y'=(y-\varepsilon b)/(1-\varepsilon)\) permanece en el polígono.
Escribiendo \(y'=\sum a_jv_j\), \(a_j\geq0\), \(\sum a_j=1\), se obtiene
\(y=\sum((1-\varepsilon)a_j+\varepsilon/13)v_j\), con todos los
coeficientes positivos.

Las diagonales desde \(v_0\) dividen el polígono en los once triángulos.
Sobre cada soporte triple, las coordenadas baricéntricas son únicas,
porque \(\langle0,i,i+1\rangle>0\); su inversa explícita aparece en
\eqref{eq:polygon-barycentric}. Cada soporte define una celda positroidal
de rango uno con coordenadas \([1:a:b]\), \(a,b>0\), en esas posiciones.
Una red dirigida planar con una fuente y tres trayectorias de pesos
\(1,a,b\) hacia los tres destinos realiza esas coordenadas por medición
de frontera. El determinante positivo prueba orientación y grado uno.
\end{proof}

La fibra genérica de la aplicación completa tiene dimensión diez:
\(\sum c_j=1\), \(\sum c_jt_j=x\) y \(\sum c_jt_j^2=y\) son tres
ecuaciones independientes en trece coordenadas. Para \(y\) interior la
fibra contiene un punto estrictamente positivo, por lo que su intersección
con el símplex tiene dimensión \(13-3=10\). El grado uno concierne a las
celdas triangulares; la forma bidimensional se obtiene mediante esas
celdas.

\subsection{Forma canónica y residuos orientados}

Para \(Y=(1,x,y)^\mathsf T\), definimos
\begin{equation}
 \ell_{ab}(Y)=\det(Z_a,Z_b,Y)
 =(t_b-t_a)\{y-(t_a+t_b)x+t_at_b\}.
 \label{eq:polygon-lines}
\end{equation}
La expresión vale también cuando \(b<a\). Para \(a<b<c\), sean
\(D_{abc}=\langle abc\rangle\) y
\begin{equation}
 \lambda_a=\frac{\ell_{bc}}{D_{abc}},\qquad
 \lambda_b=\frac{\ell_{ca}}{D_{abc}},\qquad
 \lambda_c=\frac{\ell_{ab}}{D_{abc}},\qquad
 \lambda_a+\lambda_b+\lambda_c=1.
 \label{eq:polygon-barycentric}
\end{equation}
La orientación afín es \(dx\wedge dy\). Se fija la convención de frontera
\(\operatorname{Res}^{\partial}(\eta\wedge du/u)=\eta|_{u=0}\):
la diferencial normal ocupa la última posición. En dimensión dos, el
residuo de Poincaré que coloca \(du/u\) primero tiene el signo opuesto.
Las cancelaciones siguientes usan de forma uniforme la convención declarada.

\begin{theorem}[Forma racional y cancelación de las diagonales]
\label{thm:polygon-form}
La forma canónica del triángulo orientado \([v_a,v_b,v_c]\) es
\begin{equation}
 \Omega_{abc}=
 \frac{D_{abc}^2\,dx\wedge dy}
 {\ell_{ab}(Y)\ell_{bc}(Y)\ell_{ca}(Y)}.
 \label{eq:polygon-triangle-form}
\end{equation}
La forma canónica del polígono es
\begin{equation}
 \boxed{\displaystyle
 \Omega_{P_K}(x,y)=
 \sum_{i=1}^{11}
 \frac{t_it_{i+1}(t_{i+1}-t_i)\,dx\wedge dy}
 {(y-t_ix)\{y-(t_i+t_{i+1})x+t_it_{i+1}\}(t_{i+1}x-y)}.}
 \label{eq:polygon-full-form}
\end{equation}
Los polos de las diagonales interiores se cancelan. Los únicos divisores
polares son las trece rectas de los lados exteriores. Parametrizado cada
lado desde su vértice inicial al final por \(s\in[0,1]\), su residuo de
frontera es \(ds/[s(1-s)]\).
\end{theorem}
\begin{proof}
La forma del símplex baricéntrico es
\(d\lambda_b\wedge d\lambda_c/(\lambda_a\lambda_b\lambda_c)\).
El jacobiano satisface
\(dx\wedge dy=D_{abc}\,d\lambda_b\wedge d\lambda_c\).
Sustituir \eqref{eq:polygon-barycentric} da
\eqref{eq:polygon-triangle-form}. Para \(a=0\), cancelar los tres factores
escalares de las rectas da cada sumando de \eqref{eq:polygon-full-form}.

En el lado \(ab\), pongamos \(\lambda_c=u\),
\(\lambda_b=(1-u)s\), \(\lambda_a=(1-u)(1-s)\). Entonces
\[
 \Omega_{abc}=\frac{ds}{s(1-s)}\wedge\frac{du}{u(1-u)}.
\]
El residuo es la forma del intervalo orientado. Dos triángulos contiguos
inducen orientaciones contrarias sobre su diagonal y producen residuos
opuestos. Al ser simples todos los polos, anular el residuo elimina el
polo de la diagonal como divisor. Cada lado exterior aparece una vez.
La fórmula homogénea es de grado cero como forma proyectiva y no tiene
otros divisores polares. Si dos formas racionales proyectivas tuvieran los
mismos residuos y ningún otro polo, su diferencia sería una dos-forma
holomorfa sobre \(\mathbb P^2\), que es necesariamente cero. Así quedan
probadas caracterización y unicidad.
\end{proof}

Se emplea la noción convencional de geometría positiva mediante polos
logarítmicos y residuos canónicos de frontera, desarrollada por
N.~Arkani-Hamed, Y.~Bai y T.~Lam, \emph{Positive Geometries and Canonical
Forms}, JHEP 11 (2017), 039,
\url{https://arxiv.org/abs/1703.04541}, secciones 2, 3, 5.2, 6.1 y 7.2.
La selección de los coeficientes del presente polígono precede a ese
reconocimiento y procede de \(K\).

\subsection{Refinamiento por acarreo sobre la geometría fija}

En un triángulo \(T=[A,B,C]\) de orientación positiva, escribimos
\begin{equation}
 Y(u,r)=(1-u)A+u\{(1-r)B+rC\},\qquad 0<u,r<1.
 \label{eq:polygon-blowdown}
\end{equation}
El vértice \(A\) corresponde a \(u=0\) y el lado opuesto a \(u=1\).
La forma queda
\begin{equation}
 \Omega_T=\frac{du}{u(1-u)}\wedge\frac{dr}{r(1-r)}.
 \label{eq:polygon-product-form}
\end{equation}
Sea \(B_0\geq2\) una base entera declarada, distinta del vértice \(B\).
La coordenada \(r\) admite la actualización por residuo y acarreo
\begin{equation}
 c=\lfloor B_0r\rfloor,\qquad r'=B_0r-c,\qquad
 r=\frac{c+r'}{B_0}.
 \label{eq:polygon-carry}
\end{equation}
Los extremos conservan sus marcas de frontera. El dígito \(c\) permanece
como memoria y permite invertir la actualización. Este lector usa la
misma operación afín de multiplicación, residuo y acarreo del transporte
\(L_c(x,y)=(B_0x-c,B_0y+c)\). El estado enriquecido retiene la pareja
completa y sus reglas de admisibilidad; la fórmula geométrica especifica
su realización sobre la coordenada seleccionada.

\begin{theorem}[Invariancia exacta bajo subdivisión con memoria]
\label{thm:polygon-carry}
Sean \(s_j=j/B_0\) y \(W_j=(1-s_j)B+s_jC\).
Los triángulos \(T_j=[A,W_j,W_{j+1}]\), \(0\leq j<B_0\), subdividen
\(T\) y satisfacen
\begin{equation}
 \sum_{j=0}^{B_0-1}\Omega_{T_j}=\Omega_T.
 \label{eq:polygon-carry-additivity}
\end{equation}
La igualdad permanece válida a toda profundidad finita y para cualquier
partición racional ordenada del lado. Subdividir los once triángulos
de \(P_K\) deja \(\Omega_{P_K}\) exactamente invariante.
\end{theorem}
\begin{proof}
Las regiones \(r\in[s_j,s_{j+1}]\) cubren \(T\) con interiores disjuntos.
En las coordenadas comunes,
\[
 \Omega_{T_j}=\frac{du}{u(1-u)}\wedge
 \frac{(s_{j+1}-s_j)\,dr}{(r-s_j)(s_{j+1}-r)}.
\]
La identidad racional
\[
 \frac{s_{j+1}-s_j}{(r-s_j)(s_{j+1}-r)}
 =\frac1{r-s_j}+\frac1{s_{j+1}-r}
\]
telescopa: cada término interior cancela el del subintervalo contiguo.
Quedan \(1/r+1/(1-r)=1/[r(1-r)]\). Esto prueba igualdad de formas
racionales; iterarla proporciona toda profundidad finita.
La prueba sólo utiliza \(s_0=0<s_1<\cdots<s_N=1\), de modo que cubre
particiones no uniformes. El cambio
\(r'=(r-s_j)/(s_{j+1}-s_j)\) devuelve localmente la forma unitaria.
La palabra de acarreos identifica la subcelda y su inversa recupera la
coordenada anterior. Los nuevos polos interiores se cancelan mediante
la misma identidad telescópica.
\end{proof}

La suma anterior reúne la partición geométrica completa. Cuando un lector
de supervivencia selecciona sólo algunas historias, su imagen es la unión
de las subceldas seleccionadas y su forma es la suma correspondiente;
coincide con la forma del polígono completo precisamente cuando esa unión
lo cubre con multiplicidad orientada uno. Por tanto, el criterio de
cobertura para una selección concreta de historias queda expresado por
una condición comprobable sobre su cadena de celdas.

Este refinamiento conserva el polígono fijo. Insertar un nuevo punto
\((t,t^2)\) de la parábola entre dos parámetros consecutivos \(a<t<b\)
produce otra geometría: la altura de la cuerda menos \(t^2\) es
\((t-a)(b-t)>0\), y el nuevo punto queda fuera del polígono previo.
La invariancia demostrada corresponde a la subdivisión rectilínea de
celdas, cuya frontera exterior permanece fija.

\subsection{Cambios de triangulación y reparametrización positiva}

Para cuatro vértices ordenados \(A,B,C,D\), se cumple
\begin{equation}
 \Omega_{ABC}+\Omega_{ACD}=\Omega_{ABD}+\Omega_{BCD}.
 \label{eq:polygon-flip}
\end{equation}
Los dos miembros tienen los mismos residuos exteriores y cancelan su
diagonal interior; su igualdad procede de la unicidad demostrada.
También se comprueba multiplicando por sus denominadores lineales.
Todo cambio entre triangulaciones de un polígono convexo se descompone
en estos cambios de diagonal, por lo que la forma es independiente
de la triangulación. Una identificación con coordenadas cluster debe
conservar adicionalmente la variedad y sus relaciones de cambio.

Si \(D=\operatorname{diag}(d_0,\ldots,d_{12})\), \(d_j>0\), entonces
\begin{equation}
 \mathcal A_{13,1,2}(ZD)=\mathcal A_{13,1,2}(Z).
 \label{eq:polygon-rescale}
\end{equation}
El automorfismo \([c]\mapsto[Dc]\) reparametriza la aplicación cinemática.
En la carta afín, \(c_j\) cambia a
\(c_jd_j/\sum_\ell c_\ell d_\ell\), que recorre nuevamente el símplex.
La imagen y su forma canónica permanecen iguales. El desplazamiento
de una historia parametrizada y el cambio del conjunto geométrico
completo son así operaciones distinguibles.

\subsection{Extensión de rango uno a dimensiones superiores}

Para \(2\leq m\leq12\), definimos
\begin{equation}
 Z_i^{(m)}=(1,t_i,t_i^2,\ldots,t_i^m)^\mathsf T,\qquad
 P_K^{(m)}=\operatorname{conv}
 \{(t_i,t_i^2,\ldots,t_i^m):0\leq i\leq12\}.
 \label{eq:polygon-higher-lift}
\end{equation}
Los menores máximos ordenados son los productos de Vandermonde
\(\prod_{a<b}(t_{i_b}-t_{i_a})>0\). La matriz tiene rango \(m+1\).
La imagen de \(\operatorname{Gr}_{\geq0}(1,13)\) es el polítopo cíclico
\(P_K^{(m)}=\mathcal A_{13,1,m}(Z^{(m)})\), de dimensión \(m\); su fibra
genérica tiene dimensión \(12-m\). Aquí se usa la extensión matemática
de la definición amplituhedral a cualquier \(m\); las aplicaciones
habituales a amplitudes imponen sus propios valores de \(m\).

La triangulación se determina por el orden fijo de etiquetas. Se elige
el menor vértice de un polítopo, se triangulan recursivamente todas sus
facetas que no lo contienen y se toman conos desde ese vértice. La
inducción comienza con los puntos. Una recta desde el vértice inicial
hasta un punto interior corta la frontera opuesta en una faceta; ello
prueba la cobertura. La unicidad del punto de intersección y la
compatibilidad de las restricciones inducidas por el mismo orden prueban
que los interiores son disjuntos y que las intersecciones son caras
comunes. En cada paso disminuye la dimensión, por lo que el procedimiento
termina. Esta es una triangulación por tracción determinada por las
etiquetas.

Para un símplex orientado
\(T=[p_0,\ldots,p_m]\) de esa triangulación, sean
\[
 A_T=(p_1-p_0\ \cdots\ p_m-p_0),\qquad
 (\lambda_1,\ldots,\lambda_m)^\mathsf T=A_T^{-1}(x-p_0),
 \qquad\lambda_0=1-\sum_{j=1}^m\lambda_j.
\]
La orientación de sus vértices se toma de modo que \(\det A_T>0\).
Su forma canónica es
\begin{equation}
 \Omega_T(x)=
 \frac{dx_1\wedge\cdots\wedge dx_m}
 {\det(A_T)\prod_{j=0}^m\lambda_j(x)},\qquad
 \Omega_{P_K^{(m)}}=\sum_{T}\Omega_T.
 \label{eq:polygon-higher-form}
\end{equation}
El cambio baricéntrico prueba la fórmula y que la restricción cinemática
sobre el soporte de \(T\) tiene grado uno. Los residuos son las formas
de las caras con los signos inducidos por la orientación. En una cara
interior, las dos orientaciones son contrarias y el residuo total se
anula. En una cara exterior queda la suma de las formas de la
triangulación de esa cara. La inducción sobre la dimensión demuestra
que dicha suma es su forma canónica. Repetir el residuo a lo largo de
una bandera de caras alcanza los vértices, donde la normalización es
\(\pm1\). Se obtiene así la correspondencia de residuos en todas las
codimensiones para esta familia de rango uno.

La misma demostración cubre cualquier subdivisión simplicial compatible
del polítopo fijo: sus caras interiores se cancelan y las exteriores
conservan los residuos. La inexistencia de formas holomorfas de grado
máximo sobre \(\mathbb P^m\) asegura la unicidad de la forma racional.
En particular, la estabilidad bajo refinamiento, la cobertura y los
residuos quedan cerrados en esta familia de dimensiones \(2,\ldots,12\).
Los grados grassmannianos \(k>1\) conservan sus aplicaciones y sus
obligaciones demostrativas específicas.

\begin{figure}[htbp]
\centering
\begin{tikzpicture}[x=1.9cm,y=1.05cm,>=Stealth,font=\small]
\foreach \j/\lab in {0/0,1/1,2/2,3/3,5/12}{
  \node[circle,fill=black,inner sep=1.6pt,label=above:{$a_{\lab}$}] (a\j) at (\j,1.4) {};
  \node[circle,fill=black,inner sep=1.6pt,label=below right:{$b_{\lab}$}] (b\j) at (\j,0) {};
  \node (q\j) at (\j,-1) {$q_{\lab}$};
  \draw[->] (b\j) -- (q\j);
}
\foreach \j/\prev in {1/0,2/1,3/2}{
 \draw[->] (a\prev)--(a\j);\draw[->] (b\prev)--(b\j);
 \draw[->,blue!65!black,thick] (a\j)--node[right]{$d_{\j}$}(b\j);
}
\draw[->,blue!65!black,thick] (a5)--node[right]{$d_{12}$}(b5);
\draw[dashed,->] (a3)--(a5);\draw[dashed,->] (b3)--(b5);
\node at (4,.7) {$\cdots$};
\node[left] at (a0) {$s_2$};\node[left] at (b0) {$s_1$};
\end{tikzpicture}
\caption{Red planar de dos filas. Las aristas horizontales y de salida
tienen peso uno. Cada camino desde la fuente superior desciende una sola
vez. Para dos destinos \(i<j\), las parejas disjuntas descienden después
de \(i\) y antes de \(j\); su peso total es
\(\Delta_{ij}=d_{i+1}+\cdots+d_j\). Los trazos discontinuos conservan
las ocho posiciones intermedias omitidas únicamente en el dibujo.}
\end{figure}

\begin{figure}[htbp]
\centering
\begin{tikzpicture}[x=9cm,y=6cm,font=\small]
\foreach \j/\num in {0/0,1/234,2/777,3/917,4/1646,5/2305,6/3129,7/3750,8/3808,9/4722,10/5516,11/5662,12/6263}{
 \pgfmathsetmacro{\tx}{\num/6263}
 \coordinate (v\j) at (\tx,\tx*\tx);
}
\fill[blue!5] (v0)--(v1)--(v2)--(v3)--(v4)--(v5)--(v6)--(v7)--(v8)--(v9)--(v10)--(v11)--(v12)--cycle;
\foreach \j in {2,...,11}{\draw[blue!45,thin] (v0)--(v\j);}
\draw[blue!80!black,thick] (v0)--(v1)--(v2)--(v3)--(v4)--(v5)--(v6)--(v7)--(v8)--(v9)--(v10)--(v11)--(v12)--cycle;
\foreach \j in {0,...,12}{\fill (v\j) circle[radius=.035cm];}
\node[below left] at (v0) {$v_0$};
\node[below right] at (v3) {$v_3$};
\node[below right] at (v6) {$v_6$};
\node[below right] at (v9) {$v_9$};
\node[above] at (v12) {$v_{12}$};
\end{tikzpicture}
\caption{Polígono racional determinado por los trece parámetros
\(t_j\) de \eqref{eq:polygon-Z}, dibujado con sus coordenadas efectivas
\((t_j,t_j^2)\). La triangulación desde \(v_0\) tiene once celdas.
Cada diagonal interior recibe dos residuos opuestos. Se rotulan cinco
vértices para facilitar la lectura; los trece están representados.}
\end{figure}

\section{Flujo de pesos, geometría de información y dualidad de Legendre}
\label{sec:convex-refinement}

La deformación de los intervalos admite una descripción variacional
exacta. Sean \(d_j=K_j/Q>0\), \(\sum_jd_j=1\), y \(u_j\) las
coordenadas de la dirección \eqref{eq:u-main}. Para \(s\in\RR\),
se definen
\begin{equation}
 Z(s)=\sum_{j=1}^{12}d_j e^{su_j},\qquad
 \psi(s)=\log Z(s),\qquad
 d_j(s)=\frac{d_j e^{su_j}}{Z(s)}.
 \label{eq:convex-partition}
\end{equation}
Todos los argumentos de la exponencial son adimensionales. El parámetro
\(s\) es la coordenada de esta deformación geométrica. La asignación de
una duración física pertenece a un lector temporal declarado.

\begin{theorem}[Convexidad estricta y coordenada dual]
La función \(\psi\) es analítica y estrictamente convexa sobre \(\RR\).
Si \(u_- =\min_j u_j\) y \(u_+=\max_j u_j\), la aplicación
\(v=\psi'(s)\) es un difeomorfismo de \(\RR\) sobre \((u_-,u_+)\).
Para cada \(v\) en ese intervalo existe una única coordenada \(s(v)\), y
\[
 \psi^*(v)=s(v)v-\psi(s(v)),\qquad
 (\psi^*)'(v)=s(v),\qquad
 (\psi^*)''(v)=\frac1{\psi''(s(v))}>0.
\]
\end{theorem}
\begin{proof}
Las sumas finitas permiten diferenciar término a término. Se obtiene
\begin{align}
 \psi'(s)&=\sum_jd_j(s)u_j=:\overline u(s),\notag\\
 \psi''(s)&=\sum_jd_j(s)(u_j-\overline u(s))^2.
 \label{eq:variance}
\end{align}
Los pesos son estrictamente positivos y los \(u_j\) no son todos iguales;
la varianza es por ello estrictamente positiva. Cuando \(s\) tiende a
\(+\infty\), dividir numerador y denominador de \(\psi'\) por
\(e^{su_+}\) muestra que sólo sobreviven los índices con \(u_j=u_+\).
Así \(\psi'(s)\to u_+\). El argumento con \(u_-\) da el límite
opuesto. La monotonía y el teorema de la función inversa proporcionan la
biyectividad suave. La función \(sv-\psi(s)\) tiene su máximo único
en \(s(v)\); diferenciar el valor máximo demuestra las dos identidades
duales.
\end{proof}

La evolución de cada intervalo es
\begin{equation}
 \dot d_j(s)=d_j(s)\bigl(u_j-\overline u(s)\bigr).
 \label{eq:replicator}
\end{equation}
La suma de las derivadas es cero. Por tanto la longitud total permanece
unitaria mientras se redistribuyen las longitudes internas. Esta conservación
de escala global permite atribuir el cambio de las razones dobles a una
deformación de posiciones y no a una dilatación común.

En el simplex abierto de pesos, la métrica
\(g_d(\xi,\zeta)=\sum_j\xi_j\zeta_j/d_j\), sobre vectores tangentes
de suma cero, es positiva. La función lineal
\(F(d)=\sum_jd_ju_j\) tiene como gradiente para esta métrica
precisamente el campo de \eqref{eq:replicator}, pues
\[
 g_d\bigl(d_j(u_j-\overline u),\xi\bigr)
 =\sum_j(u_j-\overline u)\xi_j=\sum_ju_j\xi_j=dF(\xi).
\]
En consecuencia \(dF/ds=\psi''(s)>0\) a lo largo del flujo. El
carácter ascendente corresponde a esta función y a esta métrica concretas.

\begin{theorem}[Invariancia variacional por subdivisión heredada]
Para cada intervalo \(j\), sean \(r_{j,c}>0\),
\(\sum_c r_{j,c}=1\), los pesos relativos de una subdivisión fija.
Asigne a los hijos
\[
 d_{j,c}=r_{j,c}d_j,\qquad u_{j,c}=u_j.
\]
La función de partición, su logaritmo, todas sus derivadas y la transformada
de Legendre coinciden exactamente con los de la partición madre.
Además, el flujo conmuta con la suma de hijos:
\[
 \sum_c d_{j,c}(s)=d_j(s),\qquad
 d_{j,c}(s)=r_{j,c}d_j(s).
\]
\end{theorem}
\begin{proof}
Por agrupación de la suma finita,
\[
 Z_{\rm ref}(s)=\sum_{j,c}r_{j,c}d_j e^{su_j}
 =\sum_j d_j e^{su_j}=Z(s).
\]
Todas las conclusiones se deducen de esta identidad y de dividir cada
peso por el mismo \(Z(s)\). El mismo argumento vale después de cualquier
número finito de subdivisiones. Para una subdivisión numerable, la suma de
términos positivos converge al mismo valor mediante su agrupación por padres.
\end{proof}

La hipótesis de dirección heredada tiene contenido verificable. Si un
intervalo con dirección \(a\) se divide en dos mitades de direcciones
\(a+b\) y \(a-b\), su contribución cambia de
\(d e^{sa}\) a \(d e^{sa}\cosh(sb)\). Para \(b\ne0\) la igualdad
de funciones de partición falla aunque el peso total en \(s=0\) sea
el mismo. Así, la conservación geométrica y variacional exige conservar
también la información direccional pertinente. La identidad demostrada
se aplica a la subdivisión heredada; su transporte a una actualización
TPK que cambie esas direcciones se decide componiendo el lector de dicha
actualización.

Existe una formulación hamiltoniana regular de la coordenada dual. En el
plano cotangente con coordenadas \((q,p)\), tome
\(H(q,p)=\psi(p)\). Las ecuaciones son
\(\dot q=\psi'(p)\), \(\dot p=0\). Para velocidades
\(\dot q\in(u_-,u_+)\), la transformación de Legendre da
\[
 L(q,\dot q)=\psi^*(\dot q),\qquad
 \frac{\partial^2L}{\partial\dot q^2}>0.
\]
Este sistema realiza variacionalmente el potencial del flujo de intervalos.
Su ecuación \(\dot p=0\) muestra que el tiempo hamiltoniano es distinto
del parámetro \(s\) que recorre la familia \eqref{eq:convex-partition}.
La evaluación de \(H\) y \(L\) se conserva bajo la subdivisión del
teorema, porque \(\psi\) y \(\psi^*\) permanecen idénticas.

Por su parte, cualquier flujo \(\dot d=F(d)\) admite en una carta local
del simplex una elevación cotangente con
\(\mathscr H(d,p)=p\cdot F(d)\). Esta segunda construcción reproduce
el propio flujo \eqref{eq:replicator} en su base, pero su Hessiano respecto
de \(p\) es nulo. Las dos realizaciones responden a preguntas distintas:
una produce una dualidad convexa regular y la otra transporta una dinámica
de base dada. La distinción preserva el significado de Hamiltoniano y
Lagrangiano al comparar esta geometría con los lectores físicos del corpus.

% Desarrollo acotado: lector energético intervalar del registro K.
% Estatuto: formalización añadida sobre K y sobre las identidades de Schur
% ya presentes en el corpus. No se atribuye novedad histórica a Schur.
\section{Formas de Dirichlet y equivalencia variacional del refinamiento}
\label{sec:energia-refinamiento-k}

La publicación positiva del registro dodecafásico define una partición
ordenada mediante
\[
 K=(234,543,140,729,659,824,621,58,914,794,146,601),
 \qquad d_j=K_j/6263,
 \qquad t_i=\sum_{j=1}^id_j.
\]
El registro se recibe de la composición
\(x_{\rm term}\mapsto\mathcal R_{12}\mapsto
(b^{90},b^{120},Q_{\rm TPK})\mapsto U_{\rm sgn}\mapsto K\),
posterior a APP--TRIT--TPK y a su estado enriquecido, conforme a los lectores
reconstruidos en \cite{hmtX}. Esta sección aplica
un lector cuadrático a la partición resultante. La comparación entre
profundidades se efectuará por restricciones a extremos heredados, de modo
que cada coeficiente de la energía conserve su origen en una separación
positiva ya producida.

El corpus contiene una energía ponderada de diferencias para el grafo de
las dos memorias y una reducción de Schur de las formas completas. Aquí se
reúnen esas operaciones sobre la red intervalar del registro. La
especialización a conductancias \(1/d_j\), el operador de interpolación
y su compatibilidad con la deformación de las separaciones constituyen
la formalización desarrollada en esta sección.

\subsection{Reconstrucción de las separaciones a partir de la matriz de rigidez}

Sea \(d=(d_1,\ldots,d_m)\), con \(d_j>0\), y sean
\(f=(f_0,\ldots,f_m)\in\mathbb C^{m+1}\) los valores de un observable
en los extremos de la partición. Definimos
\begin{equation}
 \mathcal E_d(f)=\sum_{j=1}^{m}\frac{|f_j-f_{j-1}|^2}{d_j},
 \qquad L_d=D^*\operatorname{diag}(d_1^{-1},\ldots,d_m^{-1})D,
 \label{eq:energia-intervalar}
\end{equation}
donde \(D:\mathbb C^{m+1}\to\mathbb C^m\) es el operador de
incidencia orientada, \((Df)_j=f_j-f_{j-1}\), y el asterisco denota la
adjunción hermítica. Así \(\mathcal E_d(f)=f^*L_df\).
La matriz de rigidez tiene entradas
\[
 (L_d)_{00}=d_1^{-1},\quad (L_d)_{mm}=d_m^{-1},\quad
 (L_d)_{ii}=d_i^{-1}+d_{i+1}^{-1}\quad(1\le i<m),
\]
\[
 (L_d)_{i-1,i}=(L_d)_{i,i-1}=-d_i^{-1},
\]
y las demás son nulas. Es hermítica positiva semidefinida, su núcleo
consiste en las funciones constantes y su forma es positiva definida
en el cociente por ese núcleo. La matriz completa recupera las
separaciones mediante
\begin{equation}
 d_j=-\frac1{(L_d)_{j-1,j}}.
 \label{eq:recuperacion-d-rigidez}
\end{equation}

\begin{proposition}[Equivalencia entre separaciones y lector nodal completo]
\label{prop:energia-bidireccional}
La aplicación \(d\mapsto L_d\) es una biyección entre
\((0,\infty)^m\) y el conjunto de matrices reales simétricas
tridiagonales cuyas entradas inmediatamente adyacentes a la diagonal son
negativas y cuya suma en cada fila es cero. Su inversa es
\eqref{eq:recuperacion-d-rigidez}. La normalización \(\sum_jd_j=1\)
equivale a \(\sum_{j=1}^m-1/(L_d)_{j-1,j}=1\).
\end{proposition}
\begin{proof}
Las fórmulas explícitas de \(L_d\) verifican todas las condiciones.
Recíprocamente, las entradas adyacentes negativas determinan separaciones
positivas únicas por \eqref{eq:recuperacion-d-rigidez}. La suma nula de
cada fila determina sus entradas diagonales como la suma de las
conductancias incidentes. La tridiagonalidad fija en cero las restantes
entradas. Se recupera exactamente la matriz de
\eqref{eq:energia-intervalar}; su semidefinitud positiva sigue entonces
de la factorización \(D^*\operatorname{diag}(1/d_j)D\).
\end{proof}

La derivada discreta ponderada
\(J_j=(f_j-f_{j-1})/d_j\) es el flujo orientado de la arista \(j\).
En los vértices interiores, \((L_df)_i=J_i-J_{i+1}\); en los extremos,
\((L_df)_0=-J_1\) y \((L_df)_m=J_m\). La ecuación interior
\(L_df=0\) expresa que el flujo tiene el mismo valor en cada arista.
Esta definición es la de un observable cuadrático del lector intervalar;
su realización como magnitud física requiere la carta dimensional pertinente.

\subsection{Interpolación armónica, complemento de Schur y reconstrucción del detalle}

Refinemos cada separación como
\(d_j=\sum_{c=1}^{r_j}\delta_{j,c}\), con \(\delta_{j,c}>0\).
Sea \(\widetilde d\) la concatenación de esas separaciones, \(R\) la
restricción a los extremos antiguos y \(H\) la interpolación definida por
\begin{equation}
 (Hf)_{j,c}=(1-\theta_{j,c})f_{j-1}+\theta_{j,c}f_j,
 \qquad
 \theta_{j,c}=\frac{\sum_{a=1}^c\delta_{j,a}}{d_j}.
 \label{eq:interpolacion-armonica-k}
\end{equation}
Los extremos compartidos se cuentan una sola vez. Se cumple \(RH=I\).

\begin{theorem}[Refinamiento isométrico de la energía]
\label{thm:energia-schur-refinamiento}
Para todo vector de valores antiguos \(f\),
\begin{equation}
 H^*L_{\widetilde d}H=L_d,
 \qquad L_{\widetilde d}H=R^*L_d,
 \qquad
 \min_{Rg=f}\mathcal E_{\widetilde d}(g)=\mathcal E_d(f).
 \label{eq:isometria-energia-k}
\end{equation}
El minimizador es único y vale \(Hf\). Todo \(g\) admite la
descomposición única
\begin{equation}
 g=H(Rg)+z,\qquad Rz=0,
 \qquad
 \mathcal E_{\widetilde d}(g)
 =\mathcal E_d(Rg)+\mathcal E_{\widetilde d}(z).
 \label{eq:detalle-energia-k}
\end{equation}
\end{theorem}
\begin{proof}
En una arista antigua escribamos \(h_c=g_{j,c}-g_{j,c-1}\), de modo
que \(\sum_ch_c=f_j-f_{j-1}=a\). Si
\(h_c=\delta_{j,c}a/d_j+r_c\), entonces \(\sum_cr_c=0\) y
\[
 \sum_c\frac{|h_c|^2}{\delta_{j,c}}
 =\frac{|a|^2}{d_j}+\sum_c\frac{|r_c|^2}{\delta_{j,c}}.
\]
La igualdad se obtiene expandiendo; el término cruzado es proporcional a
\(\sum_cr_c\) y se anula. El mínimo se alcanza exactamente cuando
todos los \(r_c\) son cero, condición que determina
\eqref{eq:interpolacion-armonica-k}. Al sumar las aristas se obtienen
el mínimo y la descomposición energética. El flujo de \(Hf\) es
constante dentro de cada arista antigua, por lo que sus componentes
interiores de \(L_{\widetilde d}Hf\) son cero y sus componentes heredadas
coinciden con \(L_df\). Esto prueba la segunda identidad matricial; al
multiplicar por \(H^*\), usando \(RH=I\), se obtiene la primera.
\end{proof}

El término \(\mathcal E_{\widetilde d}(z)\) es no negativo y se anula
exactamente para \(z=0\): un detalle de energía nula sería constante y,
al anularse en los extremos heredados, se anularía en toda la red.
La pareja \((Rg,z)\) conserva el vector refinado completo y permite
recuperarlo por \(g=H(Rg)+z\). La minimización selecciona su
componente armónica; el archivo del detalle \(z\) mantiene la
información que la sola forma reducida no distingue. La conservación
de la energía observada y la conservación del estado refinado quedan
así expresadas por mapas diferentes y explícitos.

Ordenando los vértices como extremos heredados \(E\) e interiores \(I\),
escribimos
\[
 L_{\widetilde d}=\begin{pmatrix}A&B^*\\B&C\end{pmatrix}.
\]
Para un refinamiento con vértices interiores, el bloque \(C\) es
positivo definido: una función nula en los extremos
que tenga energía cero es constante en cada subintervalo y, por tanto,
idénticamente nula. Se obtiene
\begin{equation}
 H=\begin{pmatrix}I\\-C^{-1}B\end{pmatrix},\qquad
 \operatorname{Schur}_{I}(L_{\widetilde d})
 =A-B^*C^{-1}B=L_d.
 \label{eq:schur-intervalar-k}
\end{equation}
Si el refinamiento conserva todos los vértices y no introduce ninguno,
la misma afirmación se reduce a \(H=I\) y \(L_{\widetilde d}=L_d\).

\begin{proposition}[Naturalidad de las prolongaciones y de la eliminación]
\label{prop:energia-naturalidad-sucesiva}
Para tres particiones anidadas, las interpolaciones armónicas y los
complementos de Schur satisfacen
\begin{equation}
 H_{n+2\leftarrow n}
 =H_{n+2\leftarrow n+1}H_{n+1\leftarrow n},
 \qquad
 \operatorname{Schur}_{n+1\to n}
 \bigl(\operatorname{Schur}_{n+2\to n+1}(L_{n+2})\bigr)=L_n.
 \label{eq:energia-naturalidad-sucesiva}
\end{equation}
\end{proposition}
\begin{proof}
La interpolación de valores afines sobre un intervalo sigue siendo la
misma función afín después de cualquier subdivisión de ese intervalo.
Esto prueba la primera identidad. Minimizar primero sobre los vértices
incorporados en el nivel último y después sobre los del nivel intermedio
es minimizar sobre todos ellos, con los mismos extremos heredados fijados.
La identidad variacional de \eqref{eq:isometria-energia-k} y la unicidad
del minimizador prueban la segunda identidad. La restricción del mínimo
se realiza, por tanto, mediante la misma composición que la prolongación
armónica de los datos.
\end{proof}

\subsection{Respuesta de Dirichlet a Neumann y defecto energético de una extensión}

Si sólo se conservan los extremos inicial y final, sea
\(D_0=\sum_{j=1}^md_j\). Para valores \((a,b)\), la solución
armónica es
\[
 f_i=a+(b-a)\frac{\sum_{j=1}^id_j}{D_0},\qquad
 J_j=(b-a)/D_0.
\]
Su operador de Dirichlet a Neumann, que transforma valores de frontera en
flujos salientes, es
\begin{equation}
 \Lambda_d=\frac1{D_0}
 \begin{pmatrix}1&-1\\-1&1\end{pmatrix},\qquad
 \min\mathcal E_d=\frac{|b-a|^2}{D_0}.
 \label{eq:dtn-intervalar-k}
\end{equation}
Cuando se prescribe el valor en todos los vértices heredados, su operador
de respuesta multipuntual es, por el mismo argumento, \(\Lambda=L_d\).

Una extensión aproximada \(g=(f,y)\), donde \(y\) reúne los valores
en los vértices interiores, tiene residual interior
\(r=Bf+Cy\). La identidad exacta
\begin{equation}
 \mathcal E_{\widetilde d}(g)-f^*\Lambda f
 =r^*C^{-1}r\ge0
 \label{eq:residual-dtn-k}
\end{equation}
se deduce de \(y+C^{-1}Bf=C^{-1}r\) y de la descomposición
\[
 g^*L_{\widetilde d}g
 =f^*(A-B^*C^{-1}B)f
 +(y+C^{-1}Bf)^*C(y+C^{-1}Bf).
\]
Así, el valor efectivo y el detalle mantienen una descomposición
hermítica completa. En particular, si una
extensión lineal es \(\widetilde H=\begin{pmatrix}I\\Z\end{pmatrix}\), su matriz de
residuales es \(\mathscr R=B+CZ\) y
\[
 \widetilde H^*L_{\widetilde d}\widetilde H-\Lambda
 =\mathscr R^*C^{-1}\mathscr R\succeq0.
\]
El residual nulo caracteriza la interpolación armónica exacta. Este
residual energético vive en el espacio de vértices interiores; la
cancelación de residuos de formas meromorfas pertenece a otro mapa de
frontera y conserva su definición propia.

\subsection{Conmutación del refinamiento y del flujo de separaciones}

Sea \(u=P_3P_{11}K\) la dirección algebraica ya recuperada de los
lectores del registro \cite{hmtX}. Para el parámetro real \(s\), definamos
\[
 d_j(s)=\frac{K_je^{su_j}}{\sum_\ell K_\ell e^{su_\ell}}.
\]
Fijadas fracciones positivas \(\vartheta_{j,c}\) de suma uno para cada
\(j\), el refinamiento heredado es
\[
 \delta_{j,c}(s)=\vartheta_{j,c}d_j(s),\qquad
 u_{j,c}=u_j.
\]
Producir primero los pesos refinados \(K_j\vartheta_{j,c}\) y
deformarlos da la misma familia: los hijos tienen idéntico exponencial y
sus coeficientes suman \(K_j\). Para todo \(s\),
\begin{equation}
 \operatorname{Schur}_{I}(L_{\widetilde d(s)})=L_{d(s)},
 \qquad H^*L_{\widetilde d(s)}H=L_{d(s)}.
 \label{eq:schur-flujo-k}
\end{equation}
En esta comparación padre--hijos, \(H\) es constante: sus coeficientes
son las sumas parciales de \(\vartheta_{j,c}\). En cambio, la
interpolación desde los dos extremos globales usa
\(t_i(s)=\sum_{j\le i}d_j(s)\) y cambia con el parámetro. Ambas
afirmaciones corresponden a restricciones diferentes de la misma red.

La normalización \(\sum_jd_j(s)=1\) hace que el operador global de dos
extremos permanezca igual a
\(\left(\begin{smallmatrix}1&-1\\-1&1\end{smallmatrix}\right)\).
La matriz de todos los nodos \(L_{d(s)}\) sí conserva la deformación y
la recupera mediante \eqref{eq:recuperacion-d-rigidez}. La pérdida de
información en el lector de dos extremos queda identificada exactamente:
retiene la longitud total y sus flujos, mientras el lector nodal retiene
todas las separaciones. El archivo de detalles proporciona el complemento
necesario para reconstruir el observable refinado.

\section{Reconstrucción analítica de la energía y memoria del refinamiento}
\label{sec:limite-energetico-k}

Las separaciones positivas construidas a partir del registro dodecafásico
determinan una partición ordenada y, por la sección anterior, una forma
de Dirichlet sobre sus valores nodales. Su refinamiento posee dos
operaciones complementarias: la interpolación armónica conserva la energía
de los datos heredados, mientras cada detalle registra la variación
incorporada entre esos datos. El paso a profundidad ilimitada consiste en
reunir esta familia compatible mediante una norma que controle la suma de
sus detalles. El resultado concierne a la realización analítica
unidimensional de la energía; conserva como antecedente la construcción
HMT del registro y de la estructura discreta conjunta del continuo.

\subsection{Particiones anidadas y espacios de interpolación}

Fijemos la partición inicial
\(P_0=\{0=t_0<t_1<\cdots<t_{12}=1\}\), cuyas longitudes son
\(d_j=t_j-t_{j-1}=K_j/\sum_iK_i>0\). Sea
\((P_n)_{n\geq0}\) una sucesión de particiones finitas anidadas de
\([0,1]\), obtenidas por subdivisión de los intervalos anteriores.
Escribiremos
\[
 |P_n|=\max\{b-a:[a,b]\text{ es un intervalo de }P_n\},
 \qquad \lim_{n\to\infty}|P_n|=0.
\]
La subdivisión reiterada de cada intervalo en un número fijo de hijos
iguales, al menos dos, satisface esta condición. Para otras reglas, la
contracción de la malla es una hipótesis explícita de este lector.
Las etiquetas de los hijos y sus datos de acarreo acompañan la
subdivisión y conservan la procedencia del refinamiento.

Sea \(V_n\) el espacio complejo de funciones continuas sobre \([0,1]\)
que son afines en cada intervalo de \(P_n\). La energía y la norma
energética se definen por
\begin{equation}
 \mathcal E_n(f)
 =\sum_{[a,b]\in P_n}\frac{|f(b)-f(a)|^2}{b-a}
 =\int_0^1|f'(x)|^2\,dx,
 \qquad
 \|f\|_{\mathcal E}^2=|f(0)|^2+\mathcal E_n(f).
 \label{eq:energia-afines-k}
\end{equation}
La igualdad integral resulta de que la derivada de una función afín es
constante en cada intervalo. La energía coincide así exactamente con
la forma nodal ya construida. La inclusión \(V_n\subset V_{n+1}\)
interpreta la misma función afín en una partición más fina; conserva
tanto \(f(0)\) como la energía y es isométrica para
\(\|\cdot\|_{\mathcal E}\).

El valor inicial tiene una función precisa: fija la componente constante
que la forma de Dirichlet identifica en su núcleo. La energía controla
las diferencias, y el par formado por el valor inicial y esas diferencias
determina el observable. En el espacio de Sobolev
\(H^1([0,1];\mathbb C)\), cuyos elementos tienen representante
absolutamente continuo con derivada débil en \(L^2\), esta norma es
equivalente a la norma usual. En efecto,
\(f(x)=f(0)+\int_0^xf'(t)\,dt\) controla la norma de \(f\) a
partir de \(|f(0)|\) y \(\|f'\|_{L^2}\); la integración de la
misma identidad controla a su vez \(|f(0)|\) mediante
\(\|f\|_{L^2}+\|f'\|_{L^2}\).

\subsection{Reconstrucción de una función a partir de sus valores compatibles}

\begin{theorem}[Reconstrucción energética y completación de las interpolaciones]
\label{thm:limite-energia-h1-k}
Sea \(f_n\in V_n\) una familia compatible por restricción a los
nodos heredados:
\[
 f_{n+1}(a)=f_n(a)
 \quad\text{para todo nodo }a\text{ de }P_n.
\]
Si \(\sup_n\mathcal E_n(f_n)<\infty\), existe una única función
\(f\in H^1([0,1];\mathbb C)\), considerada en su representante
continuo, cuyos valores en todos esos nodos son los prescritos. Además,
\(f_n\to f\) uniformemente y en \(H^1\), y
\begin{equation}
 \int_0^1|f'(x)|^2\,dx
 =\lim_{n\to\infty}\mathcal E_n(f_n).
 \label{eq:limite-energia-h1-k}
\end{equation}
Recíprocamente, la interpolación nodal de cualquier elemento de \(H^1\)
produce una familia de esta clase. Por tanto, la completación de
\(\bigcup_n V_n\) en la norma energética es \(H^1([0,1];\mathbb C)\).
\end{theorem}

\begin{proof}
Pongamos \(g_n=f_n'\in L^2([0,1];\mathbb C)\). Para
\(m\geq n\), la compatibilidad de los valores en los extremos de cada
intervalo \([a,b]\) de \(P_n\) proporciona
\[
 \frac1{b-a}\int_a^bg_m(x)\,dx
 =\frac{f_m(b)-f_m(a)}{b-a}
 =\frac{f_n(b)-f_n(a)}{b-a}
 =g_n\big|_{(a,b)}.
\]
Sea \(W_n\) el espacio de funciones de \(L^2\) constantes en cada
intervalo de \(P_n\), y sea \(Q_n:L^2\to W_n\) la proyección
ortogonal dada por los promedios sobre dichos intervalos. La identidad
anterior afirma \(Q_ng_m=g_n\). En consecuencia,
\(g_m-g_n\) es ortogonal a \(g_n\) y
\begin{equation}
 \|g_m-g_n\|_{L^2}^2
 =\mathcal E_m(f_m)-\mathcal E_n(f_n).
 \label{eq:pitagoras-gradientes-k}
\end{equation}
Las energías son crecientes y acotadas. La identidad de Pitágoras
prueba que \((g_n)\) es una sucesión de Cauchy en \(L^2\), con
límite \(g\). Definamos
\[
 f(x)=f_0(0)+\int_0^xg(t)\,dt.
\]
Esta función pertenece a \(H^1\), tiene derivada débil \(g\) y
satisface, por la desigualdad de Cauchy--Schwarz,
\begin{equation}
 \sup_{x\in[0,1]}|f_n(x)-f(x)|
 \leq\|g_n-g\|_{L^2}\longrightarrow0.
 \label{eq:convergencia-uniforme-energia-k}
\end{equation}
La convergencia uniforme conserva cada valor nodal prescrito; junto con
la convergencia de las derivadas en \(L^2\), da la convergencia en
\(H^1\). La convergencia de las normas de las derivadas prueba
\eqref{eq:limite-energia-h1-k}. La unión de los nodos es densa porque
\(|P_n|\to0\); dos representantes continuos con los mismos valores
en esa unión coinciden. Se obtiene la unicidad.

Para la recíproca, tomemos \(f\in H^1\) en su representante
continuo y sea \(f_n\) su interpolación en los nodos de \(P_n\).
Su derivada es \(g_n=Q_nf'\). Las proyecciones \(Q_n\) son
contractivas en \(L^2\). Para una función continua \(v\), el
error \(|Q_nv-v|\) está acotado por su módulo de continuidad
evaluado en \(|P_n|\), que tiende a cero. La densidad de las
funciones continuas en \(L^2\) y la contractividad implican
\(Q_ng\to g\) para todo \(g\in L^2\): dado un aproximante
continuo \(v\), se emplea
\[
 \|Q_ng-g\|_{L^2}
 \leq 2\|g-v\|_{L^2}+\|Q_nv-v\|_{L^2}.
\]
Aplicando esta convergencia a \(g=f'\), se obtienen la convergencia
de las derivadas y, mediante \eqref{eq:convergencia-uniforme-energia-k},
las de las funciones. La contractividad acota sus energías por
\(\|f'\|_{L^2}^2\), y la convergencia de las normas prueba de
nuevo la identidad del límite. Toda función de \(H^1\) es, por
tanto, límite energético de elementos de \(\bigcup_nV_n\).
La equivalencia de normas y la completitud de \(H^1\) identifican
la completación enunciada.
\end{proof}

El argumento describe concretamente qué se conserva al pasar de la red
a la función. Cada derivada finita registra el promedio de todas las
derivadas posteriores en su intervalo; al refinar, se incorporan
componentes ortogonales que mantienen ese promedio. La energía acotada
controla la suma total de esas incorporaciones. El valor continuo se
recupera entonces integrando la derivada límite y conservando el valor
inicial. La realización analítica reúne, por esta vía explícita, la
información compatible de los niveles finitos.

\subsection{Descomposición ortogonal de la memoria por escalas}

Sea \(H_n\) la interpolación armónica de los valores de \(P_n\)
sobre los nodos de \(P_{n+1}\), interpretada como función de
\(V_{n+1}\), y definamos el detalle
\[
 z_{n+1}=f_{n+1}-H_nf_n.
\]
Por \eqref{eq:detalle-energia-k}, el detalle se anula en los nodos
heredados. Su derivada tiene promedio cero en cada intervalo antiguo
y es ortogonal a \(W_n\). Para niveles distintos, la derivada de
un detalle anterior pertenece a ese espacio de funciones constantes;
las derivadas de los detalles son, por tanto, ortogonales entre sí.
La misma ortogonalidad vale frente a \(f_0'\).

\begin{proposition}[Identidad de energía y recuperación de los detalles]
\label{prop:memoria-energetica-escalas-k}
Para toda profundidad finita \(N\),
\begin{equation}
 \mathcal E_N(f_N)
 =\mathcal E_0(f_0)
 +\sum_{n=0}^{N-1}\mathcal E_{n+1}(z_{n+1}).
 \label{eq:energia-telescopica-detalles-k}
\end{equation}
Si la familia satisface las hipótesis del
Teorema~\ref{thm:limite-energia-h1-k}, la serie de la derecha
converge y su suma es \(\int_0^1|f'|^2\). Los datos \(f_0\)
y todos los detalles determinan la familia compatible completa y su
realización analítica.
\end{proposition}
\begin{proof}
La identidad finita resulta de sumar
\(\mathcal E_{n+1}(f_{n+1})
=\mathcal E_n(f_n)+\mathcal E_{n+1}(z_{n+1})\), que es la
descomposición ortogonal de la sección anterior. Los sumandos son
no negativos y las sumas parciales convergen por
\eqref{eq:limite-energia-h1-k}. La reconstrucción finita se efectúa
recursivamente mediante \(f_{n+1}=H_nf_n+z_{n+1}\); el teorema
precedente determina después la función límite. Recíprocamente, una
función de \(H^1\) determina sus interpolaciones y, por diferencia,
todos estos detalles. Cada flecha conserva así sus datos propios.
\end{proof}

La eliminación variacional de nodos interiores selecciona el detalle
nulo en el nivel eliminado. La conservación de los detalles permite
recuperar cualquier configuración compatible de energía finita. La
respuesta efectiva y la configuración completa poseen, de esta manera,
una relación precisa: la primera registra el mínimo sobre las variables
interiores; la segunda añade las componentes ortogonales y su energía.

\subsection{Transporte de la realización analítica bajo el flujo positivo}

La dirección \(u=P_3P_{11}K\), recuperada antes de la evaluación
energética \cite{hmtX}, determina la familia de separaciones de
\eqref{eq:schur-flujo-k}. Su parámetro \(s\in\mathbb R\)
describe esta deformación analítica. Escribamos
\[
 d_j(s)=\frac{K_je^{su_j}}{\sum_\ell K_\ell e^{su_\ell}},
 \qquad
 t_i(s)=\sum_{j=1}^id_j(s),
 \qquad
 a_j(s)=\frac{d_j(s)}{d_j(0)}>0.
\]
Sea \(F_s:[0,1]\to[0,1]\) la aplicación que envía cada
\(t_i(0)\) a \(t_i(s)\) y es afín en los intervalos iniciales.
Es un homeomorfismo creciente, de pendiente \(a_j(s)\) en el
intervalo \([t_{j-1}(0),t_j(0)]\). Su inversa es también afín
por intervalos y ambas aplicaciones son Lipschitz.

Supongamos que cada subdivisión hereda las fracciones de longitud y la
dirección de su intervalo padre, como en la sección anterior. Todos
los nodos refinados se transportan entonces mediante el mismo \(F_s\).
Denotemos sus particiones por \(P_n(s)=F_s(P_n(0))\) y sus
espacios de interpolación por \(V_n(s)\). La aplicación
\[
 U_s f=f\circ F_s^{-1}
\]
transporta \(V_n(0)\) a \(V_n(s)\), preserva los valores
nodales correspondientes y prolonga esta acción a \(H^1\).

\begin{theorem}[Equivalencia energética y conmutación con el refinamiento]
\label{thm:transporte-h1-flujo-k}
Para todo \(f\in H^1([0,1];\mathbb C)\),
\begin{equation}
 \int_0^1|(U_sf)'(y)|^2\,dy
 =\sum_{j=1}^{12}\frac1{a_j(s)}
 \int_{t_{j-1}(0)}^{t_j(0)}|f'(x)|^2\,dx.
 \label{eq:transporte-energia-h1-k}
\end{equation}
Para \(s\) en cualquier intervalo compacto, las normas energéticas
antes y después del transporte son uniformemente equivalentes y
\(|P_n(s)|\to0\) uniformemente en ese compacto. La interpolación
armónica y la restricción a nodos heredados conmutan con \(U_s\);
la reconstrucción y la minimización de los detalles se transportan
tanto a profundidad finita como en la realización \(H^1\).
\end{theorem}

\begin{proof}
En el intervalo inicial \(j\), la regla de la cadena y el cambio
de variable \(y=F_s(x)\) proporcionan
\((U_sf)'(F_s(x))=f'(x)/a_j(s)\) y \(dy=a_j(s)\,dx\).
La integración y la suma sobre los intervalos prueban
\eqref{eq:transporte-energia-h1-k}. La fórmula vale para los
representantes absolutamente continuos, y su lado derecho es finito.

Fijado un intervalo compacto de parámetros, la continuidad y positividad
de los doce factores dan constantes \(0<a_-\leq a_+<\infty\)
tales que \(a_-\leq a_j(s)\leq a_+\) para todo \(j\) y
todo \(s\) de ese compacto. Por tanto,
\[
 \frac1{a_+}\mathcal E(f)
 \leq\mathcal E(U_sf)
 \leq\frac1{a_-}\mathcal E(f),
 \qquad
 |P_n(s)|\leq a_+|P_n(0)|.
\]
Como \((U_sf)(0)=f(0)\), se obtiene además
\[
 \min\{1,a_+^{-1}\}\|f\|_{\mathcal E}^2
 \leq\|U_sf\|_{\mathcal E}^2
 \leq\max\{1,a_-^{-1}\}\|f\|_{\mathcal E}^2.
\]
Estas cotas prueban la equivalencia uniforme y la conservación de la
condición de malla decreciente.

En cada arista padre, los coeficientes de interpolación son las sumas
parciales de las fracciones heredadas, independientes de \(s\).
Si \(H_n(s)\) denota su interpolación, se tiene
\(U_sH_n(0)=H_n(s)U_s\) en los espacios correspondientes.
La restricción conmuta porque los nodos se transportan con sus valores.
La descomposición en función armónica y detalle se conserva, y la
identidad de Schur de \eqref{eq:schur-flujo-k} preserva su
caracterización variacional. Finalmente, la continuidad acotada de
\(U_s\) permite pasar al límite \(H^1\). En particular, una
familia reconstruida antes del transporte converge después a \(U_sf\),
con los mismos valores nodales transportados.
\end{proof}

La ley \eqref{eq:transporte-energia-h1-k} identifica los factores
métricos exactos de la deformación: el cambio de longitud de cada
intervalo modifica inversamente su contribución energética. La
equivalencia de normas conserva las configuraciones de energía finita,
y el archivo de detalles conserva su reconstrucción. Se obtiene así una
realización analítica completa de este lector del registro positivo.
Su compatibilidad con las formas canónicas se articula sobre las
subdivisiones de un dominio fijo: la suma de formas conserva el objeto
meromorfo, mientras la reducción variacional y los detalles conservan
la respuesta energética y el observable refinado. La identificación
con formas amplituhedrales de rango superior conserva las condiciones
específicas de cobertura, grado y correspondencia de residuos
establecidas para aquella realización.

\section{Transportes nonádicos y conservación de la memoria}
\label{sec:nonadic}

La reconstrucción de la sección~\ref{sec:memory} proporciona un enlace
efectivo entre la memoria del transporte y la geometría del registro
dodecafásico. El registro generado por APP--TRIT--TPK determina sus
lecturas de incidencia; las dos compresiones conservan, en componentes
separadas, la publicación terminal y sus complementos; la reunión de
las memorias con la carga restituye las doce coordenadas. La dirección
geométrica y el proyector dimensional se obtienen después de esa
restitución. El desarrollo siguiente establece el balance operatorio
que sostiene esta cadena y su comportamiento bajo refinamiento,
realización espectral y cambio de métrica
\cite{HMT-VIII,hmtX}.

\subsection{Calendario de nueve fases y descomposición ortogonal}

Sea \(\mathcal H\) un espacio de Hilbert complejo y sea
\(C\in\mathcal B(\mathcal H)\). La suma ortogonal de nueve copias,
\(\mathcal H^{\oplus9}\), conserva por separado las nueve posiciones
del calendario. Definimos
\begin{equation}
 \begin{split}
 S_C(v_0,\ldots,v_8)&=(Cv_8,v_0,\ldots,v_7),\\
 J:\mathcal H\longrightarrow\mathcal H^{\oplus9},
 \qquad Jv&=\frac13(v,\ldots,v).
 \end{split}
 \label{nt:calendar}
\end{equation}
La normalización \(1/3=9^{-1/2}\) da \(J^*J=I\). El operador
\(P=JJ^*\) es la proyección ortogonal sobre la sección uniforme;
su acción sustituye las nueve componentes por su media. Cada
componente atraviesa una vez la unión que contiene \(C\) durante una
vuelta, de modo que
\begin{equation}
 S_C^9=I_9\otimes C.
 \label{nt:full-return}
\end{equation}
El retorno de la posición del calendario realiza así una acción de
\(C\) sobre la fibra de memoria.

\begin{definition}
La compresión uniforme del transporte y su complemento son
\begin{equation}
 T_C=J^*S_CJ=\frac{8I+C}{9},
 \qquad
 \eta_C=(I-P)S_CJ.
 \label{nt:compression}
\end{equation}
La primera aplicación tiene codominio \(\mathcal H\); la segunda
tiene codominio \(J(\mathcal H)^\perp\subset\mathcal H^{\oplus9}\).
\end{definition}

Para \(z=(C-I)v\), el complemento se escribe
\begin{equation}
 \eta_Cv=\frac1{27}(8z,-z,-z,-z,-z,-z,-z,-z,-z).
 \label{nt:complement}
\end{equation}
La suma de sus componentes es cero. Los coeficientes de la
descomposición proceden, por tanto, del promedio de ocho
continuaciones idénticas y una continuación transportada.

\begin{proposition}[Balance de la compresión y de la unión de memoria]
Para todo \(C\in\mathcal B(\mathcal H)\),
\begin{align}
 \eta_C^*\eta_C
   &=\frac8{81}(I-C)^*(I-C),\label{nt:eta-gram}\\
 I-T_C^*T_C
   &=\eta_C^*\eta_C+\frac19(I-C^*C).\label{nt:general-balance}
\end{align}
Si \(C^*C=I\), la aplicación
\(v\mapsto(T_Cv,\eta_Cv)\) es isométrica.
\end{proposition}

\begin{proof}
La expresión~\eqref{nt:complement} da
\[
 \|\eta_Cv\|^2
 =\frac{8^2+8}{27^2}\|(C-I)v\|^2
 =\frac8{81}\|(C-I)v\|^2.
\]
La polarización prueba~\eqref{nt:eta-gram}. Por otra parte,
\[
 T_C^*T_C
 =\frac{64I+8C+8C^*+C^*C}{81}.
\]
Sumando \eqref{nt:eta-gram} se obtiene
\[
 T_C^*T_C+\eta_C^*\eta_C=\frac{8I+C^*C}{9}.
\]
Restar este resultado de \(I\) da~\eqref{nt:general-balance}.
Cuando \(C^*C=I\), la suma de los dos productos de Gram es la
identidad, lo que conserva todos los productos interiores.
\end{proof}

El término \(\eta_C^*\eta_C\) mide la dispersión entre las nueve
componentes. El término \((I-C^*C)/9\) registra, por separado,
la variación de la norma en la unión de memoria. La realización
unitaria sitúa toda la variación visible en el primer término y
conserva la norma en la suma ortogonal completa.

\begin{proposition}[Conservación de un observable]
Sean \(C\) unitario y \(A=A^*\in\mathcal B(\mathcal H)\) tales que
\(C^*AC=A\). Entonces
\begin{equation}
 A=T_C^*AT_C+
   \eta_C^*(I_9\otimes A)\eta_C
  =T_C^*AT_C+\frac8{81}(I-C)^*A(I-C).
 \label{nt:observable-balance}
\end{equation}
Para \(A\succeq0\), ambos términos de la suma son positivos.
\end{proposition}

\begin{proof}
La fórmula del complemento, aplicada al producto sesquilineal
definido por \(A\), prueba la segunda igualdad. La expansión de
los dos términos da
\[
 \frac1{81}\bigl(
 64A+8AC+8C^*A+C^*AC
 +8A-8AC-8C^*A+8C^*AC\bigr)
 =\frac{8A+C^*AC}{9}=A.
\]
La positividad se conserva bajo una aplicación de la forma
\(B\mapsto L^*BL\).
\end{proof}

\subsection{Iteración de la compresión y conservación del sector fijo}

Fijamos en esta subsección un operador unitario \(C\) y escribimos
\(T=T_C\), \(\eta=\eta_C\). El transporte \(S_C^N\) mantiene todas
las fases durante \(N\) pasos. La potencia \(T^N\), en cambio,
aplica sucesivamente la compresión uniforme; en cada etapa
\(\eta T^jv\) conserva el complemento de la publicación anterior.
La distinción queda expresada por los siguientes operadores.

\begin{proposition}[Registro isométrico a horizonte finito]
Para todo entero \(N\geq1\), la aplicación
\begin{equation}
 V_Nv=(T^Nv,\eta v,\eta Tv,\ldots,\eta T^{N-1}v)
 \label{nt:finite-record}
\end{equation}
de \(\mathcal H\) a
\(\mathcal H\oplus(\mathcal H^{\oplus9})^{\oplus N}\)
es isométrica. Si \(A\) satisface las hipótesis de
\eqref{nt:observable-balance}, entonces
\begin{equation}
 A=T^{*N}AT^N+
 \sum_{j=0}^{N-1}T^{*j}\eta^*(I_9\otimes A)\eta T^j.
 \label{nt:telescopy}
\end{equation}
\end{proposition}

\begin{proof}
Multiplicar~\eqref{nt:observable-balance} por \(T^{*j}\) y
\(T^j\) da
\[
 T^{*j}AT^j-T^{*(j+1)}AT^{j+1}
 =T^{*j}\eta^*(I_9\otimes A)\eta T^j.
\]
La suma telescópica prueba~\eqref{nt:telescopy}. Para \(A=I\),
el miembro derecho es el producto de Gram de \(V_N\).
El paso de \(N\) a \(N+1\) sustituye únicamente \(T^Nv\) por
\((T^{N+1}v,\eta T^Nv)\), conservando todas las componentes
de memoria previamente producidas.
\end{proof}

\begin{theorem}[Límite de la compresión y memoria ilimitada]
Sea \(P_{\mathrm{fix}}\) la proyección ortogonal sobre
\(\ker(I-C)\). Entonces
\begin{align}
 \operatorname*{s\!-\!lim}_{N\to\infty}T^N
   &=P_{\mathrm{fix}},\label{nt:strong-limit}\\
 I&=P_{\mathrm{fix}}+
 \sum_{j=0}^{\infty}T^{*j}\eta^*\eta T^j.\label{nt:infinite-record}
\end{align}
La serie converge fuertemente. Para todo observable acotado
autoadjunto que satisfaga \(C^*AC=A\), se conserva además
\begin{equation}
 A=P_{\mathrm{fix}}AP_{\mathrm{fix}}+
 \sum_{j=0}^{\infty}
 T^{*j}\eta^*(I_9\otimes A)\eta T^j.
 \label{nt:infinite-observable}
\end{equation}
\end{theorem}

\begin{proof}
El teorema espectral del operador unitario realiza \(C\) como
la función coordenada \(z\) sobre la circunferencia unidad.
La compresión corresponde a \(t(z)=(8+z)/9\), y
\begin{equation}
 1-|t(z)|^2=\frac8{81}|1-z|^2,
 \qquad |z|=1.
 \label{nt:circular-defect}
\end{equation}
Así, \(t(z)^N\) converge a cero en cada \(z\ne1\) y conserva
el valor uno en \(z=1\). La cota \(|t(z)^N|\leq1\) permite
aplicar convergencia dominada a cada medida espectral vectorial.
Se obtiene~\eqref{nt:strong-limit} y, por el mismo argumento,
\(T^{*N}\to P_{\mathrm{fix}}\) fuertemente. La acotación de
\(A\), junto con \(\|T^N\|\leq1\), da
\(T^{*N}AT^N\to P_{\mathrm{fix}}AP_{\mathrm{fix}}\).
El límite de~\eqref{nt:telescopy} prueba las dos identidades.
\end{proof}

El registro
\[
 V_\infty v=(P_{\mathrm{fix}}v,\eta v,\eta Tv,\eta T^2v,\ldots)
\]
es, por consiguiente, una isometría. El sector fijo conserva la
parte que persiste en la publicación terminal; las demás
componentes permanecen en la memoria de los complementos.

Para el desplazamiento bilateral
\(C e_n=e_{n+1}\) en \(\ell^2(\mathbb Z)\), una sucesión fija
es constante. La condición de cuadrado-sumabilidad fuerza su
anulación, por lo que \(P_{\mathrm{fix}}=0\). Toda la norma
queda entonces representada en los complementos. El espectro de
este desplazamiento es la circunferencia completa, y el máximo
espectral de \(|t(z)|^N\) es uno; de ahí
\(\|T^N\|=1\) para cada \(N\). La convergencia fuerte conserva
este comportamiento no uniforme.

Para una permutación cíclica de \(d\) posiciones, el sector fijo
es la recta de los vectores uniformes. Si la permutación tiene
varias órbitas, cada órbita aporta su propia dirección uniforme.
Esta dimensión determina la parte terminal que debe conservarse.

\begin{proposition}[Compresión de un transporte de orden tres]
Sea \(C\) unitario y \(C^3=I\). Entonces
\begin{equation}
 P_{\mathrm{fix}}=\frac{I+C+C^2}{3},
 \qquad
 T_C^*T_C=P_{\mathrm{fix}}+
 \frac{19}{27}(I-P_{\mathrm{fix}}).
 \label{nt:order-three}
\end{equation}
\end{proposition}

\begin{proof}
El promedio de las tres potencias es el proyector ortogonal
sobre los vectores fijos. Como \(C^*=C^2\),
\[
 T_C^*T_C=\frac{65I+8(C+C^2)}{81}
 =\frac{57I+24P_{\mathrm{fix}}}{81},
\]
que coincide con~\eqref{nt:order-three}.
\end{proof}

En la realización de doce posiciones de
\(\mathbb R^{12}\), \((Sx)_i=x_{i+1}\) con índices módulo
doce, \(S^3\) tiene orden cuatro y \(S^4\) tiene orden tres.
Las cuatro órbitas de \(S^4\) dan
\(\operatorname{rank}P_{\mathrm{fix}}=4\).
La realización excepcional \(g_c\) del artículo~X satisface
\(g_c^3=I\) y \(\ker(I-g_c)=0\) en su propio portador;
en ella~\eqref{nt:order-three} se reduce a
\(T_{g_c}^*T_{g_c}=19I/27\)
\cite{hmtX}. El coeficiente y el sector fijo se transportan
conjuntamente con el dominio.

La contracción radial \(M=qC\), \(0<q<1\), posee otro registro
exacto:
\begin{equation}
 L_Nv=\bigl(\sqrt{1-q^2}\,v,
 \sqrt{1-q^2}\,qCv,\ldots,
 \sqrt{1-q^2}\,q^{N-1}C^{N-1}v,q^NC^Nv\bigr).
 \label{nt:radial-record}
\end{equation}
Su producto de Gram es la identidad porque
\((1-q^2)\sum_{j=0}^{N-1}q^{2j}+q^{2N}=1\).
El terminal tiene norma \(q^N\|v\|\), y el registro infinito
conserva el producto interior. El factor \(q=5/9\) del modo
diferencial central del artículo~VIII se aplica a una vuelta
completa; su reparto uniforme entre nueve pasos utiliza
\(q^{1/9}S_C\), cuya novena potencia es
\(q(I_9\otimes C)\) \cite{HMT-VIII}.

\subsection{Reducción del registro y conservación de la traza}

La dirección fija del complemento en las nueve componentes
permite una realización mínima en dos sumandos. Definamos
\[
 D_C=\frac{\sqrt8}{9}(I-C).
\]
Las identidades anteriores dan
\(T_C^*T_C+D_C^*D_C=I\) para \(C\) unitario.
El cambio de signo respecto de la expresión de
\(\eta_C\) corresponde a una orientación de la componente
complementaria y conserva su producto de Gram.

\begin{proposition}[Canal positivo de la compresión con memoria]
Para \(C\) unitario, la aplicación
\begin{equation}
 \mathcal Q_C(X)=T_CXT_C^*+D_CXD_C^*
              =\frac{8X+CXC^*}{9}
 \label{nt:channel}
\end{equation}
es completamente positiva en \(\mathcal B(\mathcal H)\) y
conserva la identidad. Su restricción a los operadores de
clase traza conserva la traza.
\end{proposition}

\begin{proof}
Al expandir los dos términos de la primera expresión,
los productos \(CX\) y \(XC^*\) se cancelan. Los coeficientes
restantes son \(72/81=8/9\) para \(X\) y \(9/81=1/9\)
para \(CXC^*\), lo que prueba~\eqref{nt:channel}.
Para cualquier entero \(m\geq1\), la ampliación a matrices
de orden \(m\) tiene la forma
\[
 X\longmapsto
 (I_m\otimes T_C)X(I_m\otimes T_C)^*
 +(I_m\otimes D_C)X(I_m\otimes D_C)^*.
\]
Cada sumando conserva positividad; por ello la aplicación es
completamente positiva. La segunda expresión de
\eqref{nt:channel} da \(\mathcal Q_C(I)=I\).
Si \(X\) es de clase traza, la invariancia de la traza bajo
conjugación unitaria implica
\(\operatorname{Tr}\mathcal Q_C(X)=\operatorname{Tr}X\).
\end{proof}

La aplicación \eqref{nt:channel} se obtiene igualmente
preparando la isometría
\(v\mapsto T_Cv\otimes e_0+D_Cv\otimes e_1\) y efectuando
la traza parcial sobre \(\mathbb C^2\).
La conservación del vector completo pertenece a la isometría;
la conservación de la traza del estado reducido pertenece a
\(\mathcal Q_C\). Ambas afirmaciones proceden del mismo
balance ortogonal.

En la prolongación nonádica de álgebras finitas aparece una
conservación complementaria. Para
\(\mathcal A_d=M_{9^d}(\mathbb C)\), con
\(\tau_d(a)=9^{-d}\operatorname{Tr}a\), las aplicaciones
\[
 \iota_d(a)=a\otimes I_9,\qquad
 \jmath_{d,j}(a)=a\otimes p_j,
 \quad p_j=|e_j\rangle\langle e_j|
\]
satisfacen
\begin{align}
 \tau_{d+1}(\iota_d(a))&=\tau_d(a),
 &\iota_d(I)&=I,\label{nt:unital-trace}\\
 \tau_{d+1}(\jmath_{d,j}(a))&=\frac{\tau_d(a)}9,
 &\jmath_{d,j}(I)&=I\otimes p_j.\label{nt:corner-trace}
\end{align}
La demostración usa
\(\operatorname{Tr}(a\otimes b)=
\operatorname{Tr}(a)\operatorname{Tr}(b)\),
\(\operatorname{Tr}I_9=9\) y \(\operatorname{Tr}p_j=1\).
La primera inclusión conserva toda la fibra y la unidad;
la segunda conserva una continuación marcada y su soporte.

Esta distinción también se manifiesta en el árbol de
historias enriquecidas. Sean \(\mathcal H_n\) los prefijos
supervivientes de longitud \(n\), y sea \(d(h)\geq1\) el número de
emisiones salientes conservadas de \(h\). Cada emisión con multiplicidad
se representa por hijos marcados distintos; la suma siguiente recorre
ese conjunto de \(d(h)\) hijos.
La ley de pesos
\[
 \mu_{n+1}(h\epsilon)=\mu_n(h)\,p_h(\epsilon),
 \qquad p_h(\epsilon)=d(h)^{-1}
\]
conserva
\(\sum_{\epsilon}\mu_{n+1}(h\epsilon)=\mu_n(h)\).
Por tanto, en
\(\mathscr L_n=L^2(\mathcal H_n,\mu_n)\), el pullback
\[
 R_nf(h\epsilon)=f(h)
\]
es isométrico y conserva la función unidad. Su adjunto es
\[
 E_ng(h)=\sum_{\epsilon}p_h(\epsilon)g(h\epsilon),
 \qquad E_nR_n=I,\quad E_n1=1.
\]
En efecto, agrupar la suma de \(\|R_nf\|^2\) por padre
produce \(\sum_h\mu_n(h)|f(h)|^2\); el mismo agrupamiento
en el producto interior determina \(E_n\).
La memoria y las multiplicidades quedan presentes en el
dominio de las historias. Estos mapas constituyen la
realización funcional del refinamiento descrito en el
artículo~X \cite{hmtX}.

\subsection{Carta espectral y reciprocidad del radio}

El transporte anterior admite realizaciones espectrales
con dominios propios. La carta de Möbius de VIII actúa,
en coordenadas finitas, por
\begin{equation}
 \tau(z)=\frac{z-1}{z},
 \qquad z\in\mathbb C\setminus\{0\}.
 \label{nt:mobius}
\end{equation}
La inversa es \(z=(1-\tau)^{-1}\), para \(\tau\ne1\).
La extensión a la esfera de Riemann envía \(0\) a
\(\infty\) e \(\infty\) a \(1\).

\begin{proposition}[Espejo espectral y radio recíproco]
Sea \(z^\sharp=1-\overline z\). Para
\(z\notin\{0,1\}\),
\begin{equation}
 \tau(z^\sharp)=\frac1{\overline{\tau(z)}}.
 \label{nt:mirror}
\end{equation}
Si \(z=\sigma+it\), entonces
\begin{equation}
 |\tau(z)|^2
 =\frac{(\sigma-1)^2+t^2}{\sigma^2+t^2};
 \qquad
 |\tau(z)|=1\ \Longleftrightarrow\ \sigma=\frac12.
 \label{nt:critical-circle}
\end{equation}
\end{proposition}

\begin{proof}
La sustitución directa da
\[
 \tau(1-\overline z)
 =\frac{-\overline z}{1-\overline z}
 =\frac{\overline z}{\overline z-1}
 =\frac1{\overline{(z-1)/z}}.
\]
El cociente de módulos cuadrados produce
\eqref{nt:critical-circle}; la igualdad entre numerador y
denominador equivale a \(-2\sigma+1=0\).
\end{proof}

Escribiendo \(\tau(z)=r e^{i\theta}\), con \(r>0\), el
espejo preserva \(\theta\) y transforma \(r\) en \(r^{-1}\).
Su conjunto fijo en esta carta es la circunferencia unidad.
La conjugación \(z\mapsto\overline z\), por su parte,
produce \(\tau\mapsto\overline\tau\) y cambia la orientación
angular. Las dos involuciones conservan, respectivamente,
la coordenada angular y la coordenada radial.

Sobre \(z=1/2+i\gamma\), con \(\gamma\in\mathbb R\),
\begin{equation}
 \tau_\gamma=\frac{\gamma+i/2}{\gamma-i/2},
 \qquad
 \gamma=\frac{i}{2}\frac{\tau_\gamma+1}{\tau_\gamma-1}
       =\frac12\cot\frac{\theta}{2},
 \quad \tau_\gamma=e^{i\theta}.
 \label{nt:spectral-phase}
\end{equation}
La compactificación conserva la altura mediante la inversa;
el punto \(\tau=1\) corresponde al infinito.

Sea \(\mathcal D=C_c^\infty(\mathbb R;\mathbb C)\), con su topología
usual de funciones de prueba. Para la realización positiva de la forma
de Weil se utilizan expresamente su positividad semidefinida, su
invariancia por traslación y su continuidad en esa topología, en el
dominio demostrado por el lector espectral correspondiente. Si
\(\mathscr W(f,g)=\langle Ef,Eg\rangle\), su núcleo es
\(\mathcal N=\{f\in\mathcal D:\mathscr W(f,f)=0\}\), y se define
\[
 \mathcal H_{\mathscr W}
 =\overline{\mathcal D/\mathcal N}.
\]
Las traslaciones inducen un grupo unitario fuertemente
continuo \(U_t[f]=[f(\,\cdot-t)]\): la invariancia preserva la norma,
y la continuidad de traslaciones en \(\mathcal D\) y de la forma
prueba la continuidad fuerte sobre un subespacio denso y después en
toda la completación. Con la convención
\(U_t=e^{itH}\), su generador autoadjunto actúa sobre
las clases de pruebas por \(H[f]=[if']\). El Cayley de
esta realización es el operador acotado
\begin{equation}
 C_{\mathscr W}
 =(H+iI/2)(H-iI/2)^{-1}
 =I+i(H-iI/2)^{-1}.
 \label{nt:cayley-hilbert}
\end{equation}
El cálculo funcional lo hace unitario, pues el cociente
en~\eqref{nt:spectral-phase} tiene módulo uno para
\(\gamma\) real. La identidad de compresión se aplica
a \(C_{\mathscr W}\) en \(\mathcal H_{\mathscr W}\).
El espacio, el dominio del generador y la positividad
que permite su construcción son las premisas de esta
realización espectral \cite{HMT-VIII}.

La memoria bilateral actúa, en cambio, en
\(\ell^2(\mathbb Z)\), y las permutaciones \(S^3,S^4\)
actúan en \(\mathbb R^{12}\). Cada realización determina
su espectro, sus sectores fijos y la acción de una vuelta.
Los respectivos mapas de transporte especifican las
composiciones entre estos espacios.

\subsection{Cayley sobre pruebas y conservación de la forma aritmética}

La misma publicación aritmética admite un lector de Cayley
definido directamente sobre funciones de prueba. Fijemos
\(a>b>1/2\) y
\[
 \mathscr E_b=
 \left\{f\in C^\infty(\mathbb R):
   \sup_x e^{b|x|}|f^{(j)}(x)|<\infty
   \text{ para todo }j\geq0\right\}.
\]
Las funciones suaves de soporte compacto pertenecen a
\(\mathscr E_b\). Definimos
\begin{align}
 (R_a^+f)(x)&=\int_0^\infty e^{-at}f(x+t)\,dt,
 &(R_a^-f)(x)&=\int_0^\infty e^{-at}f(x-t)\,dt,\label{nt:resolvents}\\
 C_af&=f-2aR_a^+f,
 &C_a^{-1}f&=f-2aR_a^-f.\label{nt:test-cayley}
\end{align}
Cada derivada de \(R_a^\pm f\) satisface la cota del
correspondiente seminorma de \(f\), multiplicada por
\((a-b)^{-1}\). Esto sigue de
\(e^{b|x|}|f(x\pm t)|\leq M_0e^{bt}\).
Ambas aplicaciones preservan \(\mathscr E_b\).

Con la transformada
\(\widehat f(\xi)=\int_{\mathbb R}f(x)e^{-i\xi x}\,dx\),
la aplicación de Fubini da
\begin{equation}
 \widehat{C_af}(\xi)
 =c_a(\xi)\widehat f(\xi),\qquad
 c_a(\xi)=\frac{\xi-ia}{\xi+ia}.
 \label{nt:cayley-multiplier}
\end{equation}
El multiplicador de la segunda fórmula
de~\eqref{nt:test-cayley} es \(c_a^{-1}\), lo que prueba
la inversa. Para \(\xi\in\mathbb R\),
\(|c_a(\xi)|=1\); Plancherel proporciona su extensión
unitaria a \(L^2(\mathbb R)\).

Para precisar la conservación aritmética, sean
\[
 \mathcal C_{f,g}(u)
 =\int_{\mathbb R}f(x)\overline{g(x-u)}\,dx,
 \qquad
 P_f^\pm=\int_{\mathbb R}f(x)e^{\pm x/2}\,dx.
\]
Las longitudes \(\ell_{p,k}=k\log p\) y los pesos
\(w_{p,k}=(\log p)p^{-k/2}\) son las publicaciones
primo--potencia del lector aritmético. La forma completa
del artículo~VIII se escribe
\begin{equation}
 \begin{split}
 \mathscr W(f,g)={}&c_\Gamma\mathcal C_{f,g}(0)\\
 &+\int_0^\infty\frac{e^{-u/2}}{1-e^{-2u}}
 \bigl[2\mathcal C_{f,g}(0)
       -\mathcal C_{f,g}(u)-\mathcal C_{f,g}(-u)\bigr]\,du\\
 &-\sum_{p,k}w_{p,k}
   \bigl[\mathcal C_{f,g}(\ell_{p,k})
        +\mathcal C_{f,g}(-\ell_{p,k})\bigr]
 +P_f^+\overline{P_g^-}+P_f^-\overline{P_g^+},
 \end{split}
 \label{nt:weil-complete}
\end{equation}
donde \(c_\Gamma=\psi(1/4)-\log\pi\), y
\(\psi=\Gamma'/\Gamma\) es la derivada logarítmica de la
función gamma. Estas funciones y constantes intervienen
en la realización analítica posterior del lector
aritmético del corpus.

La cota
\[
 |\mathcal C_{f,g}(u)|
 \leq M_fM_g\bigl(|u|+b^{-1}\bigr)e^{-b|u|}
\]
controla la suma prima por una serie convergente del tipo
\(\sum_{n\geq2}(\log n)(1+\log n)n^{-b-1/2}\).
La diferencia simétrica de correlaciones es \(O(u^2)\)
en el origen, mientras que el peso gamma es \(O(u^{-1})\).
En infinito ese peso decrece exponencialmente.
El margen \(b>1/2\) asegura asimismo la convergencia de
los dos lectores polares. Así queda fijado el dominio
de todos los términos de~\eqref{nt:weil-complete}.

\begin{proposition}[Invariancia aritmética y balance nonádico]
Para \(f,g\in\mathscr E_b\),
\begin{equation}
 \mathscr W(C_af,C_ag)=\mathscr W(f,g).
 \label{nt:weil-invariance}
\end{equation}
Con \(\mathcal T_a=(8I+C_a)/9\), se cumple
\begin{equation}
 \mathscr W(f,g)
 =\mathscr W(\mathcal T_af,\mathcal T_ag)
 +\frac8{81}\mathscr W((I-C_a)f,(I-C_a)g).
 \label{nt:weil-balance}
\end{equation}
\end{proposition}

\begin{proof}
El multiplicador \(c_a\) es unitario en la recta real y
conmuta con las traslaciones. En consecuencia,
\(\mathcal C_{C_af,C_ag}(u)=\mathcal C_{f,g}(u)\)
para todo \(u\). Los lectores polares se transforman
por Fubini según
\[
 P_{C_af}^+
 =-\frac{a-1/2}{a+1/2}P_f^+,\qquad
 P_{C_af}^-
 =-\frac{a+1/2}{a-1/2}P_f^-.
\]
Sus factores son reales y recíprocos, de modo que
cada producto polar cruzado permanece inalterado.
La reunión de los términos de
\eqref{nt:weil-complete} prueba~\eqref{nt:weil-invariance}.
La expansión sesquilineal del segundo miembro de
\eqref{nt:weil-balance} cancela los términos mixtos y
produce
\(\bigl(8\mathscr W(f,g)+
\mathscr W(C_af,C_ag)\bigr)/9\), igual a
\(\mathscr W(f,g)\).
\end{proof}

La iteración finita conserva la forma completa:
\[
 \mathscr W(f,g)
 =\mathscr W(\mathcal T_a^Nf,\mathcal T_a^Ng)
 +\frac8{81}\sum_{j=0}^{N-1}
 \mathscr W((I-C_a)\mathcal T_a^jf,
            (I-C_a)\mathcal T_a^jg).
\]
Esta identidad sesquilineal conserva simultáneamente
gamma, primos y términos polares. Su evaluación en una
realización positiva permite interpretarla como reparto
de energía; su validez algebraica precede a esa
interpretación. El margen \(a>b>1/2\) pertenece al
dominio exponencial de pruebas de esta construcción.
El Cayley~\eqref{nt:cayley-hilbert} con parámetro \(1/2\)
pertenece al dominio hilbertiano de su generador.

\subsection{Covariancia métrica del flujo dodecafásico}

La reconstrucción de la sección~\ref{sec:memory} determina
\(u_K\in\mathbb R^{12}\), con
\(\mathbf1^{\mathsf T}u_K=0\) y \(u_K\ne0\).
Escribimos \(A_K=\operatorname{diag}(u_K)\), con la dirección y la
normalización declaradas en \eqref{eq:u-main}, y
\[
 g_s=e^{sA_K}\in\operatorname{GL}(12,\mathbb R).
\]
La traza nula de \(A_K\) implica \(\det g_s=1\).
Cada \(g_s\) transporta los pesos y las coordenadas
geométricas según la realización correspondiente.
Para conservar también el balance de memoria entre
cartas hay que transportar el producto interior.

\begin{proposition}[Transporte métrico de compresión y complemento]
Sean \(V\) un espacio hermítico finito, \(C_0\) unitario
y \(g:V\to V\) invertible. Escribimos
\(g^{-*}=(g^{-1})^*\) y definimos
\begin{equation}
 C_g=gC_0g^{-1},\qquad
 G_g=g^{-*}g^{-1},\qquad
 \langle v,w\rangle_g=\langle G_gv,w\rangle.
 \label{nt:metric-chart}
\end{equation}
Entonces \(C_g^*G_gC_g=G_g\), y
\begin{align}
 T_{C_g}g&=gT_{C_0},\label{nt:metric-naturality}\\
 \eta_{C_g}g&=g^{\oplus9}\eta_{C_0},\label{nt:eta-naturality}\\
 G_g&=T_{C_g}^*G_gT_{C_g}
 +\frac8{81}(I-C_g)^*G_g(I-C_g).\label{nt:metric-balance}
\end{align}
\end{proposition}

\begin{proof}
La positividad de \(G_g\) se obtiene de
\(\langle G_gv,v\rangle=\|g^{-1}v\|^2\).
Sustituyendo~\eqref{nt:metric-chart},
\[
 C_g^*G_gC_g
 =g^{-*}C_0^*g^*g^{-*}g^{-1}gC_0g^{-1}
 =g^{-*}C_0^*C_0g^{-1}=G_g.
\]
La identidad \(C_gg=gC_0\) da
\eqref{nt:metric-naturality} al sumar \(8g\) y dividir
por nueve. También da
\((C_g-I)g=g(C_0-I)\); su sustitución en
\eqref{nt:complement} prueba~\eqref{nt:eta-naturality}.
Finalmente, la expansión de
\eqref{nt:observable-balance}, con \(A=G_g\) y
\(C_g^*G_gC_g=G_g\), prueba
\eqref{nt:metric-balance}.
\end{proof}

Aplicando la proposición a \(g=g_s\) y a los operadores
\(C_0=S^3,S^4\) de la realización dodecafásica, las
compresiones, las memorias y sus productos interiores
se transportan por cuadrados conmutativos explícitos.
La identificación \(g_s:(V,\langle\cdot,\cdot\rangle)
\to(V,\langle\cdot,\cdot\rangle_{g_s})\) es isométrica.
Por tanto, la fórmula transporta también los sectores
fijos y los registros a toda profundidad:
\[
 P_{\mathrm{fix},s}
 =g_sP_{\mathrm{fix},0}g_s^{-1}.
\]
El proyector del miembro derecho es ortogonal respecto
de \(G_{g_s}\).

La condición de volumen y la condición métrica tienen
efectos diferenciados. Para la matriz racional
\[
 g=\begin{pmatrix}2&0\\0&1/2\end{pmatrix},
 \qquad \det g=1,
\]
la sustitución directa en la compresión da
\[
 \frac{8I+g}{9}
 =\begin{pmatrix}10/9&0\\0&17/18\end{pmatrix},
 \qquad
 T_g^*T_g+\frac8{81}(I-g)^*(I-g)
 =\frac{8I+g^*g}{9}.
\]
La primera coordenada se expande en la métrica
euclídea. El balance~\eqref{nt:metric-balance}
conserva en cambio la unidad métrica cuando se
transporta un operador unitario \(C_0\) junto con
su producto interior. El determinante controla
el volumen; \(G_g\) controla las normas y los
productos cruzados.

El flujo satisface \(g_{-s}=g_s^{-1}\).
La reciprocidad del radio de
\eqref{nt:mirror} pertenece a la carta espectral,
y la inversión de \(g_s\) pertenece a la
realización dimensional del registro. El enlace
constructivo utilizado aquí conserva las
aplicaciones explícitas de memoria a \(K\) y de
\(K\) a su dirección geométrica, seguido por
el transporte métrico de cada realización.
Esta composición mantiene, en un mismo desarrollo,
la reconstrucción exacta, el refinamiento positivo,
el registro de fronteras y la conservación
de los productos interiores.

\section{Incidencia ternaria, pegado de discriminantes y construcción de Leech}
\label{sec:lf-foundations}

La construcción reticular utiliza la información de incidencia que
acompaña al registro dodecafásico. Su dominio conserva las palabras
ternarias, el marco en que se representan, la orientación previa al
acarreo y la bandera que distingue una de las doce posiciones. El
paso a un retículo transmite esos datos mediante operaciones
sucesivas: codificación lineal, selección incidencial, pegado de
discriminantes y formación de un vecino. La clasificación reticular
interviene después de probar la positividad, paridad, autodualidad
y ausencia de raíces de la construcción obtenida.

Estas operaciones continúan la misma generación
APP--TRIT--TPK y el mismo refinamiento compatible del continuo
conjunto que anteceden a la lectura numérica \cite{HMT-IV}.
La matriz finita, el código y la métrica que se definen a
continuación son realizaciones con dominios precisos. La prueba
reticular recibe el estado marcado; el escalar final de una
constante no desempeña el papel de código de pegado ni de longitud.

\subsection{Carta elíptica y transporte de la incidencia}
\label{subsec:lf-incidence}

Las seis unidades \(\{1,2,4,5,7,8\}\) módulo nueve, conservadas
por el transporte modular, admiten una carta de seis posiciones.
La identificamos con
\(\mathbb P^1(\mathbb F_5)=\{\infty,0,1,2,3,4\}\), en ese orden.
Esta identificación fija coordenadas para realizar la relación
cuadrática del régimen elíptico del TRIT.

Sea \(\chi:\mathbb F_5\to\{0,1,-1\}\) el carácter cuadrático:
\(\chi(0)=0\), \(\chi(1)=\chi(4)=1\) y
\(\chi(2)=\chi(3)=-1\). Definimos la matriz entera
\(C_{\mathrm P}\) con diagonal cero,
\((C_{\mathrm P})_{\infty,a}=(C_{\mathrm P})_{a,\infty}=1\),
y entradas finitas \((C_{\mathrm P})_{a,b}=\chi(b-a)\).
Su cuadrado es \(5I_6\). En efecto, cada entrada diagonal del
cuadrado suma cinco unidades. Las entradas que involucran
\(\infty\) fuera de la diagonal se anulan porque la suma de
\(\chi\) es cero. Para dos índices finitos distintos, la
contribución uno de la posición \(\infty\) se cancela con
\(\sum_c\chi(c-a)\chi(c-b)=-1\). Esta última identidad se verifica
directamente sobre los cinco residuos, o mediante la sustitución
afín que lleva \((a,b)\) a \((0,1)\).

La reducción módulo tres produce la matriz simétrica
\begin{equation}
 A_W=
 \begin{pmatrix}
 0&1&1&1&1&1\\
 1&0&1&2&2&1\\
 1&1&0&1&2&2\\
 1&2&1&0&1&2\\
 1&2&2&1&0&1\\
 1&1&2&2&1&0
 \end{pmatrix},
 \qquad A_W^2=2I_6=-I_6
 \quad\text{en }\mathbb F_3.
 \label{eq:lf-aw}
\end{equation}
La igualdad conserva el tipo elíptico en este codominio finito.
El cuerpo finito y la posterior realización arquimediana de la
unidad imaginaria tienen dominios diferentes.

Para una matriz \(A\in M_6(\mathbb F_3)\) y una palabra fila
\(w\in\mathbb F_3^6\), escribimos
\[
 \gamma_A(w)=(w,wA)\in\mathbb F_3^{12}.
\]
Si \(f\in\operatorname{GL}_6(\mathbb F_3)\), entonces
\begin{equation}
 \gamma_{fAf^{-1}}(wf^{-1})
   =\gamma_A(w)(f^{-1}\oplus f^{-1}).
 \label{eq:lf-frame-covariance}
\end{equation}
La primera mitad es \(wf^{-1}\); la segunda es
\(wf^{-1}fAf^{-1}=wAf^{-1}\), lo que prueba la identidad.

Sea \(X_n^{\mathrm{mar}}\) el espacio de prefijos enriquecidos de
longitud \(n\), con marco inicial declarado, y sean
\(p_{n,m}:X_n^{\mathrm{mar}}\to X_m^{\mathrm{mar}}\) sus
truncaciones. En el paso \(j\), el prefijo contiene una palabra
visible \(w_j\) y un marco \(f_j\). La actualización del marco
tiene la forma \(f_{j+1}=g_j(x)f_j\), donde \(g_j(x)\) depende
únicamente de información ya contenida en el prefijo. Una
prolongación compatible conserva, por ello, todos los marcos
anteriores.

\begin{proposition}[Naturalidad del transporte incidencial]
\label{prop:lf-incidence-naturality}
La aplicación
\begin{equation}
 \mathcal Q_n(x)=
 \bigl((f_j,\gamma_{f_jA_Wf_j^{-1}}(w_j))\bigr)_{0\leq j<n}
 \label{eq:lf-Qn}
\end{equation}
satisface \(r_{n,m}\mathcal Q_n=\mathcal Q_mp_{n,m}\), donde \(r_{n,m}\) trunca
las primeras \(m\) componentes codificadas. Por tanto induce una
aplicación entre los límites inversos de los sistemas de
prefijos y de incidencias marcadas.
\end{proposition}
\begin{proof}
Dos prefijos compatibles tienen las mismas palabras y los mismos
operadores de transporte hasta el paso \(m-1\). Partiendo del
mismo marco inicial, la relación recursiva da por inducción los
mismos \(f_j\) hasta ese paso. Coinciden entonces las primeras
\(m\) componentes de \eqref{eq:lf-Qn}. La propiedad universal
del límite inverso produce la aplicación de prolongación.
\end{proof}

En cada marco, la primera proyección de \(\gamma_A\) recupera
exactamente la palabra visible. Este alcance es local: la
memoria profunda se conserva en el estado enriquecido y su
recuperación exige las coordenadas correspondientes. Al pasar
de palabras a soportes se retiene todavía menos información.
Asimismo, un cambio general de base en
\(\operatorname{GL}_6(\mathbb F_3)\) transporta también la métrica
de Hamming y la incidencia de la carta. Los cambios monomiales
las preservan directamente. Los soportes utilizados en esta
sección se calculan en la carta explícita \eqref{eq:lf-aw}.

La naturalidad anterior concreta la lectura incidencial que
acompaña a los cilindros numéricos. El encajamiento de cilindros
y la restricción de \eqref{eq:lf-Qn} actúan sobre un mismo
prefijo compatible, conservando respectivamente su evaluación
posicional y sus relaciones marcadas. El ambiente
\(\mathbb F_3^6\), de \(729\) palabras, proporciona el dominio
lineal de la codificación; conserva su distinción respecto de
las \(243\) palabras visibles y las \(468\) emisiones del censo
admisible de APP--TRIT--TPK.

\subsection{Código ternario autodual y enumeración de sus pesos}
\label{subsec:lf-code}

\begin{definition}[Código de la carta dodecafásica]
El código lineal asociado a \eqref{eq:lf-aw} es
\[
 C_W=\{(w,wA_W):w\in\mathbb F_3^6\}\subset\mathbb F_3^{12}.
\]
El peso \(\operatorname{wt}(c)\) cuenta las coordenadas no nulas
de \(c\), y \(\operatorname{supp}(c)\subset\{1,\ldots,12\}\)
es su conjunto de posiciones no nulas. Una hexada es el soporte
de una palabra de peso seis; su familia se denota \(\mathcal H\).
\end{definition}

\begin{theorem}[Autodualidad, distancia mínima y enumerador]
\label{thm:lf-code}
El código \(C_W\) es autodual de parámetros \([12,6,6]_3\) y
tiene enumerador
\begin{equation}
 W_{C_W}(z)=1+264z^6+440z^9+24z^{12}.
 \label{eq:lf-enumerator}
\end{equation}
\end{theorem}
\begin{proof}
La matriz generadora \([I_6\mid A_W]\) tiene rango seis.
Por la simetría de \(A_W\) y \eqref{eq:lf-aw},
\[
 [I_6\mid A_W][I_6\mid A_W]^{\mathsf T}
       =I_6+A_W^2=0.
\]
Así, el código es autoortogonal de dimensión mitad de la
longitud y, en consecuencia, autodual.

La enumeración restante es una operación entera finita sobre
la matriz explícita. Para \(r\in\{0,\ldots,6\}\), agrupamos
las palabras por \(\operatorname{wt}(w)=r\). La tabla siguiente
da el número de palabras con cada peso total:
\[
\begin{array}{c|rrrr|r}
r&0&6&9&12&\text{total}\\ \hline
0&1&0&0&0&1\\
1&0&12&0&0&12\\
2&0&60&0&0&60\\
3&0&120&40&0&160\\
4&0&60&180&0&240\\
5&0&12&180&0&192\\
6&0&0&40&24&64\\ \hline
\text{total}&1&264&440&24&729
\end{array}
\]
Cada fila contiene \(\binom6r2^r\) palabras. Para verificar
todas sus entradas, sea \(q\) un entero de \(0\) a \(728\), y
definamos sus seis cifras ternarias
\[
 d_i(q)=\left\lfloor\frac{q}{3^{6-i}}\right\rfloor\bmod3,
 \qquad
 e_j(q)=\left(\sum_{i=1}^{6}d_i(q)(A_W)_{ij}\right)\bmod3.
\]
La fila y la columna de la contribución de \(q\) son,
respectivamente,
\[
 r(q)=\sum_i\mathbf1_{\{d_i(q)\ne0\}},\qquad
 h(q)=r(q)+\sum_j\mathbf1_{\{e_j(q)\ne0\}}.
\]
La recurrencia de acumulación empieza con
\(T^{(0)}_{r,h}=0\) y aplica
\[
 T^{(q+1)}_{r,h}
 =T^{(q)}_{r,h}
   +\mathbf1_{\{r=r(q)\}}\mathbf1_{\{h=h(q)\}}
 \quad(0\leq q<729).
\]
Produce exactamente la tabla mostrada. La expansión ternaria
da una biyección entre esos \(729\) enteros y todas las palabras
de \(\mathbb F_3^6\), por lo que la operación agota el código.
La suma de sus filas prueba \eqref{eq:lf-enumerator}; el primer
peso no nulo es seis.
\end{proof}

\begin{remark}[Evaluación reproducible de la recurrencia finita]
El siguiente bucle implementa literalmente las sumas de la prueba.
Sus únicos datos son los enteros de la matriz
\eqref{eq:lf-aw}; calcula el enumerador antes de compararlo con
la igualdad obtenida.
\end{remark}
{\small
\begin{verbatim}
from collections import Counter
from itertools import product
A = ((0,1,1,1,1,1), (1,0,1,2,2,1),
     (1,1,0,1,2,2), (1,2,1,0,1,2),
     (1,2,2,1,0,1), (1,1,2,2,1,0))
rows = [Counter() for _ in range(7)]
for w in product(range(3), repeat=6):
    v = tuple(sum(w[i]*A[i][j] for i in range(6)) % 3
              for j in range(6))
    r = sum(x != 0 for x in w)
    h = r + sum(x != 0 for x in v)
    rows[r][h] += 1
total = sum(rows, Counter())
if dict(total) != {0:1, 6:264, 9:440, 12:24}:
    raise ArithmeticError("Enumerador incompatible")
\end{verbatim}
}

\begin{proposition}[Hexadas y diseño de Witt]
\label{prop:lf-witt}
La familia \(\mathcal H\) contiene \(132\) hexadas y cada
quintupla de posiciones pertenece a una única hexada.
\end{proposition}
\begin{proof}
Dos palabras de peso seis con el mismo soporte difieren por
un escalar. En efecto, entre sus seis cocientes de coordenadas
no nulas, uno de los signos \(1,-1\) aparece al menos tres
veces. Restando el múltiplo correspondiente, si las palabras
no fueran proporcionales, se obtendría una palabra no nula
de peso a lo sumo tres, en contradicción con
el teorema~\ref{thm:lf-code}. Cada soporte corresponde así a
la pareja \(\{c,-c\}\), y hay \(264/2=132\) soportes.

Si dos soportes distintos compartieran cinco posiciones,
en esa intersección uno de los dos cocientes aparecería al
menos tres veces. La combinación lineal que cancela dichas
coordenadas tendría soporte en una unión de tamaño siete
y peso a lo sumo cuatro. La distancia mínima lo impide.
Una quintupla pertenece, por tanto, a lo sumo a una hexada.
Como
\[
 132\binom65=792=\binom{12}{5},
\]
pertenece a exactamente una.
\end{proof}

El objeto construido es el código ternario extendido de Golay
con su pequeño diseño de Witt. El reconocimiento se efectúa
después de la codificación, la autodualidad y el cómputo de
distancia. Para la construcción reticular sólo se necesitan
las propiedades ya demostradas y la bandera marcada que se
obtiene a continuación.

\subsection{Bandera dodecafásica y selección del origen}
\label{subsec:lf-flag}

En el marco del registro generado se conservan las palabras
\[
 w_\pi=010211,\qquad
 w_e=201101,\qquad
 w_\varphi=121200,\qquad
 w_{\mathrm{APP}}=022020.
\]
Los subíndices \(\pi,e,\varphi\) identifican las correspondientes
regiones de lectura de las constantes; \(w_e\) conserva aquí
la región de Euler. Aplicar \(\gamma_{A_W}\) y tomar soportes da
\[
\begin{aligned}
 P_\pi&=\{2,4,5,6,7,8,9,10,11\},\\
 H_e&=\{1,3,4,6,9,10\},\\
 H_\varphi&=\{1,2,3,4,7,9\},\\
 H_{\mathrm{APP}}&=\{2,3,5,10,11,12\}.
\end{aligned}
\]
El primer soporte tiene peso nueve; los tres restantes son
hexadas. El registro
\[
 K=(234,543,140,729,659,824,621,58,914,794,146,601)
\]
selecciona
\[
 B_K=\{j:K_j\geq729\}=\{4,6,9,10\}=H_e\cap P_\pi.
\]
El umbral \(729\) pertenece al formato de bloques de la
lectura declarada. En este registro concreto, los umbrales
enteros de \(660\) a \(729\) producen el mismo soporte, de
modo que el soporte conserva la incidencia seleccionada,
mientras el registro completo conserva también los bloques
y el umbral de su lectura.

El registro firmado anterior al acarreo es
\[
\begin{split}
 Z_0={}&(7,297,353,-430,-715,-1199,\\
       &\hspace{9mm}-3,286,-894,-528,380,663).
\end{split}
\]
Su soporte negativo es la hexada
\[
 H_\alpha^-=\{4,5,6,7,9,10\}.
\]
Los signos se reciben del estado previo al acarreo y no se
extraen de la expansión del escalar final.

\begin{proposition}[Coincidencia de selección firmada e incidencial]
\label{prop:lf-flag}
La hexada \(H_\alpha^-\) es la única \(H\in\mathcal H\) que cumple
\[
\begin{gathered}
 B_K\subset H\subset P_\pi,\\
 (|H\cap H_e|,|H\cap H_\varphi|,
                   |H\cap H_{\mathrm{APP}}|)=(4,3,2).
\end{gathered}
\]
\end{proposition}
\begin{proof}
La codificación de las palabras de peso seis, o la recurrencia
finita anterior aplicada a sus soportes, da las cuatro
hexadas sobre \(B_K\). Sus pares complementarios son
\[
 \{1,3\},\qquad\{2,11\},\qquad\{5,7\},\qquad\{8,12\}.
\]
Sólo el segundo y el tercero están contenidos en \(P_\pi\).
La hexada del segundo tiene intersecciones de tamaño tres
con \(H_\varphi\) y con \(H_{\mathrm{APP}}\), y por tanto
incumple la última condición. La del tercero tiene
intersecciones \(4,3,2\), y coincide con \(H_\alpha^-\).
\end{proof}

\begin{lemma}[Origen de la bandera marcada]
\label{lem:lf-origin}
El origen seleccionado por la incidencia es
\[
 o_*=(H_e\cap H_\varphi)
                 \setminus(H_\alpha^-\cup H_{\mathrm{APP}})
     =\{1\}.
\]
La selección conmuta con toda biyección aplicada
simultáneamente a los soportes.
\end{lemma}
\begin{proof}
La intersección \(H_e\cap H_\varphi\) es
\(\{1,3,4,9\}\). La unión sustraída contiene \(3,4,9\)
y excluye \(1\). La equivariancia procede de que las
biyecciones conmutan con unión, intersección y diferencia.
\end{proof}

La bandera distingue una componente entre doce antes del
paso métrico. El vector radial se determina después,
dentro de una cámara positiva declarada. Así quedan
separadas y articuladas la selección combinatoria y la
operación reticular.

\subsection{Pegado ternario de las doce componentes}
\label{subsec:lf-glue}

Sea \(A_2=\mathbb Za_1\oplus\mathbb Za_2\), con forma
bilineal de matriz
\[
 B_2=\begin{pmatrix}2&-1\\-1&2\end{pmatrix},
 \qquad \rho=a_1+a_2,\qquad \rho^2=2.
\]
Para \(x=ua_1+va_2\),
\[
 x^2=2(u^2-uv+v^2).
\]
La identidad
\(4(u^2-uv+v^2)=(2u-v)^2+3v^2\) muestra
la positividad. El determinante es tres. Las tres
clases de \(A_2^*/A_2\) se representan mediante
\[
 \varpi_0=0,\qquad
 \varpi_1=\frac{2a_1+a_2}{3},\qquad
 \varpi_2=\frac{a_1+2a_2}{3}.
\]
El emparejamiento con \(a_1,a_2\) verifica que estos
vectores pertenecen a \(A_2^*\), y
\(2\varpi_1-\varpi_2=a_1\). Por ello
\(j\mapsto\varpi_j+A_2\) identifica aditivamente
\(\mathbb F_3\) con el discriminante.

Para \(c\in C_W\), definimos
\(\phi(c)=(\varpi_{c_1},\ldots,\varpi_{c_{12}})\)
y
\begin{equation}
 N=\bigcup_{c\in C_W}\bigl(A_2^{12}+\phi(c)\bigr).
 \label{eq:lf-glue}
\end{equation}

\begin{theorem}[Retículo de pegado par y unimodular]
\label{thm:lf-niemeier}
El conjunto \(N\) es un retículo positivo, par y
unimodular de rango \(24\). Sus vectores de norma dos
son exactamente las raíces de \(A_2^{12}\).
\end{theorem}
\begin{proof}
La linealidad del código y la aditividad del discriminante
hacen que la unión sea un subgrupo de
\((A_2^*)^{12}\) que contiene \(A_2^{12}\).
Es, por tanto, un retículo de rango \(24\).

El emparejamiento de las clases discriminantes es
\(2cd/3\) módulo \(\mathbb Z\) en cada componente.
La autoortogonalidad de \(C_W\) anula su suma y prueba
la integralidad de \(N\). La norma de una clase
representada por \(\phi(c)\) es
\(2\operatorname{wt}(c)/3\) módulo \(2\mathbb Z\).
Todos los pesos de \eqref{eq:lf-enumerator} son
múltiplos de tres, luego \(N\) es par.

Su índice sobre \(A_2^{12}\) es \(|C_W|=3^6\).
Con el determinante de la matriz de Gram,
\[
 \det N=\frac{\det(A_2^{12})}{[N:A_2^{12}]^2}
       =\frac{3^{12}}{(3^6)^2}=1.
\]
La forma ambiente es la suma ortogonal de doce formas
positivas; queda demostrada la positividad.
La integralidad y el determinante uno implican
\(N^*=N\).

En una clase no nula del discriminante local, un
vector tiene la forma \((ua_1+va_2)/3\), con residuos
\((u,v)\equiv(2,1)\) o \((1,2)\pmod3\).
El entero \(u^2-uv+v^2\) es positivo y divisible
por tres, luego es al menos tres. Los representantes
\(\varpi_1,\varpi_2\) alcanzan ese valor: la norma
mínima de cada clase no nula es \(2/3\).
Una palabra no nula del código tiene por lo menos
seis coordenadas no nulas. Todo vector de una clase
de pegado no trivial tiene, por tanto, norma
al menos \(6(2/3)=4\).

En \(A_2^{12}\), una norma dos sólo puede proceder
de una única componente, porque cada componente
tiene norma par positiva mínima dos. Las soluciones
de \(u^2-uv+v^2=1\) son
\((\pm1,0),(0,\pm1),(1,1),(-1,-1)\).
Son las seis raíces de \(A_2\), y completan el censo
de raíces de \(N\).
\end{proof}

Las propiedades anteriores identifican \(N\) como el
retículo de Niemeier de sistema de raíces \(A_2^{12}\).
La descripción de todas sus raíces es necesaria para
controlar cuáles sobreviven al paso al vecino.

\subsection{Cámara radial y vecino sin raíces}
\label{subsec:lf-neighbour}

Se fija la cámara radial
\[
 v(a)=\sum_{j=1}^{12}a_j\rho_j,\qquad
 a_j=1+3b_j,\qquad b_j\in\mathbb Z_{\geq0},
\]
donde \(\rho_j\) es la copia de \(\rho\) en la
componente \(j\). La paridad de un vecino de índice
tres requiere \(v(a)^2\equiv0\pmod{18}\).
Como
\[
 v(a)^2=2\sum_j(1+3b_j)^2
       =24+12\sum_jb_j+18\sum_jb_j^2,
\]
la congruencia equivale a \(\sum_jb_j\equiv1\pmod3\).
La norma mínima dentro de la cámara se alcanza con
un único \(b_j=1\) y los demás cero, y vale \(54\).
El origen \(o_*=1\) del lema~\ref{lem:lf-origin}
selecciona
\begin{equation}
 v_\alpha=4\rho_1+\rho_2+\cdots+\rho_{12},
 \qquad v_\alpha^2=54.
 \label{eq:lf-valpha}
\end{equation}
El calificativo de mínimo se refiere a la cámara
declarada. En la clase residual completa,
\((a_1-2a_2,\rho,\ldots,\rho)\) representa la
misma clase módulo tres y tiene norma
\(14+11\cdot2=36\). Esta distinción preserva
la condición que determina el coeficiente cuatro.

Definimos
\begin{equation}
\begin{split}
 N_0&=\{x\in N:\langle x,v_\alpha\rangle\equiv0\pmod3\},\\
 \Lambda&=N_0+\mathbb Z\,v_\alpha/3.
\end{split}
\label{eq:lf-lambda}
\end{equation}

\begin{theorem}[Vecino par autodual sin raíces]
\label{thm:lf-leech}
La construcción \eqref{eq:lf-lambda} es un
retículo positivo, par y autodual de rango \(24\)
sin vectores de norma dos. En consecuencia, es
isométrica al retículo de Leech.
\end{theorem}
\begin{proof}
El carácter \(x\mapsto\langle x,v_\alpha\rangle\bmod3\)
es no nulo. Una raíz simple de la componente \(j\)
tiene emparejamiento \(a_j\equiv1\pmod3\) con
\(a_j\rho_j\). Por tanto \([N:N_0]=3\), de donde
\(\det N_0=9\).

Para \(x\in N_0\), el emparejamiento
\(\langle x,v_\alpha/3\rangle\) es entero, y así
\(v_\alpha/3\in N_0^*\). Además,
\(\langle v_\alpha,v_\alpha\rangle=54\) implica
\(v_\alpha\in N_0\). Si \(v_\alpha/3\in N\),
la integralidad de \(N\) obligaría a que todos los
emparejamientos de \(N\) con \(v_\alpha\) fuesen
divisibles por tres, contradiciendo el carácter
no nulo. La clase de \(v_\alpha/3\) sobre \(N_0\)
tiene entonces orden exactamente tres. Resulta
\[
 [\Lambda:N_0]=3,\qquad
 \det\Lambda=\frac9{3^2}=1.
\]
Los emparejamientos de \(N_0\) con el nuevo
generador son enteros y la norma de este generador
es seis. La extensión es integral. Para
\(x\in N_0\) y \(m\in\mathbb Z\),
\[
 \left(x+m\frac{v_\alpha}{3}\right)^2
 =x^2+2m\frac{\langle x,v_\alpha\rangle}{3}
       +6m^2\in2\mathbb Z.
\]
Por ello \(\Lambda\) es par; su integralidad y
determinante uno prueban la autodualidad.
La positividad y el rango se heredan del espacio
euclídeo ambiente.

Para excluir raíces, primero consideramos la
clase \(N_0\). Todas las raíces de \(N\) pertenecen
a \(A_2^{12}\) por el teorema~\ref{thm:lf-niemeier}.
Los emparejamientos de sus seis raíces locales
con \(a_j\rho_j\) son \(\pm a_j,\pm2a_j\),
congruentes con \(\pm1,\pm2\pmod3\).
Ninguna raíz pertenece a \(N_0\).

Las otras dos clases de \(\Lambda/N_0\) pueden
representarse con \(\pm v_\alpha/3\).
En cada componente, sus vectores pertenecen a
una de las clases
\[
 A_2+\varpi_j\pm\rho/3,\qquad j=0,1,2.
\]
La componente distinguida tiene la misma
descripción porque \(4\rho/3\equiv\rho/3\pmod{A_2}\).
Para los numeradores \((u,v)\) en
\((ua_1+va_2)/3\), los seis residuos son
\[
 (1,1),(0,2),(2,0),(2,2),(1,0),(0,1)\pmod3.
\]
Todos son no nulos. La positividad de
\(u^2-uv+v^2\) obliga a que sea al menos uno,
y los representantes
\((1,1),(0,-1),(-1,0),(-1,-1),(1,0),(0,1)\)
alcanzan ese mínimo en las seis clases.
Cada componente tiene norma al menos \(2/9\).
Un vector de cualquiera de las dos clases
nuevas tiene norma al menos
\[
 12\cdot\frac29=\frac83>2.
\]
Se han excluido raíces en las tres clases.
El teorema de clasificación de los retículos
positivos pares unimodulares de rango veinticuatro
identifica el único caso sin raíces con el
retículo de Leech \cite{ConwaySloane}.
\end{proof}

La identificación final utiliza una clasificación
externa sobre un objeto cuyas propiedades ya han
sido demostradas. La selección de la componente
distinguida conserva la bandera procedente del
registro firmado; el subíndice \(\alpha\) expresa
esa procedencia. La norma y los coeficientes del
vector radial se obtienen de la métrica y de
la cámara mediante la congruencia indicada.

El retículo \(\Lambda\), su forma bilineal, sus
doce planos ambientes y el vector \(v_\alpha/3\)
quedan así determinados con todas las propiedades
utilizadas por la deformación posterior. El
código contiene asimismo la palabra
\[
 c_*=(1,1,1,1,1,1,2,1,1,1,1,1),
\]
pues \(q_*=(1,1,1,1,1,1)\) satisface
\(q_*A_W=(2,1,1,1,1,1)\). Éste es el dato
de soporte completo que permite construir y
probar la estabilidad de la isometría de orden
tres en la sección~\ref{sec:leech-dualidad}.

\section{Polarización dodecafásica, dualidad reticular y reciprocidad theta}
\label{sec:leech-dualidad}

La incidencia excepcional y la polarización del registro dodecafásico
proceden de dos realizaciones del mismo estado enriquecido. La primera
conserva las relaciones de frontera, el código de pegado y el origen
marcado que determinan el retículo; la segunda hace actuar una
representación sobre las doce coordenadas enteras ya obtenidas por los
lectores del Topological Phase Kernel. Su composición permite estudiar
una deformación métrica cuyo parámetro modifica las longitudes relativas
de doce planos y conserva simultáneamente el volumen total y una
simetría ternaria explícita.

El punto de partida causal sigue siendo la construcción
APP--TRIT--TPK del estado enriquecido y del continuo conjunto.
El registro \(K\) se recibe después del registro dodecafásico y de sus
operaciones de incidencia y orientación. La representación empleada
actúa sobre esa salida y conserva sus doce valores como datos ya
determinados. La realización reticular se recibe con su forma bilineal,
sus clases de pegado y su vector marcado. Estas estructuras, desarrolladas
respectivamente en \cite{hmtX,HMT-IV}, se mantienen al componerlas.
La familia que se construye a continuación es una realización geométrica
de esos datos con una carta común de doce posiciones.

\subsection{Polarización del registro y marcación de los doce planos}
\label{subsec:leech-polarizacion}

El registro dodecafásico y su suma son
\begin{equation}
 \begin{split}
 K={}&(234,543,140,729,659,824,\\
     &\hspace{13mm}621,58,914,794,146,601),\\
 \sum_{j=1}^{12}K_j={}&6263.
 \end{split}
 \label{eq:leech-K}
\end{equation}
Sea \(A_5\) el grupo de permutaciones pares de cinco elementos y sea
\(C_5=\langle(1\,2\,3\,4\,5)\rangle\). Sus doce clases izquierdas
determinan una representación por permutaciones
\(\varrho:A_5\to O(\mathbb R^{12})\). La carta marcada se fija mediante
los representantes siguientes, escritos como listas de imágenes:
\[
\begin{array}{c|c@{\qquad}c|c}
j&r_j&j&r_j\\ \hline
1&(4,3,5,2,1)&7&(5,4,3,2,1)\\
2&(4,2,1,3,5)&8&(1,3,4,2,5)\\
3&(1,2,4,5,3)&9&(4,2,3,5,1)\\
4&(2,5,3,4,1)&10&(3,2,5,4,1)\\
5&(4,3,1,5,2)&11&(4,5,2,3,1)\\
6&(5,1,2,3,4)&12&(5,3,2,4,1).
\end{array}
\]
Concretamente, \(\varrho(a)e_j=e_\ell\) si
\(ar_jC_5=r_\ell C_5\). El carácter tridimensional \(\chi_3\) toma
los valores
\[
\begin{array}{c|ccccc}
\text{clase}&1&2^2&3&5_a&5_b\\ \hline
\chi_3&3&-1&0&(1+\sqrt5)/2&(1-\sqrt5)/2.
\end{array}
\]
La clase \(5_a\) contiene el generador de \(C_5\) indicado y \(5_b\)
contiene su cuadrado. Esta elección distingue los dos caracteres
tridimensionales conjugados. Si \(I_{12}\) es la identidad y
\(\mathbf1\) es el vector de doce unos, se definen
\begin{equation}
 P_3=\frac3{60}\sum_{a\in A_5}\chi_3(a^{-1})\varrho(a),
 \qquad
 P_{11}=I_{12}-\frac1{12}\mathbf1\mathbf1^{\mathsf T}.
 \label{eq:leech-projectors}
\end{equation}
El campo \(\mathbb Q(\sqrt5)\) aparece aquí en la realización de la
representación que actúa sobre el registro generado.

\begin{lemma}[Polarización de suma nula]
\label{lem:leech-u}
Los operadores de \eqref{eq:leech-projectors} satisfacen
\[
 P_3=P_3^{\mathsf T}=P_3^2,\qquad
 \operatorname{rank}P_3=3,\qquad P_3\mathbf1=0,\qquad
 P_3P_{11}=P_{11}P_3=P_3.
\]
Por consiguiente,
\begin{equation}
 u=P_3P_{11}K=P_3K,\qquad
 \sum_{j=1}^{12}u_j=0,\qquad
 \|u\|^2=\frac{6638585+2275584\sqrt5}{20}>0.
 \label{eq:leech-u-norm}
\end{equation}
\end{lemma}
\begin{proof}
La suma del carácter es el proyector isótipo de la representación.
La realidad del carácter y la sustitución \(a\mapsto a^{-1}\)
dan su simetría. La multiplicidad del carácter tridimensional en la
representación de \(A_5/C_5\) es la dimensión de sus invariantes por
\(C_5\), que vale
\[
 \frac15\left(
 3+2\frac{1+\sqrt5}{2}+2\frac{1-\sqrt5}{2}
 \right)=1.
\]
El proyector tiene por ello rango tres. Su ortogonalidad al carácter
trivial implica \(P_3\mathbf1=0\) y las identidades con \(P_{11}\).
La aplicación de la suma finita al registro
\eqref{eq:leech-K} produce las doce coordenadas
\begin{equation}
\begin{aligned}
u={}&(-495/4-233\sqrt5/10,\quad403/4+169\sqrt5/10,\\
&-403/4-169\sqrt5/10,\quad495/4+233\sqrt5/10,\\
&601/4+1031\sqrt5/10,\quad203/4+211\sqrt5/5,\\
&-203/4-211\sqrt5/5,\quad-601/4-1031\sqrt5/10,\\
&313/4+569\sqrt5/10,\quad162+219\sqrt5/20,\\
&-162-219\sqrt5/20,\quad-313/4-569\sqrt5/10).
\end{aligned}
\label{eq:leech-u-coordinates}
\end{equation}
Sumarlas y elevarlas al cuadrado en \(\mathbb Q(\sqrt5)\) da
\eqref{eq:leech-u-norm}. Equivalentemente,
\[
 \|u\|^2=K^{\mathsf T}P_3K
 =\frac1{20}\sum_{a\in A_5}
          \chi_3(a^{-1})K^{\mathsf T}\varrho(a)K.
\]
Los coeficientes positivos de su evaluación prueban que \(u\ne0\).
\end{proof}

La suma nula tiene una función geométrica concreta: las expansiones de
unos planos podrán compensar las contracciones de los restantes en el
determinante total. La no nulidad de \(u\), en cambio, garantiza que esa
compensación contiene una deformación efectiva.

El espacio euclídeo ambiente de la realización reticular es
\begin{equation}
 E=\bigoplus_{j=1}^{12}E_j,\qquad
 B=\operatorname{diag}(B_2,\ldots,B_2),\qquad
 B_2=\begin{pmatrix}2&-1\\-1&2\end{pmatrix}.
 \label{eq:leech-ambient}
\end{equation}
Cada \(E_j\) es el plano real generado por el retículo de raíces
\(A_2\), expresado en su base de raíces simples. La descomposición
ortogonal pertenece al espacio ambiente. El retículo de Leech se
obtiene mediante pegado y paso a un vecino, operaciones que relacionan
los doce factores.

Sea \(C_W\subset\mathbb F_3^{12}\) el código ternario de dimensión seis
de la construcción Paley--Witt. Identificando el discriminante de cada
factor \(A_2\) con \(\mathbb F_3\), el pegado es
\[
 N=N(C_W)=
 \bigcup_{c\in C_W}\bigl(A_2^{12}+\phi(c)\bigr),
\]
donde \(\phi(c)\) elige representantes de las clases discriminantes.
La autoortogonalidad y autodualidad del código transmiten paridad e
integralidad a \(N\); su determinante es
\[
 \det N=\frac{3^{12}}{(3^6)^2}=1.
\]
La construcción excepcional fija el origen marcado \(j=1\) y el
vector
\begin{equation}
 v_\alpha=4\rho_1+\rho_2+\cdots+\rho_{12},
 \qquad \rho_j=(1,1)_j,\qquad
 \|v_\alpha\|_B^2=54.
 \label{eq:leech-marked-vector}
\end{equation}
Con él se forman
\begin{equation}
\begin{split}
 N_{\alpha,0}
  &=\{x\in N:\langle x,v_\alpha\rangle_B\equiv0\pmod3\},\\
 \Lambda&=N_{\alpha,0}+\mathbb Z\,v_\alpha/3.
\end{split}
\label{eq:leech-neighbour}
\end{equation}
La realización excepcional establece que \(\Lambda\) es un retículo
positivo, par, unimodular, sin raíces y de rango \(24\), de modo que es
isométrico al retículo de Leech \cite{HMT-IV}. En particular,
\(\Lambda^*=\Lambda\), con dual respecto de \(B\).
Estas propiedades y la pertenencia \(v_\alpha/3\in\Lambda\) son las
hipótesis reticulares precisas empleadas en los teoremas siguientes.

\subsection{Isometría ternaria y estabilidad del vecino}
\label{subsec:leech-c3}

Definimos las matrices bidimensionales
\[
 C=\begin{pmatrix}0&-1\\1&-1\end{pmatrix},
 \qquad
 R=\begin{pmatrix}0&1\\1&0\end{pmatrix},
 \qquad
 \varpi=\frac13\begin{pmatrix}2\\1\end{pmatrix}.
\]
La multiplicación directa verifica
\[
 C^3=I_2,\quad C^2+C+I_2=0,\quad
 C^{\mathsf T}B_2C=B_2,\quad
 RCR=C^{-1},\quad R^{\mathsf T}B_2R=B_2.
\]
Además, \(C\varpi-\varpi=(-1,0)^{\mathsf T}\), por lo que \(C\)
actúa trivialmente en \(A_2^*/A_2\). La reflexión \(R\) cambia el signo
de la clase discriminante. El código \(C_W\) contiene la palabra de
soporte completo
\[
 c_*=(1,1,1,1,1,1,2,1,1,1,1,1).
\]
En la presentación \(C_W=\{(q,qA_W):q\in\mathbb F_3^6\}\), dicha
palabra procede de \(q_*=(1,1,1,1,1,1)\) y
\(q_*A_W=(2,1,1,1,1,1)\).

Para cada plano se fija el marco
\begin{equation}
 Q_j=\begin{cases}
 R,&(c_*)_j=1,\\
 I_2,&(c_*)_j=2,
 \end{cases}
 \qquad
 g_c=\bigoplus_{j=1}^{12}Q_jCQ_j^{-1}.
 \label{eq:leech-c3}
\end{equation}
Los símbolos \(Q_j\) designan marcos de planos; el símbolo \(P_3\)
queda reservado para el proyector sobre el sector tridimensional.

\begin{proposition}[Simetría de orden tres del retículo marcado]
\label{prop:leech-c3}
La transformación \(g_c\) es una isometría de orden tres de \(E\),
tiene como único vector fijo el cero y satisface
\[
 g_cN=N,\qquad g_cN_{\alpha,0}=N_{\alpha,0},
 \qquad g_c\Lambda=\Lambda.
\]
\end{proposition}
\begin{proof}
Cada bloque es \(C\) o \(C^{-1}\), cuyo polinomio mínimo es
\(X^2+X+1\). Esto prueba el orden y la afirmación sobre vectores fijos,
así como la conservación de \(B\). Cada bloque actúa trivialmente en
el discriminante y conserva, por tanto, todas las clases de pegado de
\(N\).

Escribamos \(v=v_\alpha=\sum_j a_j\rho_j\), con \(a_1=4\) y
\(a_j=1\) para \(j>1\). Todos los \(a_j\) son congruentes con uno
módulo tres. Las identidades
\[
 (C-I_2)\rho=(-2,-1)^{\mathsf T},
 \qquad (C^{-1}-I_2)\rho=(-1,-2)^{\mathsf T}
\]
muestran que
\[
 y=\frac{g_cv-v}{3},\qquad z=\frac{g_c^{-1}v-v}{3}
\]
tienen clases discriminantes \(c_*\) y \(-c_*\), respectivamente.
Ambas pertenecen a \(C_W\), luego \(y,z\in N\). En cada factor,
\[
 \left\langle
 \frac{(C^{\pm1}-I_2)a_j\rho}{3},a_j\rho
 \right\rangle_{B_2}=-a_j^2.
\]
Por suma, \(\langle y,v\rangle_B=\langle z,v\rangle_B=-27\);
por ello \(y,z\in N_{\alpha,0}\).
Si \(x\in N_{\alpha,0}\), entonces
\[
 \langle g_cx,v\rangle_B
  =\langle x,g_c^{-1}v\rangle_B
  =\langle x,v\rangle_B+3\langle x,z\rangle_B
  \equiv0\pmod3.
\]
La misma operación con \(g_c^{-1}\) demuestra la igualdad
\(g_cN_{\alpha,0}=N_{\alpha,0}\). Finalmente,
\(g_c(v/3)=v/3+y\) conserva la clase que genera
\(\Lambda=N_{\alpha,0}+\mathbb Z(v/3)\).
\end{proof}

La simetría ternaria actúa, por consiguiente, sobre el retículo
pegado completo. Su conservación utiliza la palabra del código y
la congruencia del vector marcado, además de la rotación de cada
plano. Esta precisión permite transportar la simetría al flujo sin
perder la información aritmética que hizo posible su construcción.

\subsection{Flujo anisótropo y conservación del covolumen}
\label{subsec:leech-flow}

\begin{definition}[Deformación determinada por la polarización]
Para \(s\in\mathbb R\), parámetro adimensional, se define
\begin{equation}
 H_u=\bigoplus_{j=1}^{12}u_jI_2,\qquad
 D_s=\exp(sH_u)=\bigoplus_{j=1}^{12}e^{su_j}I_2,\qquad
 \Lambda_s=D_s\Lambda.
 \label{eq:leech-flow}
\end{equation}
La coordenada \(u_j\) se asigna al plano \(E_j\) de la misma carta
marcada.
\end{definition}

El factor positivo \(e^{su_j}\) conserva la orientación de cada
plano. Las doce etiquetas transmiten la correspondencia entre la
polarización y la realización reticular. Esta correspondencia
forma parte de la estructura de partida; una permutación de las
etiquetas se trata mediante su transporte simultáneo en ambos objetos.

\begin{theorem}[Covolumen y simetría del flujo]
\label{thm:leech-flow}
Para todos \(s,t\in\mathbb R\) se cumple
\[
\begin{gathered}
 D_{s+t}=D_sD_t,\qquad D_0=I,\qquad D_s^{-1}=D_{-s},\\
 D_s^\dagger=D_s,\qquad D_sg_c=g_cD_s,\qquad \det D_s=1,
\end{gathered}
\]
donde \(A^\dagger=B^{-1}A^{\mathsf T}B\).
Cada \(\Lambda_s\) es una red discreta de rango \(24\) y covolumen
uno, preservada por la isometría \(g_c\) de orden tres.
\end{theorem}
\begin{proof}
Las identidades de grupo se verifican en cada bloque exponencial.
Los bloques escalares son autoadjuntos para \(B_2\), y conmutan
con los bloques de \(g_c\). La suma nula de
\eqref{eq:leech-u-norm} da
\[
 \det D_s=\prod_{j=1}^{12}e^{2su_j}
 =\exp\left(2s\sum_{j=1}^{12}u_j\right)=1.
\]
La imagen de una red por un operador invertible es una red del
mismo rango; su covolumen se multiplica por el valor absoluto del
determinante. El covolumen de \(\Lambda\) es uno, de modo que el
de \(\Lambda_s\) también lo es. La conmutación y
la proposición~\ref{prop:leech-c3} dan
\[
 g_c\Lambda_s=g_cD_s\Lambda=D_sg_c\Lambda=D_s\Lambda.
\]
El orden y el espacio de puntos fijos de \(g_c\) permanecen
inalterados.
\end{proof}

\begin{remark}[Parámetro y simetría conservada]
La dirección \(u\) queda determinada por el registro; \(s\) recorre
una familia de realizaciones. Reemplazar \(u\) por \(u/\|u\|\)
reparametriza la misma familia. Una interpretación física que
seleccione un valor particular de \(s\) añade su regla de realización.
El teorema conserva la simetría explícita \(g_c\); el transporte de
otros automorfismos requiere sus correspondientes relaciones con
\(H_u\).
\end{remark}

\subsection{Dual euclídeo y deformación efectiva}
\label{subsec:leech-dual}

\begin{definition}[Dual euclídeo]
Para una red \(L\subset E\), su dual respecto de \(B\) es
\[
 L^*=\{y\in E:\langle y,x\rangle_B\in\mathbb Z
                         \text{ para todo }x\in L\}.
\]
\end{definition}

\begin{theorem}[Reciprocidad del dual]
\label{thm:leech-dual}
La familia \eqref{eq:leech-flow} satisface
\begin{equation}
 \Lambda_s^*=\Lambda_{-s}\qquad(s\in\mathbb R).
 \label{eq:leech-dual}
\end{equation}
Su autodualidad en un parámetro \(s\) equivale a
\(D_{2s}\Lambda=\Lambda\).
\end{theorem}
\begin{proof}
Para un operador invertible \(A\),
\[
 (AL)^*=A^{-\dagger}L^*.
\]
En efecto, \(\langle y,Ax\rangle_B=\langle A^\dagger y,x\rangle_B\)
es entero para todo \(x\in L\) precisamente cuando
\(A^\dagger y\in L^*\). Tomando \(A=D_s\), la autoadjunción,
la identidad \(D_s^{-1}=D_{-s}\) y \(\Lambda^*=\Lambda\) prueban
\eqref{eq:leech-dual}. La igualdad
\(D_s\Lambda=D_{-s}\Lambda\) equivale, multiplicando por \(D_s\),
a \(D_{2s}\Lambda=\Lambda\).
\end{proof}

El covolumen controla el volumen de un dominio fundamental; la
autodualidad controla todos los productos enteros con la red.
La primera propiedad es continua a lo largo de esta familia. La
segunda impone la condición aritmética exacta del teorema anterior.
Se empleará por ello la expresión \emph{covolumen uno} para el
flujo euclídeo y se especificará la integralidad cuando intervenga
la unimodularidad aritmética.

\begin{proposition}[Testigo de pérdida local de integralidad]
\label{prop:leech-nonintegral}
Existe \(\delta>0\) tal que, si \(0<|s|<\delta\), la red
\(\Lambda_s\) es no integral y no es isométrica al retículo de Leech.
\end{proposition}
\begin{proof}
El vector \(v_\alpha/3\) pertenece a \(\Lambda\). Su norma
transportada es
\begin{equation}
 f(s)=\left\|D_s\frac{v_\alpha}{3}\right\|_B^2
 =\frac29\left(16e^{2su_1}+
                       \sum_{j=2}^{12}e^{2su_j}\right).
 \label{eq:leech-witness}
\end{equation}
Por \(\sum_j u_j=0\) y por la primera coordenada de
\eqref{eq:leech-u-coordinates},
\begin{equation}
\begin{split}
 f(0)&=6,\\
 f'(0)&=\frac49\left(16u_1+\sum_{j=2}^{12}u_j\right)
       =\frac{20}{3}u_1
       =-825-\frac{466}{3}\sqrt5\ne0.
\end{split}
\label{eq:leech-witness-derivative}
\end{equation}
La continuidad de \(f'\) permite elegir un intervalo
\((-\delta,\delta)\) en el que \(f\) sea estrictamente monótona.
Reduciendo \(\delta\), se obtiene además
\(11/2<f(s)<13/2\). Para \(0<|s|<\delta\), resulta \(f(s)\ne6\);
en ese intervalo, seis es el único entero. Por tanto, \(\Lambda_s\)
contiene un vector de norma no entera. Una red integral tiene
norma entera en todos sus vectores. Puesto que Leech es integral
y la norma se conserva por isometrías, queda demostrada también
la segunda afirmación.
\end{proof}

La deformación preserva volumen y simetría ternaria al tiempo que
modifica efectivamente la estructura métrica aritmética. Los pares
opuestos de coordenadas de \(u\) definen una permutación del espacio
ambiente que intercambia dilataciones y contracciones. Su promoción
a un automorfismo del retículo pegado exigiría conservar el código
y el vecino. La reciprocidad del dual ya demostrada utiliza
únicamente la autoadjunción del flujo y la autodualidad inicial.

\subsection{Convergencia y reciprocidad de las funciones theta}
\label{subsec:leech-theta}

Para \(s\in\mathbb R\) y \(t>0\), definimos
\begin{equation}
 \Theta_s(t)=\sum_{\lambda\in\Lambda}
                \exp\bigl(-\pi t\|D_s\lambda\|_B^2\bigr).
 \label{eq:leech-theta-real}
\end{equation}
Es la suma gaussiana de todas las longitudes de la red deformada,
con sus multiplicidades. Si \(M=\max_j|u_j|\), se tiene
\begin{equation}
 e^{-2|s|M}\|x\|_B^2
 \leq\|D_sx\|_B^2
 \leq e^{2|s|M}\|x\|_B^2.
 \label{eq:leech-quadratic-bounds}
\end{equation}
Sobre un compacto de \(\mathbb R\times(0,\infty)\), la cota
inferior proporciona una gaussiana dominante sobre la red fija.
La serie converge absoluta y uniformemente. Las derivadas de
orden finito añaden factores polinomiales en las coordenadas de
\(\lambda\); dichos factores quedan dominados por otra gaussiana.
Las derivaciones término a término son, por consiguiente,
legítimas en esos compactos.

\begin{theorem}[Reciprocidad gaussiana]
\label{thm:leech-theta}
Con el volumen euclídeo asociado a \(B\),
\begin{equation}
 \Theta_s(t)=t^{-12}\Theta_{-s}(t^{-1}),
 \qquad s\in\mathbb R,\quad t>0.
 \label{eq:leech-theta-reciprocity}
\end{equation}
En particular, \(\Theta_s(1)=\Theta_{-s}(1)\).
\end{theorem}
\begin{proof}
Fijamos la transformada de Fourier
\[
 \widehat f(y)=\int_E
      f(x)e^{-2\pi i\langle x,y\rangle_B}\,dx.
\]
Para una red \(L\subset E\) y
\(f_t(x)=e^{-\pi t\|x\|_B^2}\), periodizamos
\[
 F_t(x)=\sum_{\ell\in L}f_t(x+\ell).
\]
La convergencia normal de la serie y sus derivadas da una función
suave sobre \(E/L\). Si \(\xi\in L^*\), su coeficiente de Fourier es
\[
\begin{split}
 c_\xi
 &=\frac1{\operatorname{covol}(L)}
   \int_{E/L}F_t(x)e^{-2\pi i\langle x,\xi\rangle_B}\,dx\\
 &=\frac{\widehat f_t(\xi)}{\operatorname{covol}(L)}.
\end{split}
\]
La segunda igualdad descompone \(E\) en traslaciones de un dominio
fundamental; el carácter tiene valor uno sobre \(L\). El producto
de las integrales gaussianas unidimensionales, en una base
ortonormal para \(B\), da
\[
 \widehat f_t(\xi)=
 t^{-12}\exp\bigl(-\pi\|\xi\|_B^2/t\bigr).
\]
La serie de Fourier es absolutamente convergente. Evaluándola en
\(x=0\), obtenemos la suma de Poisson
\[
 \sum_{\ell\in L}e^{-\pi t\|\ell\|_B^2}
 =\frac{t^{-12}}{\operatorname{covol}(L)}
      \sum_{\xi\in L^*}e^{-\pi\|\xi\|_B^2/t}.
\]
Ahora \(L=\Lambda_s\) tiene covolumen uno y
\(L^*=\Lambda_{-s}\) por el teorema~\ref{thm:leech-dual}.
La sustitución prueba \eqref{eq:leech-theta-reciprocity}.
El caso \(t=1\) es inmediato.
\end{proof}

La periodización anterior es la demostración gaussiana de Poisson;
puede cotejarse con \cite[teorema 2, pp.~10--11]{Elkies}.
La convención con signo positivo en Fourier conduce a la misma
transformada de la gaussiana, por su paridad. La igualdad de theta
afecta a las series completas y conserva el dominio analítico de
convergencia.

\begin{proposition}[Extensión holomorfa]
\label{prop:leech-theta-complex}
Para \(\tau\) en el semiplano superior
\(\mathbb H=\{\tau\in\mathbb C:\operatorname{Im}\tau>0\}\), la serie
\begin{equation}
 \theta_s(\tau)=\sum_{\lambda\in\Lambda}
             e^{\pi i\tau\|D_s\lambda\|_B^2}
 \label{eq:leech-theta-complex}
\end{equation}
converge normalmente y satisface
\begin{equation}
 \theta_s(-1/\tau)=(-i\tau)^{12}\theta_{-s}(\tau).
 \label{eq:leech-theta-modular}
\end{equation}
\end{proposition}
\begin{proof}
El valor absoluto de cada sumando es
\(e^{-\pi\operatorname{Im}\tau\|D_s\lambda\|_B^2}\).
Sobre compactos del semiplano superior, la parte imaginaria tiene
una cota inferior positiva, de modo que
\eqref{eq:leech-quadratic-bounds} prueba convergencia normal y
holomorfía. Para \(\tau=it\), con \(t>0\),
\eqref{eq:leech-theta-reciprocity} aplicada en \(t^{-1}\) da
\(\theta_s(i/t)=t^{12}\theta_{-s}(it)\).
Ambos miembros de \eqref{eq:leech-theta-modular} son holomorfos
en \(\mathbb H\); el teorema de identidad prolonga la igualdad
desde el eje imaginario positivo.
\end{proof}

La inversión modular relaciona los parámetros \(s\) y \(-s\).
La periodicidad adicional bajo \(\tau\mapsto\tau+1\) utiliza la
paridad de las normas. La proposición~\ref{prop:leech-nonintegral}
exhibe parámetros arbitrariamente próximos a cero para los que
la integralidad euclídea se pierde. Por ello la familia conserva
la reciprocidad analítica demostrada, mientras una realización
mediante formas modulares escalares o álgebras de vértices
reticulares incorpora además las hipótesis aritméticas
correspondientes. En la realización excepcional inicial, esas
hipótesis sostienen la continuación hacia las álgebras de vértices
y Monstrous Moonshine \cite{HMT-IV}.

\subsection{Forma partida y autodualidad en dimensión cuarenta y ocho}
\label{subsec:leech-split}

La pérdida de integralidad de una proyección euclídea conduce a
considerar ambas proyecciones de manera simultánea. En el espacio
\(E\oplus E\) se introduce la forma bilineal
\begin{equation}
 \mathcal B((x,p),(x',p'))=
       \langle x,p'\rangle_B+\langle p,x'\rangle_B.
 \label{eq:leech-split-form}
\end{equation}
Sus productos emparejan las dos coordenadas. Definimos
\begin{equation}
\begin{split}
 \mathcal D_s&=D_s\oplus D_{-s},\\
 \Gamma_s&=\mathcal D_s(\Lambda\oplus\Lambda)\\
 &=\{(D_s\lambda,D_{-s}\mu):\lambda,\mu\in\Lambda\}.
\end{split}
\label{eq:leech-split-lattice}
\end{equation}

\begin{theorem}[Autodualidad integral para la forma partida]
\label{thm:leech-split}
La forma \(\mathcal B\) tiene signatura \((24,24)\).
Para todo \(s\in\mathbb R\), \(\mathcal D_s\) es una isometría de
\(\mathcal B\), y \(\Gamma_s\) es una red integral, par y autodual
respecto de esta forma.
\end{theorem}
\begin{proof}
Introduciendo
\[
 p_L=\frac{x+p}{\sqrt2},\qquad
 p_R=\frac{x-p}{\sqrt2},
\]
se obtiene
\[
 \mathcal B((x,p),(x,p))=\|p_L\|_B^2-\|p_R\|_B^2.
\]
Como \(B\) es positiva en dimensión \(24\), la signatura es
\((24,24)\). La autoadjunción y reciprocidad de los bloques dan
\[
 \langle D_sx,D_{-s}p'\rangle_B=\langle x,p'\rangle_B,
\]
y análogamente para el segundo producto cruzado. Por ello
\(\mathcal B(\mathcal D_s\xi,\mathcal D_s\eta)=\mathcal B(\xi,\eta)\).

En \(\Gamma_0=\Lambda\oplus\Lambda\), todos los productos cruzados
son enteros y
\[
 \mathcal B((\lambda,\mu),(\lambda,\mu))
   =2\langle\lambda,\mu\rangle_B\in2\mathbb Z.
\]
Esto prueba integralidad y paridad. Si \((x,p)\) pertenece al dual
de \(\Gamma_0\) para \(\mathcal B\), el emparejamiento con
\((\lambda,0)\) y \((0,\mu)\) exige respectivamente
\(p,x\in\Lambda^*=\Lambda\). El dual es, por tanto, \(\Gamma_0\).
La isometría \(\mathcal D_s\) transporta estas tres propiedades
a \(\Gamma_s\).
\end{proof}

La métrica especifica el contenido de la conclusión. La familia
\(\Gamma_s\) conserva la integralidad para los productos cruzados
de \(\mathcal B\), incluso cuando \(\Lambda_s\) pierde integralidad
para \(B\). Como retículos abstractos con forma partida, los
\(\Gamma_s\) son isométricos. Sus dos proyecciones euclídeas,
consideradas con las métricas positivas respectivas, varían con
el parámetro.

\begin{proposition}[Intercambio de las proyecciones recíprocas]
\label{prop:leech-split-involution}
La involución \(\mathcal J(x,p)=(p,x)\) satisface
\[
 \mathcal J\Gamma_s=\Gamma_{-s},\qquad
 \mathcal J\mathcal D_s\mathcal J=\mathcal D_{-s}.
\]
En las coordenadas \((p_L,p_R)\), fija \(p_L\) y cambia el signo
de \(p_R\).
\end{proposition}
\begin{proof}
El intercambio convierte
\((D_s\lambda,D_{-s}\mu)\) en
\((D_{-s}\mu,D_s\lambda)\), que recorre \(\Gamma_{-s}\).
La conjugación de los dos bloques da la segunda identidad.
Finalmente, \(x+p\) se conserva y \(x-p\) cambia de signo.
\end{proof}

La involución expresa una dualidad de la construcción reticular
completa. Su eventual representación física conserva, además de
la forma partida, los operadores y los datos de identificación
propios de esa realización.

\subsection{Energía reticular y equivalencia unitaria}
\label{subsec:leech-energy}

La energía de radio recíproco
\(m^2/R^2+w^2R^2\) intercambia sus dos términos bajo
\((m,w,R)\mapsto(w,m,R^{-1})\), con una conjugación unitaria
que preserva los dominios operatorios \cite{HMT-IV}.
La realización sobre los doce planos marcados emplea
\begin{equation}
 \mathcal E_s(\lambda,\mu)=
       \|D_s\lambda\|_B^2+\|D_{-s}\mu\|_B^2,
 \qquad \lambda,\mu\in\Lambda.
 \label{eq:leech-energy}
\end{equation}
La energía es positiva y cumple
\(\mathcal E_s(\lambda,\mu)=\mathcal E_{-s}(\mu,\lambda)\).

\begin{theorem}[Operador positivo y conjugación de dominios]
\label{thm:leech-energy}
En el espacio de Hilbert \(\ell^2(\Lambda\oplus\Lambda)\), sea
\[
 (H_s\psi)(\lambda,\mu)=
       \mathcal E_s(\lambda,\mu)\psi(\lambda,\mu)
\]
con dominio
\begin{equation}
 \mathcal D(H_s)=
 \left\{\psi\in\ell^2(\Lambda\oplus\Lambda):
 \sum_{\lambda,\mu}\mathcal E_s(\lambda,\mu)^2
             |\psi(\lambda,\mu)|^2<\infty\right\}.
 \label{eq:leech-energy-domain}
\end{equation}
Entonces \(H_s\) es positivo, autoadjunto y de resolvente compacto.
La aplicación \(U\psi(\lambda,\mu)=\psi(\mu,\lambda)\) es unitaria y
\begin{equation}
 U\mathcal D(H_s)=\mathcal D(H_{-s}),\qquad
 UH_sU^{-1}=H_{-s}.
 \label{eq:leech-unitary}
\end{equation}
\end{theorem}
\begin{proof}
La multiplicación por una función real en un conjunto discreto,
con el dominio máximo \eqref{eq:leech-energy-domain}, es
autoadjunta. Puede comprobarse directamente mediante la base
ortonormal de los vectores de soporte unitario: la pertenencia
al dominio del adjunto exige que la sucesión obtenida al
multiplicar por \(\mathcal E_s\) sea cuadrado sumable, exactamente
la misma condición. Además,
\[
 \langle\psi,H_s\psi\rangle
   =\sum_{\lambda,\mu}\mathcal E_s(\lambda,\mu)
                         |\psi(\lambda,\mu)|^2\geq0.
\]
Por \eqref{eq:leech-quadratic-bounds},
\[
 \mathcal E_s(\lambda,\mu)\geq
 e^{-2|s|M}\bigl(\|\lambda\|_B^2+\|\mu\|_B^2\bigr).
\]
Una red tiene finitos puntos en cualquier conjunto acotado; por
ello cada subnivel de \(\mathcal E_s\) es finito. Los coeficientes
diagonales \((1+\mathcal E_s)^{-1}\) tienden a cero fuera de
conjuntos finitos, de modo que \((H_s+I)^{-1}\) es límite en norma
de operadores de rango finito y es compacto.

El intercambio de los dos índices conserva la norma de
\(\ell^2(\Lambda\oplus\Lambda)\), por lo que \(U\) es unitario
y \(U^{-1}=U\). La igualdad de energías al invertir \(s\) transforma
la condición de dominio de \(H_s\) en la de \(H_{-s}\), y la
evaluación sobre cada par \((\lambda,\mu)\) prueba
\eqref{eq:leech-unitary}.
\end{proof}

Quedan así articuladas tres conservaciones de distinto dominio.
En los planos euclídeos, el flujo mantiene covolumen y simetría
ternaria y relaciona cada red con el dual del parámetro opuesto.
En el espacio duplicado, los productos cruzados conservan una
red integral, par y autodual. En la realización operatoria, el
intercambio recíproco conserva la energía mediante una equivalencia
unitaria de operadores con sus dominios. Las fórmulas theta
expresan la misma reciprocidad en la suma convergente de las
longitudes de toda la red.

% Procedencia primaria (lectura de 18-09-2026):
% XROOT = output/EXTREMA_MEDIA_RAZON_NARRACION_Y_REVISION_20260916/ARTICULO/ES/source
% XROOT/libro/01_acoplamiento.tex:51-147,727-749.
% XROOT/sections/k_direccion_dimensional.tex:54-155.
% XROOT/nuclear/12.tex:9-88.
% VIROOT = output/AMPLIACION_FORMAL_LEAN_SERIE_20260916/delivery/USB_ES_EN_STAGE10_20260917/ES/PAQUETES/06_PARTICULAS_Y_MASAS_ES/documentacion_original/payload/source_es
% VIROOT/sections/05_registro_operador_masa.tex:11-119,193-294.
% VIROOT/sections/02_particula_persistencia.tex:210-235.
% XROOT/nuclear/07.tex:5-25,74-154: Hamiltoniano de transporte eliptico.
% XROOT/nuclear/11.tex:878-979: elevacion cotangente y accion discreta.
% Estatutos: cadena K->alpha->accion->masa = RESULTADO_RECUPERADO;
% carta 1+1+10 de eventos y criterio de entrelazamiento = FORMALIZACION_NUEVA;
% congruencia, espectro generalizado y control racional = pruebas completas aqui.

\section{Congruencias positivas, lectores de acción y dinámica hamiltoniana}
\label{sec:k-mass-bridge}

La relación entre el registro dodecafásico y la Ley General de Masas tiene
una composición efectiva. APP produce las dos hojas con sus residuos y
cocientes; TRIT orienta la transición y determina su régimen; TPK compone
selección, transporte y actualización conservando ruta, acarreo, frontera,
supervivencia y memoria. Del estado enriquecido y de su prolongación
compatible proceden conjuntamente las publicaciones regionales y el
registro de doce posiciones. La clausura numérica y la lectura material
comparten ese origen y conservan sus operaciones específicas.

Con origen de fase, orientación y base mil declarados, el registro es
\[
 K=(234,543,140,729,659,824,621,58,914,794,146,601).
\]
El lector posicional utiliza
\begin{equation}
 \kappa_{\rm ag}(K)=\frac{\sum_{j=1}^{12}K_j1000^{12-j}}{1000^{12}-1},
 \qquad
 \alpha=\pi_{\rm HMT}+e_{\rm HMT}-\varphi_{\rm HMT}-4-
              \kappa_{\rm ag}(K).
 \label{eq:mass-k-alpha}
\end{equation}
Los valores de la derecha son publicaciones internas del mismo proceso
coinductivo. La fórmula expresa el lector de clausura infinita compatible,
con la frontera de acarreo que lo define.

La acción recibe esas publicaciones mediante los lectores
\begin{align}
 H_5={}&\frac12\sqrt{1000\alpha/\varphi}-\alpha+
       \frac9{16}\alpha^2-\frac59\alpha^3+
       \frac7{48}\alpha^4-\frac1{54}\alpha^5,\notag\\
 D_A={}&\exp(-100\pi\alpha/9),\qquad
 \mathcal C_\pi=169\alpha^6+\frac{D_A\alpha^7}{1-D_A\alpha},\notag\\
 \eta_{\rm ret}={}&H_5-\frac{90}{\pi}\mathcal C_\pi,
 \qquad R_{\rm act}=\frac{H_5}{\eta_{\rm ret}}>0.
 \label{eq:mass-k-action}
\end{align}
Aquí y en las fórmulas siguientes, \(\pi,\varphi,e\) abrevian las
publicaciones HMT ya construidas. La norma determinantal local
\(\Sigma_H=\varphi\alpha^{16}\) fija la década \(-34\), y la
sección \(\hbar_{\rm ret}=\eta_{\rm ret}10^{-34}\mathcal U_S\)
pertenece a la recta de acción declarada. La razón \(R_{\rm act}\)
conserva la relación entre las dos secciones de acción al cambiar su base
dimensional común.

La ruta electrónica central proporciona, mediante sus lectores de
desplazamiento y autocorrelación, el triple \((80,54,6)\). Con
\[
 \Delta_4=\frac{\pi-e\log\pi}{270}\left(1+\frac\pi{729}\right),
 \qquad
 b_e^{\rm alg}=\frac{\sqrt3}{4}
       \bigl[\exp(\varphi/\pi^2)-22\alpha^3\bigr](1+15\Delta_4),
\]
su publicación material contiene
\begin{equation}
 \kappa_e=80-54\alpha+6\Delta_4,
 \qquad \varepsilon_e^\beta=b_e^{\rm alg}R_{\rm act}^{\kappa_e}.
 \label{eq:mass-k-electron}
\end{equation}
La composición \eqref{eq:mass-k-alpha}--\eqref{eq:mass-k-electron}
identifica la intervención concreta de \(K\): su lectura de clausura
participa en \(\alpha\), y esa salida actúa después en el lector de
acción y en el exponente de retorno. La elección central del electrón
pertenece a la involución de la carta APP, cuyo único punto fijo es
\((5,5)\). La centralidad de esa celda y su evaluación escalar son
dos resultados enlazados por la ruta y su registro.

\subsection{Coordenadas reversibles del registro de eventos materiales}

El registro material asigna a una ruta \(\gamma\) doce eventos
orientados \(\mathbf e(\gamma)=(e_0,\ldots,e_{11})\), uno por cada
ventana nonádica del calendario de ciento ocho posiciones. El mismo
orden temporal permite realizar estos eventos en el portador de doce
coordenadas utilizado por \(K\). Esta identificación requiere conservar
el origen y la orientación del calendario.

Sean \(\mathbf1\) el vector uniforme,
\(P_{11}=I-\mathbf1\mathbf1^*/12\), y \(P_3\) el proyector
tridimensional de la carta incidencial declarada. La construcción
geométrica del registro determina
\[
 u_K=P_3P_{11}K,\qquad
 \|u_K\|^2=\frac{6638585+2275584\sqrt5}{20}>0,
 \qquad u_K\perp\mathbf1,
\]
y el proyector transversal
\[
 P_{10}(K)=P_{11}-\frac{u_Ku_K^*}{\|u_K\|^2}.
\]

\begin{proposition}[Descomposición material relativa a la dirección de \(K\)]
Para todo vector de eventos \(\mathbf e\), la aplicación
\begin{equation}
 \mathbf e\longmapsto
 \left(s_C=\mathbf1^*\mathbf e,\quad
       a_K=\frac{u_K^*\mathbf e}{\|u_K\|^2},\quad
       v_K=P_{10}(K)\mathbf e\right)
 \label{eq:mass-k-event-chart}
\end{equation}
es invertible sobre su imagen y satisface
\begin{equation}
 \mathbf e=\frac{s_C}{12}\mathbf1+a_Ku_K+v_K,
 \qquad
 \|\mathbf e\|^2=\frac{|s_C|^2}{12}
          +|a_K|^2\|u_K\|^2+\|v_K\|^2.
 \label{eq:mass-k-event-inverse}
\end{equation}
\end{proposition}
\begin{proof}
Los proyectores \(\mathbf1\mathbf1^*/12\),
\(u_Ku_K^*/\|u_K\|^2\) y \(P_{10}(K)\) son ortogonales dos a
dos y suman la identidad. Aplicar esa igualdad a \(\mathbf e\)
demuestra la inversión. La ortogonalidad de los tres sumandos da la
identidad de normas.
\end{proof}

Esta carta conserva el vector de eventos completo. Por ello conserva
también los funcionales que dependen de sus signos, sus sumas y sus
correlaciones, una vez realizada la reconstrucción
\eqref{eq:mass-k-event-inverse}. El registro de predecesores de los nueves,
la cuenta de acción y las coordenadas de torre mantienen su residencia
en la trayectoria enriquecida: la carta anterior modifica únicamente
las coordenadas de los doce eventos. El carácter se evalúa sobre la misma
firma recuperada. La construcción proporciona así una representación
geométrica de la información utilizada por la masa, con una dirección
distinguida por \(K\) y diez componentes transversales.

La cantidad de información que retiene cada publicación puede probarse
exactamente. En el dominio algebraico de registros, \(K'=K+\mathbf1\)
satisface
\[
 u_{K'}=u_K,\qquad P_{10}(K')=P_{10}(K),\qquad
 \kappa_{\rm ag}(K')-\kappa_{\rm ag}(K)=\frac1{999}.
\]
La última identidad resulta de sumar la progresión geométrica de los
doce pesos posicionales. Si las publicaciones regionales se mantienen
fijas, la lectura de \(\alpha\) cambia en \(-1/999\).
Este control algebraico determina el alcance exacto de la dirección:
su proyector conserva una polarización, mientras la clausura escalar
utiliza además la información uniforme y posicional del registro.
La admisibilidad de \(K'\) como otra historia terminal requiere su
propia construcción APP--TRIT--TPK.

\subsection{Congruencia positiva y espectro del operador material}

Para una multisección \(M\), sea \(\mathscr H_M\) el espacio finito
con base ortonormal etiquetada por sus rutas. El lector de retorno
\(r\in\{D,\Omega\}\) produce un operador autoadjunto
\(B_M^r e_\gamma=\kappa_\gamma^r e_\gamma\). El carácter refinado,
evaluado sobre diferencias respecto del electrón, da el operador basal
\[
 M_M^{(0)}e_\gamma=
 m_e^{(0)}\chi_{\rm ref}(\sigma_\gamma-\sigma_e,
                  \eta_\gamma-\eta_e)e_\gamma.
\]
Los coeficientes de torre y su convergencia pertenecen al lector de cada
ruta. Su publicación beta es la congruencia
\begin{equation}
 D_M^r=R_{\rm act}^{B_M^r/2},\qquad
 M_M^{\beta,r}=D_M^rM_M^{(0)}D_M^r.
 \label{eq:mass-k-congruence}
\end{equation}

\begin{proposition}[Invariantes exactos de la congruencia material]
El operador \(D_M^r\) es positivo e invertible. La congruencia
\eqref{eq:mass-k-congruence} conserva rango e inercia; en particular,
conserva la positividad. Si \(M_M^{(0)}\) y \(B_M^r\) son diagonales
en la base de rutas, entonces
\[
 m_\gamma^\beta=m_\gamma^{(0)}R_{\rm act}^{\kappa_\gamma^r}.
\]
Para dos rutas en la misma carta dimensional se obtiene
\[
 \frac{m_Y^\beta}{m_X^\beta}=
 \chi_{\rm ref}(\sigma_Y-\sigma_X,\eta_Y-\eta_X)
 R_{\rm act}^{\kappa_Y^r-\kappa_X^r}.
\]
\end{proposition}
\begin{proof}
El cálculo funcional da
\(D_M^r=\exp((\log R_{\rm act})B_M^r/2)\), con espectro
estrictamente positivo. Su inversa es la exponencial del opuesto.
Para todo \(v\), la forma cuadrática transformada vale
\((D_M^rv)^*M_M^{(0)}(D_M^rv)\). El cambio invertible de variable
preserva las dimensiones máximas de subespacios positivos, negativos
y nulos; éstas determinan la inercia y el rango. En una base común
diagonal, los dos factores aportan conjuntamente
\(R_{\rm act}^{\kappa_\gamma^r}\). Dividir dos valores cancela
la sección dimensional común y utiliza la ley multiplicativa del
carácter.
\end{proof}

La distinción entre transporte de una forma y modificación de un
operador respecto de una métrica fija es comprobable sin aproximaciones.
Para
\[
 M_0=\begin{pmatrix}1&0\\0&4\end{pmatrix},\qquad
 D=\begin{pmatrix}2&0\\0&1/2\end{pmatrix},\qquad
 M'=D^*M_0D=\begin{pmatrix}4&0\\0&1\end{pmatrix},
\]
el determinante de \(D\) es uno y los valores asociados a las dos
rutas se intercambian. Incluso el espectro no ordenado cambia si
\(M_0=I\): se transforma de \(\{1,1\}\) en \(\{4,1/4\}\).
Ambos controles preservan positividad, rango e inercia.

\begin{proposition}[Invarianza espectral con métrica transportada]
Sean \(G_0>0\) una métrica y \(D\) invertible. Si
\[
 M'=D^*M_0D,\qquad G'=D^*G_0D,
\]
los pares \((M',G')\) y \((M_0,G_0)\) tienen el mismo espectro
generalizado, con las mismas multiplicidades.
\end{proposition}
\begin{proof}
Se tiene
\[
 \det(M'-\lambda G')=|\det D|^2\det(M_0-\lambda G_0),
 \qquad (G')^{-1}M'=D^{-1}G_0^{-1}M_0D.
\]
La primera igualdad preserva las raíces y sus multiplicidades; la
segunda exhibe una semejanza de los operadores asociados.
\end{proof}

El flujo diagonal positivo inducido por \(K\) posee, por tanto, dos
realizaciones matemáticas precisas sobre una forma material. Cuando
transporta simultáneamente la forma y su métrica describe el mismo
objeto en otra carta. Cuando actúa sobre la forma conservando la métrica
de rutas, produce una familia operatoria que puede cambiar las masas
relativas. Su identificación física se decide por el lector de rutas y
por la acción declarada.

\subsection{Compatibilidad del flujo dodecafásico con el retorno material}

Sea
\[
 A_K=(\log\rho_A)\operatorname{diag}(\widehat u_{K,1},\ldots,
       \widehat u_{K,12}),\quad
 \widehat u_K=u_K/\|u_K\|,\quad \rho_A=\pi/\varphi^2,
\]
de modo que \(g_K(t)=e^{tA_K}\). La familia de retorno material
en la multisección \(M\) tiene generador
\[
 C_M^r=\tfrac12(\log R_{\rm act})B_M^r,\qquad
 S_M^r(t)=e^{tC_M^r}.
\]

\begin{proposition}[Transporte de la congruencia por un lector entrelazante]
Una aplicación lineal declarada
\(J_M:\mathbb C^{12}\to\mathscr H_M\) que cumple
\begin{equation}
 C_M^rJ_M=J_MA_K
 \label{eq:mass-k-intertwiner}
\end{equation}
satisface
\[
 S_M^r(t)J_M=J_Mg_K(t),
\]
y, para \(M_M(t)=S_M^r(t)M_M^{(0)}S_M^r(t)\),
\begin{equation}
 J_M^*M_M(t)J_M
 =g_K(t)^*\bigl(J_M^*M_M^{(0)}J_M\bigr)g_K(t).
 \label{eq:mass-k-pullback}
\end{equation}
\end{proposition}
\begin{proof}
La identidad \eqref{eq:mass-k-intertwiner} se itera a cada potencia
y se suma en la serie exponencial convergente. Tomar adjuntos y
sustituir en la forma material demuestra
\eqref{eq:mass-k-pullback}.
\end{proof}

El criterio es finito y admite una prueba de esfuerzo directamente
operatoria. En las bases diagonales declaradas, si \(a_j\) y
\(c_\gamma\) son los autovalores de \(A_K\) y \(C_M^r\),
respectivamente, se debe cumplir
\[
 (c_\gamma-a_j)(J_M)_{\gamma j}=0
 \qquad\text{para cada }(\gamma,j).
\]
Cada coeficiente no nulo del lector exige la igualdad correspondiente
de autovalores. La composición recuperada de
\eqref{eq:mass-k-alpha}--\eqref{eq:mass-k-electron} y la carta reversible
\eqref{eq:mass-k-event-chart} son resultados efectivos. La identificación
adicional de sus dos flujos requiere un lector \(J_M\) con la propiedad
\eqref{eq:mass-k-intertwiner}; esta proposición determina exactamente
qué debe comprobar una realización que la sostenga.

\subsection{Transporte elíptico y transformación regular de Legendre}

La realización variacional del transporte tiene un antecedente explícito
en la elipse de acción. El par angular ya publicado determina
\(\eta=\operatorname{artanh}(C^*/A)\), en la cámara
\(A\ne0\), \(|C^*/A|<1\). Las coordenadas \((Q,P)\), ambas de
dimensión raíz de acción, realizan la forma
\[
 \mathcal I_\eta(Q,P)=\frac12\bigl(e^{-\eta}Q^2+e^\eta P^2\bigr),
 \qquad \lambda=P\,dQ,\quad \Omega=dQ\wedge dP.
\]
El transporte entre formas, con
\(M_\eta=\operatorname{diag}(e^{\eta/2},e^{-\eta/2})\) y
\(J_2=\left(\begin{smallmatrix}0&-1\\1&0\end{smallmatrix}\right)\),
es \(M_{\eta'}e^{-\vartheta J_2}M_\eta^{-1}\). Una
interpolación \(\eta(t)\), \(\vartheta(t)\), con
\(\dot\vartheta=\omega(t)>0\), tiene el Hamiltoniano recuperado
del artículo X:
\begin{equation}
 H_{\rm tr}(Q,P,t)=\omega(t)\mathcal I_{\eta(t)}(Q,P)
                 +\frac{\dot\eta(t)}2QP.
 \label{eq:mass-k-hamiltonian}
\end{equation}
El segundo término procede de \(\dot M_\eta M_\eta^{-1}\).
Las ecuaciones
\[
 \dot Q=\omega e^\eta P+\frac{\dot\eta}2Q,
 \qquad
 \dot P=-\omega e^{-\eta}Q-\frac{\dot\eta}2P
\]
conservan \(\mathcal I_{\eta(t)}(Q(t),P(t))\): la variación
explícita de \(\eta\) y la contribución del término de conexión
se cancelan. Este resultado y la identidad
\(P\dot Q-H_{\rm tr}=p\dot q-\omega(q^2+p^2)/2\), para
\(Q=e^{\eta/2}q\), \(P=e^{-\eta/2}p\), están demostrados
en la sección de conservación variacional de la obra de extrema y
media razón holográfica.

\begin{proposition}[Lagrangiano regular del transporte elíptico]
Supongamos \(\eta\in C^2\), \(\omega\in C^1\) y \(\omega>0\).
La transformación de Legendre respecto de \(P\) es invertible y
produce
\begin{equation}
 L_{\rm tr}(Q,\dot Q,t)=
 \frac{\bigl(\dot Q-\dot\eta Q/2\bigr)^2}{2\omega e^\eta}
 -\frac{\omega e^{-\eta}}2Q^2.
 \label{eq:mass-k-lagrangian}
\end{equation}
\end{proposition}
\begin{proof}
Escribamos \(a=\dot\eta/2\), \(b=\omega e^\eta>0\) y
\(c=\omega e^{-\eta}\). Entonces
\(H_{\rm tr}=bP^2/2+aQP+cQ^2/2\), y
\[
 \dot Q=\partial_PH_{\rm tr}=bP+aQ,
 \qquad P=\frac{\dot Q-aQ}{b}.
\]
Al sustituir esta expresión en \(P\dot Q-H_{\rm tr}\), los
términos que contienen \(P\) se reúnen en
\(P(\dot Q-aQ)-bP^2/2=(\dot Q-aQ)^2/(2b)\), lo cual
demuestra \eqref{eq:mass-k-lagrangian}. Además,
\[
 \partial_{\dot Q}L_{\rm tr}=P,
 \qquad
 \partial_P^2H_{\rm tr}=\omega e^\eta>0,
 \qquad
 \partial_{\dot Q}^2L_{\rm tr}=\frac1{\omega e^\eta}>0.
\]
La transformación es regular y sus ecuaciones de Euler--Lagrange
equivalen a las ecuaciones hamiltonianas anteriores.
\end{proof}

Esta eliminación algebraica reúne el Hamiltoniano ya construido con
su forma lagrangiana de segundo orden. La regularidad procede del
término cuadrático de la realización elíptica. Para un levantamiento
cotangente puramente lineal \(H_K(q,p)=p^{\mathsf T}A_Kq\), la
ecuación \(\dot q=A_Kq\) es independiente de \(p\) y el Hessiano
\(\partial_p^2H_K\) es nulo. Su formulación variacional utiliza la
acción de primer orden \(\int p^{\mathsf T}(\dot q-A_Kq)\,dt\).
El corpus contiene también su antecedente discreto
\(L_n=p_{n+1}^{\mathsf T}(q_{n+1}-U_nq_n)\), cuyas ecuaciones
producen el transporte \(q_{n+1}=U_nq_n\) y el transporte dual
\(p_{n+1}=U_n^{-\mathsf T}p_n\). Las dos formulaciones poseen
dominios y operadores determinados. La interpolación temporal y la
sección de acción se conservan como parte de la realización física
declarada.

\begin{figure}[htbp]
\centering
\begin{tikzpicture}[
 box/.style={draw=black!70,rounded corners=1pt,align=center,
             font=\small,inner sep=5pt,text width=40mm},
 arrow/.style={-{Stealth[length=2mm]},thick}]
\node[box] (memory) at (0,0) {Carga y memorias\\\((q,m_0,m_1)\)};
\node[box] (register) at (0,-1.65) {Registro reconstruido\\\(K=(q/12)\mathbf1+Lm\)};
\node[box] (weights) at (-3.8,-3.55) {Partición ordenada\\\(d_j=K_j/q\)};
\node[box] (direction) at (3.8,-3.55) {Dirección de suma nula\\\(u=P_3P_{11}K\)};
\node[box] (geometry) at (-4.6,-6) {Medición de frontera\\y forma canónica\\\(C(d),\ \Omega_{P_K}\)};
\node[box] (energy) at (0,-6) {Energía nodal\\y detalles\\\(L_d,\ \mathcal E(z)\)};
\node[box] (lattice) at (4.6,-6) {Deformación reticular\\y red dual\\\(\Lambda_s^*=\Lambda_{-s}\)};
\draw[arrow] (memory)--(register);
\draw[arrow] (register)--(weights);
\draw[arrow] (register)--(direction);
\draw[arrow] (weights)--(geometry);
\draw[arrow] (weights)--(energy);
\draw[arrow] (direction)--(lattice);
\draw[arrow] (direction) to[bend left=12] node[right,font=\scriptsize,align=left,fill=white,inner sep=1pt]{flujo\\heredado} (energy);
\end{tikzpicture}
\caption{Realizaciones posteriores del registro recuperable. La rama
intervalar determina los menores de caminos, la forma canónica del
polígono y el laplaciano nodal. La dirección algebraica determina el
flujo de intervalos y, junto con la realización reticular marcada, la
familia dual de veinticuatro dimensiones. Cada flecha conserva el dominio
y las hipótesis de su construcción; los codominios inferiores son
distintos.}
\label{fig:multimorphology}
\end{figure}


\section{Conclusiones: reconstrucción, invariantes y dualidad}
El resultado estructural prioritario es la recuperabilidad de la forma.
La carga uniforme y las dos memorias nonádicas reconstruyen el registro
mediante una inversa explícita; el cono positivo describe exactamente
el dominio donde esa reconstrucción induce menores estrictamente positivos.
La cota de estabilidad cuantifica la persistencia de ese dominio. Por
tanto, los pesos geométricos admiten una procedencia y una recuperación
operatorias, además de su evaluación numérica.

La composición desarrollada tiene una secuencia explícita. El registro
terminal genera doce pesos positivos; sus sumas parciales dan trece puntos
ordenados; una red planar realiza sus menores; la curva de momentos
construye polígonos y politopos; sus triangulaciones proporcionan cobertura
y formas con residuos orientados. El refinamiento de una celda fija conserva
la suma de sus formas. La introducción de nuevos vértices en la curva de
momentos produce, por su parte, otra geometría y exige comparar dominios.

La dirección algebraica del registro distingue dos operaciones. El
reescalado positivo de columnas conserva las razones proyectivas pertinentes.
Su acción normalizada sobre los intervalos genera en cambio una deformación
efectiva, certificada por una derivada exacta no nula. Esta deformación admite
un potencial convexo cuya transformada de Legendre se conserva al subdividir
con dirección heredada. La invariancia tiene así un mecanismo comprobable:
cada hijo conserva el generador de su padre y la suma de sus pesos recupera
el peso anterior.

El refinamiento energético proporciona una segunda conservación exacta.
La reducción de Schur recupera el laplaciano efectivo y la prolongación
minimizante separa cada configuración en su parte observada y un detalle
reconstructible. Las energías de los detalles se suman ortogonalmente.
Con mallas decrecientes y energía uniformemente acotada, esa sucesión
converge en la realización de Sobolev declarada. La invariancia de la
forma canónica, la identidad del potencial convexo y la conservación de
energía son tres resultados distintos de compatibilidad, demostrados
para los refinamientos especificados.

La prolongación a los doce planos reticulares conserva una simetría
ternaria y el volumen. El resultado decisivo es la reciprocidad
\(\Lambda_s^*=\Lambda_{-s}\), que induce la inversión
\(\Theta_s(t)=t^{-12}\Theta_{-s}(t^{-1})\). La realización partida de
signatura \((24,24)\) conserva integralidad, paridad y autodualidad para
todo el flujo. En la proyección euclídea, la variación explícita de una
norma distingue deformación y automorfismo de la red inicial. Esta
distinción determina el contenido exacto de la conexión con la dualidad
de momentos y devanados.

La realización material conserva sus propios lectores. La publicación de
acción y el carácter de masas reciben el registro por composiciones ya
construidas; la congruencia preserva positividad e inercia, mientras la
métrica determina la interpretación de sus autovalores. Estas diferencias
precisan el significado de multimorfología: varias formas matemáticas
proceden de una genealogía común y se relacionan por operaciones verificables.
La identificación completa de amplituedros de rango superior requerirá
las pruebas de cobertura y residuos para esos mapas concretos. El resultado
de rango uno demostrado aquí fija un caso completo, con coordenadas
racionales, triangulación y controles reproducibles.

La orientación temporal del TRIT integra el retorno con memoria, la
actualización de umbral y las prolongaciones conjugadas. Sus generadores
elíptico, nilpotente e hiperbólico proporcionan leyes distintas de
evolución; la composición del TPK conserva la genealogía que las
concatena. En particular, la actividad del régimen cero queda expresada
por un transporte unipotente, y la doble proyección conserva simultáneamente
la convergencia de las coordenadas y la incidencia excepcional. La tesis
geométrica del artículo reside en esta articulación: la estructura de
partida genera registros recuperables y éstos sostienen realizaciones
positivas cuyas transformaciones, fronteras y memorias se pueden calcular
y comparar con precisión.

\begin{thebibliography}{99}
\bibitem{hmtI} O. Haidara Fall y R. Ramos Balsa,
\emph{Generación coinductiva de \(\pi,\varphi,e\) a profundidad arbitraria,
derivación de la constante de estructura fina y estructura algebraica del
electrón}. Serie HMT--MD, artículo I, edición consultada de septiembre de 2026.
\bibitem{hmtX} O. Haidara Fall y R. Ramos Balsa,
\emph{Extrema y media razón holográfica}. Desarrollo del acoplamiento
aritmético--geométrico, registro dodecafásico y dirección dimensional,
edición consultada de septiembre de 2026.
\bibitem{hmtVI} O. Haidara Fall y R. Ramos Balsa,
\emph{Partículas, persistencia y espectro de masas}. Serie HMT--MD, artículo VI,
edición consultada de septiembre de 2026.
\bibitem{HMT-IV} O. Haidara Fall y R. Ramos Balsa,
\emph{Moonshine, dualidad T y teoría M desde la estructura discreta del
continuo}. Serie HMT--MD, artículo IV, edición consultada de septiembre de 2026.
\bibitem{HMT-VIII} O. Haidara Fall y R. Ramos Balsa,
\emph{Aritmética genealógica de los números primos y estructura espectral
de la función zeta}. Serie HMT--MD, artículo VIII, edición consultada de septiembre de 2026.
\bibitem{HMT-IX} O. Haidara Fall y R. Ramos Balsa,
\emph{Resolución de la hipótesis del continuo en la clausura genealógica
holográfica}. Serie HMT--MD, artículo IX, edición consultada de septiembre de 2026.
\bibitem{Elkies} N. D. Elkies,
\emph{Rational lattices and their theta functions}. Arizona Winter School,
2009, teorema 2.
\href{https://people.math.harvard.edu/~elkies/aws09.pdf}{Notas del autor}.
\bibitem{ConwaySloane} J. H. Conway y N. J. A. Sloane,
\emph{Sphere Packings, Lattices and Groups}, tercera edición, Springer, 1999.
\href{https://doi.org/10.1007/978-1-4757-6568-7}{DOI: 10.1007/978-1-4757-6568-7}.
\bibitem{postnikov} A. Postnikov,
\emph{Total positivity, Grassmannians, and networks}, 2006,
\href{https://arxiv.org/abs/math/0609764}{arXiv:math/0609764}.
\bibitem{positive} N. Arkani-Hamed, Y. Bai y T. Lam,
\emph{Positive Geometries and Canonical Forms}, 2017,
\href{https://arxiv.org/abs/1703.04541}{arXiv:1703.04541}.
\end{thebibliography}
\end{document}
